%%%%%%%%%%%%%%%%%%%%%%%%%%%%%%%%%%%%%%%%%%%%%%%%%%%%%%%%%%%%%%%%%%%%%%%%%
%
% Plantilla para libro de texto de matemáticas.
%
% Esta plantilla ha sido desarrollada desde cero, pero utiliza algunas partes
% del código de la plantilla original utilizada en apuntesDGIIM
% (https://github.com/libreim/apuntesDGIIM), basada a su vez en las plantillas
% 'Short Sectioned Assignment' de Frits Wenneker (http://www.howtotex.com),
% 'Plantilla de Trabajo' de Mario Román y 'Plantilla básica de Latex en Español'
% de Andrés Herrera Poyatos (https://github.com/andreshp). También recoge
% ideas de la plantilla 'Multi-Purpose Large Font Title Page' de
% Frits Wenneker y Vel (vel@latextemplates.com).
%
% Licencia:
% CC BY-NC-SA 4.0 (https://creativecommons.org/licenses/by-nc-sa/4.0/)
%
%%%%%%%%%%%%%%%%%%%%%%%%%%%%%%%%%%%%%%%%%%%%%%%%%%%%%%%%%%%%%%%%%%%%%%%%%


% ---------------------------------------------------------------------------
% CONFIGURACIÓN BÁSICA DEL DOCUMENTO
% ---------------------------------------------------------------------------

%\documentclass[11pt, a4paper, twoside]{article} % Usar para imprimir
\documentclass[10pt, a4paper]{article}

\linespread{1.3}            % Espaciado entre líneas.
\setlength\parindent{0pt}   % No indentar el texto por defecto.
\setlength\parskip{7pt}

% ---------------------------------------------------------------------------
% PAQUETES BÁSICOS
% ---------------------------------------------------------------------------

% IDIOMA
\usepackage[utf8]{inputenc}
\usepackage[spanish, es-tabla, es-lcroman, es-noquoting]{babel}

% MATEMÁTICAS
\usepackage{amsmath}    % Paquete básico de matemáticas
\usepackage{amsthm}     % Teoremas
\usepackage{mathrsfs}   % Fuente para ciertas letras utilizadas en matemáticas

% FUENTES
\usepackage{newpxtext, newpxmath}   % Fuente similar a Palatino
\usepackage{FiraSans}                 % Fuente sans serif
\usepackage[T1]{fontenc}
\usepackage[italic]{mathastext}     % Utiliza la fuente del documento
                                    % en los entornos matemáticos

% MÁRGENES
\usepackage[margin=2.5cm, top=3cm]{geometry}

% LISTAS
\usepackage{enumitem}       % Mejores listas
\setlist{leftmargin=.5in}   % Especifica la indentación para las listas.

%  OTROS
\usepackage{hyperref}   % Enlaces
\usepackage{graphicx}   % Permite incluir gráficos en el documento

% ---------------------------------------------------------------------------
% COLORES
% ---------------------------------------------------------------------------

\usepackage{xcolor}     % Permite definir y utilizar colores

\definecolor{50}{HTML}{ECEFF1}
\definecolor{100}{HTML}{CFD8DC}
\definecolor{200}{HTML}{B0BEC5}
\definecolor{300}{HTML}{90A4AE}
\definecolor{400}{HTML}{78909C}
\definecolor{500}{HTML}{607D8B}
\definecolor{600}{HTML}{546E7A}
\definecolor{700}{HTML}{455A64}
\definecolor{800}{HTML}{37474F}
\definecolor{900}{HTML}{263238}

% ---------------------------------------------------------------------------
% FUENTE DE MATEMÁTICAS ADAPTABLE EN SANS-SERIF
% ---------------------------------------------------------------------------

\usepackage{sansmath}

\newif\IfInSansMode
\let\oldsf\sffamily
\renewcommand*{\sffamily}{\oldsf\sansmath\InSansModetrue}

% ---------------------------------------------------------------------------
% DISEÑO DE PÁGINA
% ---------------------------------------------------------------------------

\usepackage{pagecolor}
\usepackage{afterpage}

% ---------------------------------------------------------------------------
% CABECERA Y PIE DE PÁGINA
% ---------------------------------------------------------------------------

\usepackage{fancyhdr}   % Paquete para cabeceras y pies de página

% Indica que las páginas usarán la configuración de fancyhdr
\pagestyle{fancy}
\fancyhf{}

% Representa la sección de la cabecera
\renewcommand{\sectionmark}[1]{%
\markboth{#1}{}}

% Representa la subsección de la cabecera
\renewcommand{\subsectionmark}[1]{%
\markright{#1}{}}

% Parte derecha de la cabecera
\fancyhead[LE,RO]{\sffamily\textsl{\rightmark} \hspace{1em}  \textcolor{500}{\rule[-0.4ex]{0.2ex}{1.2em}} \hspace{1em} \thepage}

% Parte izquierda de la cabecera
\fancyhead[RE,LO]{\sffamily{\leftmark}}

% Elimina la línea de la cabecera
\renewcommand{\headrulewidth}{0pt}

% Controla la altura de la cabecera para que no haya errores
\setlength{\headheight}{14pt}

% ---------------------------------------------------------------------------
% TÍTULOS DE PARTES Y SECCIONES
% ---------------------------------------------------------------------------

\usepackage{titlesec}

% Estilo de los títulos de las partes
\titleformat{\part}[hang]{\Huge\bfseries\sffamily}{\thepart\hspace{20pt}\textcolor{500}{|}\hspace{20pt}}{0pt}{\Huge\bfseries}
\titlespacing*{\part}{0cm}{-2em}{2em}[0pt]

% Reiniciamos el contador de secciones entre partes (opcional)
\makeatletter
\@addtoreset{section}{part}
\makeatother

% Estilo de los títulos de las secciones, subsecciones y subsubsecciones
\titleformat{\section}
  {\Large\bfseries\sffamily}{\thesection}{1em}{}

\titleformat{\subsection}
  {\Large\sffamily}{\thesubsection}{1em}{}[\vspace{.5em}]

\titleformat{\subsubsection}
  {\sffamily}{\thesubsubsection}{1em}{}

% ---------------------------------------------------------------------------
% ENTORNOS PERSONALIZADOS
% ---------------------------------------------------------------------------

\usepackage{mdframed}

%% DEFINICIONES DE LOS ESTILOS

% Nuevo estilo para definiciones
\newtheoremstyle{definition-style}  % Nombre del estilo
{}                                  % Espacio por encima
{}                                  % Espacio por debajo
{}                                  % Fuente del cuerpo
{}                                  % Identación
{\bf\sffamily}                      % Fuente para la cabecera
{.}                                 % Puntuación tras la cabecera
{.5em}                              % Espacio tras la cabecera
{\thmname{#1}\thmnumber{ #2}\thmnote{ (#3)}}  % Especificación de la cabecera

% Nuevo estilo para notas
\newtheoremstyle{remark-style}
{10pt}
{10pt}
{}
{}
{\itshape \sffamily}
{.}
{.5em}
{}

% Nuevo estilo para teoremas y proposiciones
\newtheoremstyle{theorem-style}
{}
{}
{}
{}
{\bfseries \sffamily}
{.}
{.5em}
{\thmname{#1}\thmnumber{ #2}\thmnote{ (#3)}}

% Nuevo estilo para ejemplos
\newtheoremstyle{example-style}
{10pt}
{10pt}
{}
{}
{\bf \sffamily}
{}
{.5em}
{\thmname{#1}\thmnumber{ #2.}\thmnote{ #3.}}

% Nuevo estilo para la demostración

\makeatletter
\renewenvironment{proof}[1][\proofname] {\par\pushQED{\qed}\normalfont\topsep6\p@\@plus6\p@\relax\trivlist\item[\hskip\labelsep\itshape\sffamily#1\@addpunct{.}]\ignorespaces}{\popQED\endtrivlist\@endpefalse}
\makeatother

%% ASIGNACIÓN DE LOS ESTILOS

% Teoremas, proposiciones y corolarios
\theoremstyle{theorem-style}
\newtheorem{nth}{Teorema}[section]
\newtheorem{nprop}{Proposición}[section]
\newtheorem{ncor}{Corolario}[section]
\newtheorem{lema}{Lema}[section]

% Definiciones
\theoremstyle{definition-style}
\newtheorem{ndef}{Definición}[section]

% Notas
\theoremstyle{remark-style}
\newtheorem*{nota}{Nota}

% Ejemplos
\theoremstyle{example-style}
\newtheorem{ejemplo}{Ejemplo}[section]

% Ejercicios y solución
\theoremstyle{definition-style}
\newtheorem{ejer}{Ejercicio}[section]

\theoremstyle{remark-style}
\newtheorem*{sol}{Solución}

%% MARCOS DE LOS ESTILOS

% Configuración general de mdframe, los estilos de los teoremas, etc
\mdfsetup{
  skipabove=1em,
  skipbelow=1em,
  innertopmargin=1em,
  innerbottommargin=1em,
  splittopskip=2\topsep,
}

% Definimos los marcos de los estilos

\mdfdefinestyle{nth-frame}{
	linewidth=2pt, %
	linecolor= 500, %
	topline=false, %
	bottomline=false, %
	rightline=false,%
	leftmargin=0em, %
	innerleftmargin=1em, %
  innerrightmargin=1em,
	rightmargin=0em, %
}%

\mdfdefinestyle{nprop-frame}{
	linewidth=2pt, %
	linecolor= 300, %
	topline=false, %
	bottomline=false, %
	rightline=false,%
	leftmargin=0pt, %
	innerleftmargin=1em, %
	innerrightmargin=1em,
	rightmargin=0pt, %
}%

\mdfdefinestyle{ndef-frame}{
	linewidth=2pt, %
	linecolor= 500, %
	backgroundcolor= 50,
	topline=false, %
	bottomline=false, %
	rightline=false,%
	leftmargin=0pt, %
	innerleftmargin=1em, %
	innerrightmargin=1em,
	rightmargin=0pt, %
}%

\mdfdefinestyle{ejer-frame}{
	linewidth=2pt, %
	linecolor= 300, %
	backgroundcolor= 50,
	topline=false, %
	bottomline=false, %
	rightline=false,%
	leftmargin=0pt, %
	innerleftmargin=1em, %
	innerrightmargin=1em,
	rightmargin=0pt, %
}%

\mdfdefinestyle{ejemplo-frame}{
	linewidth=0pt, %
	linecolor= 300, %
	leftline=false, %
	rightline=false, %
	leftmargin=0pt, %
	innerleftmargin=1.3em, %
	innerrightmargin=1em,
	rightmargin=0pt, %
	innertopmargin=0em,%
	innerbottommargin=0em, %
	splittopskip=\topskip, %
}%

% Asignamos los marcos a los estilos
\surroundwithmdframed[style=nth-frame]{nth}
\surroundwithmdframed[style=nprop-frame]{nprop}
\surroundwithmdframed[style=nprop-frame]{ncor}
\surroundwithmdframed[style=ndef-frame]{ndef}
\surroundwithmdframed[style=ejer-frame]{ejer}
\surroundwithmdframed[style=ejemplo-frame]{ejemplo}
\surroundwithmdframed[style=ejemplo-frame]{sol}

% ---------------------------------------------------------------------------
% CONFIGURACIÓN DE LA PORTADA
% ---------------------------------------------------------------------------

\newcommand{\asignatura}{%
  Asignatura de\\
  nombre largo\\
}

\newcommand{\autor}{Autor}

\newcommand{\grado}{Grado en el que se imparte}

\newcommand{\universidad}{Universidad}

\newcommand{\enlaceweb}{ejemplo.github.io}

% ---------------------------------------------------------------------------
% CONFIGURACIÓN PERSONALIZADA
% ---------------------------------------------------------------------------

\usepackage{config}

%%%%%%%%%%%%%%%%%%%%%%%%%%%%%%%%%%%%%%%%%%%%%%%%%%%%%%%%%%%%%%%%%%%%%%%%%%%%%
% ---------------------------------------------------------------------------
% COMIENZO DEL DOCUMENTO
% ---------------------------------------------------------------------------
%%%%%%%%%%%%%%%%%%%%%%%%%%%%%%%%%%%%%%%%%%%%%%%%%%%%%%%%%%%%%%%%%%%%%%%%%%%%%

\begin{document}

% ---------------------------------------------------------------------------
% PORTADA EXTERIOR
% ---------------------------------------------------------------------------

\newpagecolor{500}\afterpage{\restorepagecolor} % Color de la página
\begin{titlepage}

  % Título del documento
	\parbox[t]{\textwidth}{
			\raggedright % Texto alineado a la izquierda
			\fontsize{50pt}{50pt}\selectfont\sffamily\color{white}{
			  \textbf{\asignatura}
      }
	}

	\vfill

	%% Autor e información del documento
	\parbox[t]{\textwidth}{
		\raggedright % Texto alineado a la izquierda
		\sffamily\large\color{white}
		{\Large \autor }\\[4pt]
		\grado\\
		\universidad\\[4pt]
		\texttt{\enlaceweb}
	}

\end{titlepage}

% ---------------------------------------------------------------------------
% PÁGINA DE LICENCIA
% ---------------------------------------------------------------------------

\thispagestyle{empty}
\null
\vfill

%% Información sobre la licencia
\parbox[t]{\textwidth}{
  \includegraphics{by-nc-sa.pdf}\\[4pt]
  \raggedright % Texto alineado a la izquierda
  \sffamily\large
  {\Large Este libro se distribuye bajo una licencia CC BY-NC-SA 4.0.}\\[4pt]
  Eres libre de distribuir y adaptar el material siempre que reconozcas a los\\
  autores originales del documento, no lo utilices para fines comerciales\\
  y lo distribuyas bajo la misma licencia.\\[4pt]
  \texttt{creativecommons.org/licenses/by-nc-sa/4.0/}
}

% ---------------------------------------------------------------------------
% PORTADA INTERIOR
% ---------------------------------------------------------------------------

\begin{titlepage}

  % Título del documento
	\parbox[t]{\textwidth}{
			\raggedright % Texto alineado a la izquierda
			\fontsize{50pt}{50pt}\selectfont\sffamily\color{500}{
			  \textbf{\asignatura}
      }
	}

	\vfill

	%% Autor e información del documento
	\parbox[t]{\textwidth}{
		\raggedright % Texto alineado a la izquierda
		\sffamily\large
		{\Large \autor}\\[4pt]
		\grado\\
		\universidad\\[4pt]
		\texttt{\enlaceweb}
	}

\end{titlepage}

% ---------------------------------------------------------------------------
% ÍNDICE
% ---------------------------------------------------------------------------

\thispagestyle{empty}
\tableofcontents
\newpage

% ---------------------------------------------------------------------------
% CONTENIDO
% ---------------------------------------------------------------------------

\part{Teoría}
%%%%%%%%%%%%%%%%%%%%%%%%%%%%%%%%%%%%%%%%%%%%%%%%%%%%%%%%%%%%%%%%
%
% Apuntes de la asignatura Álgebra I.
% Doble Grado de Informática y Matemáticas.
% Universidad de Granada.
% Curso 2016/17.
%
%
%
% Agradecimientos:
% Andrés Herrera (@andreshp) y Mario Román (@M42) por
% las plantillas base.
%
% Sitio original:
% https://github.com/libreim/apuntesDGIIM/
%
% Licencia:
% CC BY-NC-SA 4.0 (https://creativecommons.org/licenses/by-nc-sa/4.0/)
%
%%%%%%%%%%%%%%%%%%%%%%%%%%%%%%%%%%%%%%%%%%%%%%%%%%%%%%%%%%%%%%%


%------------------------------------------------------------------------------
%   ACKNOWLEDGMENTS
%------------------------------------------------------------------------------

%%%%%%%%%%%%%%%%%%%%%%%%%%%%%%%%%%%%%%%%%%%%%%%%%%%%%%%%%%%%%%%%%%%%%%%%
% Plantilla básica de Latex en Español.
%
% Autor: Andrés Herrera Poyatos (https://github.com/andreshp)
%
% Es una plantilla básica para redactar documentos. Utiliza el paquete  fancyhdr para darle un
% estilo moderno pero serio.
%
% La plantilla se encuentra adaptada al español.
%
%%%%%%%%%%%%%%%%%%%%%%%%%%%%%%%%%%%%%%%%%%%%%%%%%%%%%%%%%%%%%%%%%%%%%%%%%

%%%
% Plantilla de Trabajo
% Modificación de una plantilla de Latex de Frits Wenneker para adaptarla
% al castellano y a las necesidades de escribir informática y matemáticas.
%
% Editada por: Mario Román
%
% License:
% CC BY-NC-SA 3.0 (http://creativecommons.org/licenses/by-nc-sa/3.0/)
%%%

%%%%%%%%%%%%%%%%%%%%%%%%%%%%%%%%%%%%%%%%
% Short Sectioned Assignment
% LaTeX Template
% Version 1.0 (5/5/12)
%
% This template has been downloaded from:
% http://www.LaTeXTemplates.com
%
% Original author:
% Frits Wenneker (http://www.howtotex.com)
%
% License:
% CC BY-NC-SA 3.0 (http://creativecommons.org/licenses/by-nc-sa/3.0/)
%
%%%%%%%%%%%%%%%%%%%%%%%%%%%%%%%%%%%%%%%%%


% Tipo de documento y opciones.
\documentclass[11pt, a4paper, titlepage]{article}


%---------------------------------------------------------------------------
%   PAQUETES
%---------------------------------------------------------------------------

% Idioma y codificación para Español.
\usepackage[utf8]{inputenc}
\usepackage[spanish, es-tabla, es-lcroman, es-noquoting]{babel}
\selectlanguage{spanish}
% \usepackage[T1]{fontenc}

% Fuente utilizada.
\usepackage{courier}    % Fuente Courier.
\usepackage{microtype}  % Mejora la letra final de cara al lector.

% Diseño de página.
\usepackage{fancyhdr}   % Utilizado para hacer títulos propios.
\usepackage{titlesec} 	% Utilizado para hacer títulos propios.
\usepackage{lastpage}   % Referencia a la última página.
\usepackage{extramarks} % Marcas extras. Utilizado en pie de página y cabecera.
\usepackage[parfill]{parskip}    % Crea una nueva línea entre párrafos.
\usepackage{geometry}            % Geometría de las páginas.

% Símbolos y matemáticas.
\usepackage{amssymb, amsmath, amsthm, amsfonts, amscd}
\usepackage{upgreek}

\usepackage{mdframed}

% Otros.
\usepackage{enumitem}   % Listas mejoradas.
\usepackage[hidelinks]{hyperref}
\usepackage{pgfplots}
\usepackage{graphicx}   % Gráficos.

% Dibujos Tikz
\usepackage{pgf,tikz}

% Fuentes personalizadas
\usepackage[scaled=.85]{newpxtext,newpxmath}
\usepackage[scaled=.85]{FiraSans}
\usepackage[T1]{fontenc}

% Código para ajustar las fuentes matemáticas al estilo del texto
% que le rodea

\DeclareMathVersion{sans}
\SetSymbolFont{operators}{sans}{OT1}{cmbr}{m}{n}
\SetSymbolFont{letters}{sans}{OML}{cmbr}{m}{it}
\SetSymbolFont{symbols}{sans}{OMS}{cmbrs}{m}{n}
\SetMathAlphabet{\mathit}{sans}{OT1}{cmbr}{m}{sl}
\SetMathAlphabet{\mathbf}{sans}{OT1}{cmbr}{bx}{n}
\SetMathAlphabet{\mathtt}{sans}{OT1}{cmtl}{m}{n}
\SetSymbolFont{largesymbols}{sans}{OMX}{iwona}{m}{n}

\DeclareMathVersion{boldsans}
\SetSymbolFont{operators}{boldsans}{OT1}{cmbr}{b}{n}
\SetSymbolFont{letters}{boldsans}{OML}{cmbrm}{b}{it}
\SetSymbolFont{symbols}{boldsans}{OMS}{cmbrs}{b}{n}
\SetMathAlphabet{\mathit}{boldsans}{OT1}{cmbr}{b}{sl}
\SetMathAlphabet{\mathbf}{boldsans}{OT1}{cmbr}{bx}{n}
\SetMathAlphabet{\mathtt}{boldsans}{OT1}{cmtl}{b}{n}
\SetSymbolFont{largesymbols}{boldsans}{OMX}{iwona}{bx}{n}

\newif\IfInSansMode
\let\oldsf\sffamily
\renewcommand*{\sffamily}{\oldsf\mathversion{sans}\InSansModetrue}
\let\oldmd\mdseries
\renewcommand*{\mdseries}{\oldmd\IfInSansMode\mathversion{sans}\fi\relax}
\let\oldbf\bfseries
\renewcommand*{\bfseries}{\oldbf\IfInSansMode\mathversion{boldsans}\else%
   \mathversion{bold}\fi\relax}
\let\oldnorm\normalfont
\renewcommand*{\normalfont}{\oldnorm\InSansModefalse\mathversion{normal}}
\let\oldrm\rmfamily
\renewcommand*{\rmfamily}{\oldrm\InSansModefalse\mathversion{normal}}

% Colores

\definecolor{50}{HTML}{FFEBEE}
\definecolor{300}{HTML}{E57373}
\definecolor{500}{HTML}{F44336}
\definecolor{700}{HTML}{D32F2F}
\definecolor{900}{HTML}{B71C1C}

%---------------------------------------------------------------------------
%   OPCIONES PERSONALIZADAS
%---------------------------------------------------------------------------

% Redefinir letra griega épsilon.
\let\epsilon\upvarepsilon

% Formato de texto.
\linespread{1.1}            % Espaciado entre líneas.
\setlength\parindent{0pt}   % No indentar el texto por defecto.
\setlist{leftmargin=.5in}   % Indentación para las listas.

% Estilo de página.
\pagestyle{fancy}
\fancyhf{}
\geometry{left=3cm,right=3cm,top=3cm,bottom=3cm}   % Márgenes y cabecera.

% Ruta donde buscar gráficos
\graphicspath{{../Recursos/Plantillas/}, {Recursos/Plantillas/}}

% Redefinir entorno de demostración (reducir espacio superior)
\makeatletter
\renewenvironment{proof}[1][\proofname] {\vspace{-15pt}\par\pushQED{\qed}\normalfont\topsep6\p@\@plus6\p@\relax\trivlist\item[\hskip\labelsep\it#1\@addpunct{.}]\ignorespaces}{\popQED\endtrivlist\@endpefalse}
\makeatother

% Estilo de las cabeceras

\titleformat{\section}
  {\Large\bfseries\sffamily}{\thesection}{1em}{}
\titleformat{\subsection}
  {\large\sffamily}{\thesubsection}{1em}{}
\titleformat{\subsubsection}
  {\normalsize\sffamily}{\thesubsubsection}{1em}{}

% Estilo de enlaces
\hypersetup{
  % hidelinks = true,   % Oculta todos los enlaces.
  colorlinks = true,   % Muestra todos los enlaces, sin bordes alrededor.
  linkcolor=black,     % Color de enlaces genéricos
  citecolor={500!50!black},   % Color de enlaces de referencias
  urlcolor={500!80!black}     % Color de enlaces de URL
}

% Aumentar el tamaño del interlineado
\linespread{1.3}
%---------------------------------------------------------------------------
%   COMANDOS PERSONALIZADOS
%---------------------------------------------------------------------------

% Números enteros: \ent
\providecommand{\ent}{\mathbb{Z}}

% Números racionales: \rac
\providecommand{\rac}{\mathbb{Q}}

% Números naturales: \nat
\providecommand{\nat}{\mathbb{N}}


% Valor absoluto: \abs{}
\providecommand{\abs}[1]{\lvert#1\rvert}

% Fracción grande: \ddfrac{}{}
\newcommand\ddfrac[2]{\frac{\displaystyle #1}{\displaystyle #2}}

% Texto en negrita en modo matemática: \bm{}
\newcommand{\bm}[1]{\boldsymbol{#1}}

% Línea horizontal.
\newcommand{\horrule}[1]{\rule{\linewidth}{#1}}

% Restricción de una aplicación.
\newcommand\restr[2]{{% we make the whole thing an ordinary symbol
  \left.\kern-\nulldelimiterspace % automatically resize the bar with \right
  #1 % the function
  \vphantom{\big|} % pretend it's a little taller at normal size
  \right|_{#2} % this is the delimiter
  }}

% Imagen de una aplicación.
\DeclareMathOperator*{\img}{img}

% Divisores de un elemento de un anillo
\DeclareMathOperator*{\rdiv}{div}

%---------------------------------------------------------------------------
%   CABECERA Y PIE DE PÁGINA
%---------------------------------------------------------------------------

\renewcommand{\sectionmark}[1]{%
\markboth{#1}{}}

\renewcommand{\subsectionmark}[1]{%
\markright{#1}{}}

%\addtolength{\headheight}{4ex}
\renewcommand{\headrulewidth}{0pt}

\addtolength{\headwidth}{\marginparsep}
%\addtolength{\headwidth}{\marginparwidth}

\fancypagestyle{section}{%
  \fancyhead{}
  %\addtolength{\headheight}{-10ex}
  %\renewcommand{\headheight}{0pt}%
  %\setlength{\footskip}{-48pt}%
  \fancyfoot[LE,RO]{\Large\sffamily\thepage}
  \renewcommand{\headrulewidth}{0pt}%
  \renewcommand{\footrulewidth}{0pt}%
}

\let\originalsection\section
\RenewDocumentCommand{\section}{som}{%
  \IfBooleanTF{#1}
    {\originalsection*{#3}}
    {\IfNoValueTF{#2}
      {\originalsection{#3}}
      {\originalsection[#2]{#3}}%
    }%
  \thispagestyle{section}%
}

% Texto del pie de página.
\fancyhead[LE,RO]{\rule[-4ex]{0pt}{2ex}\sffamily\textsl{\rightmark}}
\cfoot{}                                                 % Centro.
\fancyhead[LO,RE]{\sffamily{\leftmark}}
\fancyfoot[LE,RO]{\Large\sffamily\thepage}

%---------------------------------------------------------------------------
%   ENTORNOS PARA MATEMÁTICAS
%---------------------------------------------------------------------------

% Nuevo estilo para definiciones.
\newtheoremstyle{definition-style} % Nombre del estilo.
{}               % Espacio por encima.
{}               % Espacio por debajo.
{}                   % Fuente del cuerpo.
{}                   % Identación.
{\bf\sffamily}                % Fuente para la cabecera.
{.}                  % Puntuación tras la cabecera.
{.5em}               % Espacio tras la cabecera.
{\thmname{#1}\thmnumber{ #2}\thmnote{ (#3)}}     % Especificación de la cabecera (actual: nombre en negrita).

% Nuevo estilo para notas.
\newtheoremstyle{remark-style}
{10pt}
{10pt}
{}
{}
{\itshape \sffamily}
{.}
{.5em}
{}

% Nuevo estilo para teoremas y proposiciones.
\newtheoremstyle{theorem-style}
{}
{}
{}
{}
{\bfseries \sffamily}
{.}
{.5em}
{\thmname{#1}\thmnumber{ #2}\thmnote{ (#3)}}

% Nuevo estilo para la demostración
%\renewenvironment{proof}{{\itshape \sffamily Demostración \\}}{\hspace*{\fill}\qed}

\makeatletter
\renewenvironment{proof}[1][\proofname] {\par\pushQED{\qed}\normalfont\topsep6\p@\@plus6\p@\relax\trivlist\item[\hskip\labelsep\itshape\sffamily#1\@addpunct{.}]\ignorespaces}{\popQED\endtrivlist\@endpefalse}
\makeatother

% Configuración general de mdframe, los estilos de los teoremas, etc
\mdfsetup{skipabove=12pt,skipbelow=12pt,innertopmargin=12pt,innerbottommargin=4pt}

% Creamos los 'marcos' de los estilos

\mdfdefinestyle{nth-frame}{
	linewidth=2pt, %
	linecolor= 500, %
	%linecolor=black,
	topline=false, %
	bottomline=false, %
	rightline=false,%
	leftmargin=0pt, %
	innerleftmargin=1em, %
	rightmargin=0pt, %
	innerrightmargin=0pt, % %
	%splittopskip=\topskip, %
	%splitbottomskip=\topskip, %
}%

\mdfdefinestyle{nprop-frame}{
	linewidth=2pt, %
	linecolor= 300, %
	%linecolor= gray,
	topline=false, %
	bottomline=false, %
	rightline=false,%
	leftmargin=0pt, %
	innerleftmargin=1em, %
	innerrightmargin=0em,
	rightmargin=0pt, %
	%splittopskip=\topskip, %
	%splitbottomskip=\topskip, %
}%

\mdfdefinestyle{ndef-frame}{
	linewidth=2pt, %
	linecolor= 300, %
	%linecolor= gray!50,
	backgroundcolor= 50,
	%backgroundcolor= gray!5,
	topline=false, %
	bottomline=false, %
	rightline=false,%
	leftmargin=0pt, %
	innerleftmargin=1em, %
	innerrightmargin=1em,
	rightmargin=0pt, %
	innertopmargin=1.5em,%
	innerbottommargin=1em, %
	splittopskip=\topskip, %
	%splitbottomskip=\topskip, %
}%

% Asignamos los marcos a los estilos

\surroundwithmdframed[style=nth-frame]{nth}
\surroundwithmdframed[style=nprop-frame]{nprop}
\surroundwithmdframed[style=nprop-frame]{ncor}
\surroundwithmdframed[style=ndef-frame]{ndef}

% Nuevo estilo para ejemplos.
\newtheoremstyle{example-style}
{10pt}
{10pt}
{}
{}
{\bf \sffamily}
{:}
{.5em}
{}

% Teoremas, proposiciones y corolarios.
\theoremstyle{theorem-style}
\newtheorem{nth}{Teorema}[section]
\newtheorem{nprop}{Proposición}[section]
\newtheorem{ncor}{Corolario}[section]
\newtheorem{lema}{Lema}[section]
% Definiciones.
\theoremstyle{definition-style}
\newtheorem{ndef}{Definición}[section]

% Notas.
\theoremstyle{remark-style}
\newtheorem*{nota}{Nota}

% Ejemplos.
\theoremstyle{example-style}
\newtheorem{ejemplo}{Ejemplo}[section]

% Listas ordenadas con números romanos (i), (ii), etc.
\newenvironment{nlist}
{\begin{enumerate}
\renewcommand\labelenumi{(\emph{\roman{enumi})}}}
{\end{enumerate}}

% División por casos con llave a la derecha.
\newenvironment{rcases}
  {\left.\begin{aligned}}
  {\end{aligned}\right\rbrace}



%---------------------------------------------------------------------------
%   PÁGINA DE TÍTULO
%---------------------------------------------------------------------------

% Título del documento.
\newcommand{\subject}{Álgebra I}

% Autor del documento.
\newcommand{\docauthor}{Doble Grado de Informática y Matemáticas}

% Título
\title{
  \normalfont \normalsize
  \textsc{Universidad de Granada} \\ [25pt]    % Texto por encima.
  \horrule{0.5pt} \\[0.4cm] % Línea horizontal fina.
  \huge \subject\\ % Título.
  \horrule{2pt} \\[0.5cm] % Línea horizontal gruesa.
}

% Autor.
\author{\Large{\docauthor}}

% Fecha.
\date{\vspace{-1.5em} \normalsize Curso 2016/17}


%---------------------------------------------------------------------------
%   COMIENZO DEL DOCUMENTO
%---------------------------------------------------------------------------
\begin{document}
\maketitle
\tableofcontents
\vfill
\begin{center}
\includegraphics{by-nc-sa}  % Licencia.
\end{center}
\newpage

\section{Anillo conmutativo}

\begin{ndef}[Monoide]
  Un monoide es un par $(X, \cdot)$, donde $X$ es un conjunto y $\cdot$ es una aplicación $\cdot: X\times X\to X$ con las siguientes propiedades:
  \begin{nlist}
  \item Asociatividad: $a\cdot (b \cdot c) = (a \cdot b) \cdot c \ \forall a,b,c \in X$
  \item Existencia de elemento neutro: $\exists e \in X : a \cdot e = e \cdot a \ \forall a \in X$
  \end{nlist}

  En ocasiones, si la operación $\cdot$ queda clara por el contexto, se dirá simplemente que ``X es un monoide''.
\end{ndef}

Nótese que la definición de $\cdot$ implica la \emph{clausura} de la misma, es decir, para cualesquiera $a,b \in X$, $a\cdot b \in X$.

Un grupo se define similarmente, como un conjunto con una operación y ciertas propiedades.

\begin{ndef}[Grupo]
  Un grupo es un monoide $(X, \cdot)$, con la siguiente propiedad adicional:

  \begin{itemize}
  \item Existencia de elementos inversos: $\forall a\in X \ \exists a^{-1} : a\cdot a^{-1} = a^{-1}\cdot a = e$, donde $e$ es el elemento neutro de $X$ como monoide con $\cdot$.
  \end{itemize}

\end{ndef}

\begin{ndef}[Monoide conmutativo]
  Un monoide $(X, \cdot)$ se dice \emph{conmutativo} o \emph{abeliano} si $\cdot$ es una aplicación simétrica.
\end{ndef}

\begin{ndef}[Grupo conmutativo]
  Un grupo $(X, \cdot)$ se dice \emph{conmutativo} o \emph{abeliano} si $\cdot$ es una aplicación simétrica.
\end{ndef}


\begin{ndef}[Anillo]

  Un anillo es una tripla $(X, +, \cdot)$ tal que:

  \begin{nlist}
  \item $(X, +)$ es un grupo conmutativo.
  \item $(X, \cdot)$ es un monoide.
  \item $\cdot$ es distributiva a izquierda y a derecha respecto de $+$:
    \[
    \begin{cases}
      (a + b) \cdot c = (a \cdot c) + (b \cdot c)\\
      c \cdot (a + b) = (c \cdot a) + (c \cdot b)
    \end{cases} \ \forall a,b,c \in X
    \]

  \end{nlist}

\end{ndef}

\begin{ndef}[Anillo conmutativo]
  Un anillo $(X, +, \cdot)$ se dice \emph{conmutativo} o \emph{abeliano} si $(X, \cdot)$ es un monoide conmutativo.
\end{ndef}

Nótese que en este caso, la distributividad de $\cdot$ a izquierda y a derecha son dos propiedades equivalentes.

\begin{nota}
  En todas las estructuras provistas de una operación de las notadas con $\cdot$, se notará $a \cdot b = ab$.
\end{nota}


\section*{Caracterización de $\mathbb{Z}_{n}$}

Si $a,n\in \mathbb{Z}$, entonces existen $q,r \in \mathbb{Z},\ 0 \le r < n$ tales que:

\[
    a = qn + r
\]

Llamaremos $R_n:\mathbb{N} \rightarrow \mathbb{Z}_n$ a la aplicación definida como:

\[
R_n(a) = a - qn = a- nE\left(\frac{a}{n}\right)
\]
Para esta aplicación, observamos las siguientes propiedades:

\begin{itemize}
\item Si $0 \leq a < n  \Rightarrow R_n(a) = a$
\item $\forall a,b \in \mathbb{N}$

\begin{itemize}
	\item  $R_n(a+b) = R_n(R_n(a) + R_n(b))$
	\item $R_n(ab) = R_n(R_n(a)*R_n(b))$
\end{itemize}

\end{itemize}


Una vez observadas estas propiedades de la aplicación $R_n$, definimos las suma y el producto de $\mathbb{Z}_n$.
\begin{ndef}[Suma y producto en $\ent_n$] Se define la suma y el producto en $\ent_n$ de la forma:
	\begin{itemize}
	\item $a\oplus b = R_n(a+b)$
	\item $a\otimes b = R_n(ab)$
\end{itemize}

\end{ndef}



Es fácil verificar que $\mathbb{Z}_n$ es un anillo conmutativo con estas operaciones.\\

\begin{ndef}[Unidad]
  Si A es un anillo conmutativo (a.c.) $a \in A$ es ``una unidad'' o ``invertible'' si $\exists a^{-1}$ tal que $ aa^{-1} = 1$.

  El conjunto de las unidades de $A$ se nota $U(A) = \{ a \in A : a \text{ es una unidad} \}$.
\end{ndef}



\begin{ndef}[Cuerpo] Se dice que $A$ es un cuerpo si siendo un anillo conmutativo, $U(A) = A - \{0\}$, es decir, $\forall a \in A \ \exists a^{-1}$  con $a \neq 0$.
\end{ndef}


\begin{nprop}[Asociatividad generalizada]
	Sea A un anillo conmutativo, y $a_1, ..., a_n$ una lista de elementos de A.
La propiedad de la \textbf{asociatividad generalizada} nos dice que:
$\forall m$ tal que $1 \leq m < n$ se verifican: \[\sum_{i=1}^{n}a_i = \left(\sum_{i=1}^{m}a_i\right) + \left(\sum_{i=m+1}^{n}a_i\right)\]\[\prod_{i=1}^{n}a_i = \left(\prod_{i=1}^{m}a_i\right)\left(\prod_{i=m+1}^{n}a_i\right)\]\\
\end{nprop}


\begin{nprop}[Distributividad generalizada]
Se tiene la propiedad análoga a la asociatividad generalizada para la distributividad:
	\[\left(\sum_{i=1}^{n} a_i\right)\left(\sum_{j=1}^{m}b_j\right) = \sum_{i=1}^{n}\sum_{j=1}^{m}a_i b_j \hspace{1cm} \forall a,b\in A\]
\end{nprop}

\begin{ndef}[Subanillo]

Si $A$ es un anillo conmutativo y $B \subseteq A$, se dice que $B$ es un subanillo de $A$ ($B \leq A$) si $\left(B, \restr{+}{B\times B}, \restr{\cdot}{B\times B}\right)$ es un anillo conmutativo.

Esto equivale a:
\begin{itemize}
\item $1,-1 \in B$
\item B es cerrado para la suma y el producto.
\end{itemize}

\end{ndef}

\newpage

\subsection*{Anillos de números cuadráticos}

\begin{itemize}
\item $\mathbb{Z}[\sqrt{n}]$.
Definimos este conjunto de la siguiente forma:
\begin{center}$\mathbb{Z}[\sqrt{n}] = \{a+b\sqrt{n} \in \mathbb{C}$ : $a,b \in \mathbb{Z} \}$ $\leq \mathbb{C}$\end{center}

Podemos definir también $\mathbb{Q}[\sqrt{n}]$ de la misma forma:

\begin{center}$\mathbb{Q}[\sqrt{n}] = \{a+b\sqrt{n} \in \mathbb{C}$ : $a,b \in \mathbb{Q} \}$ $\leq \mathbb{C}$\end{center}

Se puede comprobar que $\mathbb{Z}[\sqrt{n}]$ $\leq$ $\mathbb{Q}[\sqrt{n}]$ y que $\mathbb{Q}[\sqrt{n}]$ es un cuerpo.

\begin{ndef}[Conjugado]
	Si $\alpha = a+b\sqrt{n}\in \mathbb{Q}[\sqrt{n}]$ se define su conjugado como $\bar{\alpha} = a - b\sqrt{n}$. Este verifica que:
\begin{enumerate}

\item $ \overline{(\alpha+ \beta)} $ = $\bar{\alpha} + \bar{\beta}$
\item $\overline{\alpha \beta}$ = $\bar{\alpha}\bar{\beta}$
\item $\alpha = \bar{\alpha} \Leftrightarrow b = 0$

\end{enumerate}
\end{ndef}

\begin{ndef}[Norma]
	Si $\alpha \in \mathbb{Z}[\sqrt{n}]$ o $\alpha \in \mathbb{Q}[\sqrt{n}]$, se define la norma $N(\alpha) = \alpha \bar{\alpha} = a^2 - nb^2 \in \rac $. Así:
\begin{enumerate}
	\item $N(\alpha \beta) = N(\alpha) * N(\beta)$
	\item $N(\alpha) = 0 \iff \alpha = 0$
\end{enumerate}
\end{ndef}


\begin{nprop}
	$\alpha = a + b \sqrt{n} \in \ent[\sqrt{n]}$ es invertible $\iff$ $N(\alpha) \in \{-1,1\}$
\end{nprop}

\end{itemize}

\subsection*{Anillos de series.}
\begin{ndef}
  Si $A$ es un anillo conmutativo y $x$ es un símbolo que no denota ningún elemento de $A$, el anillo de series con coeficientes en $A$, denotado con $A[[x]]$ esta definido como:

\[
    A[[x]] = \left\{a = \displaystyle\sum_{n\in \mathbb{N}} a_n x^n , a_n \in A\ \forall	n \in \mathbb{N}\right\}
\]

Y definimos la suma y el producto de la siguiente forma:\\
\[(a+b) = \sum_{n\in \mathbb{N}}(a_n+b_n)x^n\]

\[ (ab) = \sum_{n\in \mathbb{N}} \left(\sum_{k=0}^n a_k b_{n-k}\right)x^n\]

\end{ndef}

Se puede probar que con estas operaciones de suma y producto, $A[[x]]$ es un anillo.

\section{Homomorfismos}

\begin{ndef}
Si $(A, +, \cdot), (B, \ddagger, \ast)$ son anillos conmutativos, una aplicación $\varphi: A \to B$ es un homomorfismo si:

\begin{enumerate}
	\item $\varphi(1_A) = 1_B$, donde $1_A$ y $1_B$ son los elementos neutros de $\cdot$ en $A$ y $\ast$ en $B$, respectivamente.
	\item $\varphi(a+b) = \varphi(a) \ddagger \varphi(b)$
	\item $\varphi(ab) = \varphi(a)  \varphi(b)$

\end{enumerate}

\end{ndef}

Además, decimos que:

\begin{enumerate}
	\item Es monomorfismo si es inyectivo.
	\item Es epimorfismo si es sobreyectivo.
	\item Es isomorfismo si es biyectivo.
\end{enumerate}

\textbf{Propiedades de los homomorfismos}

\begin{itemize}

\item $\varphi(0) = 0$

\item $\varphi(-a) = -\varphi(a) $

\item $\varphi(\sum_{i = 1}^n a_i) = \sum_{i = 1}^n\varphi(a_i)$.

$\varphi(\prod_{i = 1}^n a_i) = \prod_{i = 1}^n\varphi(a_i)$

\item $\varphi(na) = n\varphi(a)$

\item $\varphi(a^n) = \varphi(a)^n$

\end{itemize}


Estas propiedades nos dicen que $\img(\varphi) = \{ \varphi(x): x \in A\} \leq B$ es un subanillo.

\begin{nprop}
	Si $\varphi$ es monomorfismo, entonces la aplicación restringida:
\[
A \to \img(\varphi)\]
\[
a \mapsto \varphi(a)
\]

es un epimorfismo y por ello es un isomorfismo, podemos decir que $A \cong \img(\varphi)$.\\
\end{nprop}


\begin{nota}
	Se puede probar que $R_n: \ent \to \ent_n$ es un homomorfismo, llamado \emph{Homomorfismo de reducción módulo n}\\
\end{nota}

\begin{ndef}[Anillo de polinomios sobre un anillo conmutativo]
  Sean $A$ un anillo conmutativo y $x$ un símbolo que no denote ningún elemento de $A$. Entonces se define el anillo de polinomios sobre $A$ (también ``con coeficientes en $A$'') como:

  \[
       A[x] = \{ a_nx^n + \dots + a_1x + a_0 : n \in \mathbb{N}, a_i \in A \ \forall i \in \{0, \dots, n\} \}
  \]

  Las operaciones de suma y producto en $A[x]$ son las usuales. Con las convenciones de que $x^0 = 1, x^1 = x$ y $\sum_{i=m}^n a_i = 0$ si $m>n$:

  \[
  \begin{array}{lllll}
    \sum_{i=0}^n a_ix^i + \sum_{i=0}^m b_ix^i & = & \sum_{i=0}^m (a_i + b_i)x^i + \sum_{i=m+1}^n a_ix^i, \text{ con $m \le n$.} & & \\
    \left(\sum_{i=0}^n a_ix^i\right) \cdot \left(\sum_{i=0}^m b_ix^i\right) & = & \sum_{i=0}^n \sum_{j=0}^m a_i b_j x^{i+j} =  \sum_{k=0}^{m+n} \sum_{i+j=k} a_ib_j x^k
  \end{array}
  \]

\end{ndef}

$A[x]$ es un subanillo de $A[[x]]$ (los elementos de $A[x]$ son los de $A[[x]]$ para los que $a_n = 0$ a partir de un cierto $n\in \mathbb{N}$).

\begin{nprop}[Homomorfismo de cambio de coeficientes)(1]
 Dado $A$ cualquier anillo conmutativo, conocido $A[x]$.

 Si $\varphi:A \to B$ es un homomorfismo de anillos conmutativos, entonces:
 \[
 \exists \psi: A[x] \to B[x] : \psi\left(\sum_i a_i x^i\right) = \sum_i\varphi(a_i) x^i
 \]


\end{nprop}

\begin{nprop}[Sustitución en un polinomio)(2]

Si $A$ es un anillo y $a \in A$ entonces: existe un homomorfismo $E_a: A[x] \to A$ tal que $E_a(\sum_i a_i x^i) = \sum_i a_i a^i$.

\end{nprop}

\begin{nprop}[3]
Si $A \leq B$ es un subanillo y $b\in B$, la aplicación $E_b:A[x] \to B$ definida como $E_b(\sum_i a_i x^i) = \sum_i a_i b^i$ es un homomorfismo.

\end{nprop}

\begin{nprop}[Engloba a las anteriores]
Si $\varphi:A \to B$ es un homomorfismo y $b\in B$, la aplicación $\Phi:A[x] \to B$ definida como $\Phi ( \sum_i a_i x^i) = \sum_i \varphi(a_i)b^i \in B$ es un homomorfismo.

\end{nprop}

\begin{proof}
Veamos primero cómo (4) engloba a las demás:
\begin{nlist}

\item $4 \Rightarrow 3$. Se ve tomando como $\varphi$ la inclusión en $B$
\item $4 \Rightarrow 2$. Tomamos esta vez como $\varphi$ la identidad
\item $4 \Rightarrow 1$. Suponemos 4 válido. Probaremos que $\exists\psi: A \to B[x]$ que lleva $a \mapsto \psi(a) = \varphi(a)$.
  En efecto, basta tomar $\psi = i_{B[x]} \circ \varphi$, donde $i_{B[x]} : B \to B[x]$ es la inclusión en $B[x]$.


De esta forma, tomamos $x\in B[x]$. Entonces:

\[
\begin{array}{lll}

  A[x] & \to & B[x] \\
  \sum_i a_i x^i & \mapsto & \sum_i \varphi(a_i)x^i
\end{array}
\]

es un homomorfismo, que es justamente el enunciado de la primera proposición.

\end{nlist}

Pasamos ahora a la demostración de la proposición 4.


Sean $f = \sum a_i x^i$ y $g = \sum b_i x^i \in A[x]$.
Si ahora aplicamos $\Phi(f+g) = \sum \varphi(a_i + b_i)b^i$, como $\varphi$ es homomorfismo, eso es igual a:
\[
    \sum (\varphi(a_i) + \varphi(b_i))b^i
\]

Usando que $B$ es un anillo y por ello hay distributividad, eso es igual a

\[
    \sum (\varphi(a_i)b^i + \varphi(b_i)b^i).
\]

Por la asociatividad generalizada, es igual a:

\[
\sum \varphi(a_i)b^i + \sum \varphi(b_i)b^i = \Phi(f) + \Phi(g)
\]

Por lo que queda probado para la suma.

Ahora probaremos el producto. Recordemos que, de la definición:

\[
fg = \sum c_k x^k \text{ con $c_k$ }  = \sum_{i+j = k} a_ib_j
\]

 Así: \[\Phi(f*g) = \sum_n\varphi(c_n)b^n = \sum \varphi\left(\sum_{i+j = n} a_ib_j\right)b^n = \sum_n \left( \sum_{i+j = n} \varphi(a_i)\varphi(b_j)\right)b^n\]

 Desarrollamos por otro lado

 \[
 \begin{array}{>{\displaystyle}l}
   \Phi(f) * \Phi(g) = \left(\sum_i \varphi(a_i)b^i\right)\left(\sum_j \varphi(b_j)b^j\right) =^{(1)} \sum_{i,j} \varphi(a_i)b^i\varphi(b_j)b^j \\
   =^{(2)} \sum_{i,j} \varphi(a_ib_j)b^{i+j} = \sum_n\left(\sum_{i,j: i+j = n} \varphi(a_ib_j)b^n\right)
 \end{array}
 \]

 Donde en $(1)$ hemos usado la distributividad general y en $(2)$ hemos usado que estamos en un anillo conmutativo y que $\varphi$ es un homomorfismo.

 Hemos llegado a dos expresiones que son iguales, probando así el resultado.

\end{proof}


Sabemos que cada polinomio $f(x)$ constituye una función de evaluación $f(x) \in A[x]$
\[
f(x):B \to B
\]
\[
\hspace{1cm}b \mapsto f(b)
\]

Sin embargo, un polinomio es mucho más que la función de evaluación que él mismo define.

Estudiaremos el anillo de polinomios en varias variables.

\begin{ndef}[Monomio en $r$ variables]
  Sean $A$ un anillo conmutativo, $p_1, \dots, p_r \in \mathbb{N}$ y $x_1, \dots, x_r$ símbolos que no denoten ningún elemento de $A$.
  Entonces

  \[
      \prod_{i=1}^r x_i^{p_i}
  \]

  es un monomio en $r$ variables con coeficientes en $A$. Si $p = (p_1, \dots, p_r)$, denotamos $X^p = \prod_{i=1}^r x_i^{p_i}$.
\end{ndef}

\begin{ndef}[Anillo de polinomios en $r$ variables sobre un anillo conmutativo]
  Sean $A$ un anillo conmutativo y $x_1, \dots, x_r$, con $r \in \mathbb{N}$ símbolos que no denoten ningún elemento de $A$. Entonces, se define el anillo de polinomios en
  $r$ variables sobre $A$ como

  \[
      A[x_1, \dots, x_r] = \left\{ \sum_{i=1}^n b_iX^{p_i} :
      \begin{array}{l}
        B = \{b_1, \dots, b_n\} \subseteq A \\
        P = \{p_1, \dots, p_n\} \subset \mathbb{N}^r \\
        n \in \mathbb{N}
      \end{array} \right\}
   \]

      Es decir, el conjunto de las combinaciones lineales finitas con coeficientes en $A$ de monomios en $r$ variables sobre $A$.
\end{ndef}

Este anillo de polinomios es igual al que se puede obtener tomando anillos de polinomios sobre anillos de polinomios sucesivamente,
es decir:

\[
    A[x_1, \dots, x_{s}][x_{s+1}, \dots, x_{r}] = A[x_1, \dots, x_r]
\]

para cualquier número de variables $x_1, \dots, x_r$ y cualquier permutación de ellas.

Ilustraremos este hecho en el caso $r=2$, $A[x][y]$.

Definimos

\[
f=\sum f_i y^i,\ f_i = \sum_j a_{ij}x^j
\]

Luego

\[
f = \sum_i \left(\sum_j a_{ij}x^j \right)y^i = \sum_{i,j} a_{ij}x^iy^j
\]

Como $A$ es un anillo conmutativo, se tiene que $A[x][y] = A[y][x] = A[x,y]$.

Después de esto, vamos a demostrarlo en el caso general.

\begin{nprop}
  Sea $A$ un anillo conmutativo. Entonces

  \[
     A[x_1, \dots, x_{r-1}][x_r] \text{ y } A[x_1, \dots, x_r]
  \]

   son isomorfos.
\end{nprop}

\begin{proof}
  Si $f \in A[x_1, \dots, x_{r-1}][x_r]$,

  \[
       f = \sum_{i=0}^n f_i x_r^i,\ f_i = \sum_{j\in J} a_{ij} X^j,\ J = \{j_0, \dots, j_m\} \subset \mathbb{N}^{r-1}
  \]

  Luego

  \[
      f = \sum_{(i,j) \in \{1,\dots,n\}\times J} a_{ij}X^j x_r^i
  \]

  Definimos $\Phi : A[x_1, \dots, x_{r-1}][x_r] \to A[x_1, \dots, x_r]$:

  \[
    \Phi \left(\sum_{(i,j) \in \{1,\dots,n\}\times J} a_{ij}X^j x_r^i\right) = \sum_{i=0}^{nm} b_i X^{p_i},\
    \begin{array}{l}
      B = \{ b_k = a_{ij} : jn + i = k \} \\
      P = \{ p_k = q_{ij} : q_{ij} \in \{1, \dots, n \} \times J, jn+i = k \}
    \end{array}
  \]

  Se puede comprobar que este es un isomorfismo de anillos conmutativos.

\end{proof}


\begin{nprop}

Si $\varphi :A \to B$ es un homomorfismo, $\forall(b_1,...,b_n) \in B^n$ la aplicación:

\[
\Phi:A[x_1,...x_n] \to B \iff \Phi\left(\sum_{i1,...,in}a_{i1}...a_{in}x^{i1}... x^{in}\right) = \sum_{i1,...,in}a_{i1}...a_{in}b^{i1}... b^{in} \in B
\]

es un homomorfismo de anillos conmutativos. Es conocido como evaluación de un polinomio en n variables.

\end{nprop}



\begin{nprop}
	Si $\varphi:A \to B$ es un homomorfismo, $\forall\  b \in B$ existe un único homomorfismo definido como:

\[
\Phi:A[x] \to B : \begin{cases}
	\Phi(a) = \varphi(a)\ \forall a \in A \\
	\Phi(x) = b
\end{cases}
\]
\[
\Phi\left(\sum a_i x^i\right) = \sum \Phi(a_i x^i) =  \sum \Phi(a_i) \Phi(x)^i = \sum \varphi(a_i) b^i
\]

Además, ya se probó que esto es un homomorfismo de anillos conmutativos.

\end{nprop}

\begin{ncor}
	$A \leq B$ subanillo, $\forall b \in B$ existe un único homomorfismo
\[
E_b: A[x] \to B: \begin{cases}
E_b(a) = a \hspace{0.3cm} \forall a \in A\\
E_b(x) = b

\end{cases}
\]
\end{ncor}

\begin{nota}
	En la anterior definición, tendremos en cuenta que $x$ es un elemento concreto de $A[x]$ (el polinomio de grado uno con coeficientes cero y la unidad), no un elemento cualquiera de $A[x]$.
\end{nota}

\begin{nota}
	Si $f(x)\in A[x]$ denota un polinomio de $A[x]$, notaremos: $E_b(f(x)) = f(b)$.
	De la misma forma, si $f(x) = \sum a_i x^i \Rightarrow E_b(f(x)) = \sum a_i b^i$
\end{nota}


\begin{nprop}[Evaluación en $r$ variables]
  Si $\varphi : A \longrightarrow B$ es un homomorfismo de anillos conmutativos, y $b_1,\cdots, b_r \in B$ una lista ordenada. Entonces
  \[
    \exists!\ \phi : A[x_1,\cdots,x_r] \longrightarrow B :\begin{cases}
      \phi(a) = \varphi(a) & \forall a\in A\\
      \phi(x_1) = b_1\\
      \vdots\\
      \phi(x_r) = b_r
    \end{cases}
  \]
\end{nprop}

\begin{proof}
  Si $r=1$, ya está probado. Para $r>1$:
  \[
    \exists\ \psi : A[x_1,\cdots,x_{r-1}] \longrightarrow B : \begin{cases}
      \psi(a) = \varphi(a)\\
      \psi(x_i) = b_i & \forall i=1,\cdots,r-1
    \end{cases}
  \]
  \[  \exists\ \phi : A[x_1,\cdots,x_r] \longrightarrow B \begin{cases}
      \phi(a) = \psi(a) = \varphi(a)\\
      \phi(x_i) = \psi(x_i) = b_i &\forall i=1,\cdots,r-1\\
      \phi(x_r) = b_r
    \end{cases}
  \]

  ¿Es único?
  \[
    \phi\left(\sum_{i1,\cdots,ir} a_{1i}\cdots a_{ir}x_1^{i1}\cdots x_r^{ir} \right)= \sum _{i1,\cdots,ir} \varphi(a_{i1}\cdots a_{ir})b_1^{i1}\cdots b_r^{ir}
  \]
\end{proof}

\begin{nprop}[Evaluación en subanillos $r$ variables]
	Si $A\leq B, \forall b_1,\cdots,b_r \in B$ lista ordenada:

\[
\exists !\ E_{b_1,\cdots,b_r}: A[x_1,\cdots,x_r] \to B :
\begin{cases}
	a \mapsto a\\
	x_i \mapsto b_i
\end{cases}
\]

Se suele notar $f(x_1,\cdots,x_r) \to f(b_1,\cdots,b_r)$
\end{nprop}

\section{Dominio de integridad}

\begin{ndef}[Dominio de integridad]
Un anillo conmutativo $A$ es un dominio de integridad si verifica la propiedad:
\[
a \neq 0\ \wedge\ b \neq 0 \Rightarrow ab \neq 0 \ \ \forall a,b \in A
\]

Equivalentemente,

\[ \hspace{0.1cm} ab=0 \implies \begin{cases}
	a = 0,\text{ o}\\
	b = 0
\end{cases} \ \forall a,b \in A
\]

\end{ndef}

\begin{nprop}[Propiedad de simplificación]

A es un dominio de integridad $\iff$ se verifica: $ax=ay$ con $a\neq 0 \Rightarrow x = y$
\end{nprop}
\begin{proof}\hfill\\
	\boxed{\Rightarrow} $a(x-y) = 0$, por ser A dominio de integridad, $x-y = 0 \Rightarrow x=y$ \\
	\boxed{\Leftarrow} $ab = 0$ con $a\neq 0$, tomo $x,y \in A : b = x-y$, entonces $ab = a(x-y) = 0 \implies x=y \implies b = 0$.
\end{proof}

\begin{ndef}[Divisor de 0]
	$a\in A$ es divisor de 0 si $\exists b\neq 0 : ab = 0$.
\end{ndef}

\begin{nprop}
	Si A es un dominio de integridad $\Rightarrow$ el 0 es el único divisor de 0.

	Equivalentemente: A es dominio de integridad $\iff$ no tiene divisores de cero no nulos.
\begin{nlist}
	\item $A\leq B$ y $B$ es D.I. $\Rightarrow$ A es D.I.
	\item Todo cuerpo es D.I.
	\item Si $u \in U(A) \Rightarrow$ u no es divisor de 0 (supongamos $u*b = 0 \Rightarrow u*u^{-1}*b = u^{-1}*0 \Rightarrow b=0$)
\end{nlist}
\end{nprop}

\begin{nprop}
	Si $A$ es finito, $A$ es dominio de integridad $\iff$ $A$ es un cuerpo
\end{nprop}
\begin{proof}\hfill\\
	\boxed{\Leftarrow} Trivial\\
	\boxed{\Rightarrow} $0 \neq a \in A$.
Tomo $\{1,a,a^2,\dots\} = \{a^n: n\in \mathbb{N}\} \subseteq A$
Como $A$ tiene cardinal finito y es cerrado para el producto: $\exists n,k \in \mathbb{N} : a^n = a^{n+k}$.\\
Pero, por ello: $a^n = a^n a^k; a^n * 1 = a^n * a^k$, luego $a^n$ no es 0 porque $A$ es dominio de integridad y por ser D.I. entonces:
\[
1=a^k \begin{cases}
	k = 1 \Rightarrow a = 1\\
	k > 1 \Rightarrow a^{k-1}*a = 1

\end{cases}\]
Con lo que $\exists$ inverso de $a = a^{k-1}$ y como $a$ es un elemento cualquiera, todo elemento tiene inverso, luego es un cuerpo.
\end{proof}


\begin{nprop}
	Todo D.I. es un subanillo de un cuerpo.
\end{nprop}


Antes de probar esta proposición, presentaremos otros conceptos. Comenzaremos por repasar brevemente los conceptos de relación binaria y
relación de equivalencia, y definir el cociente de un conjunto por una relación de equivalencia. Esto nos servirá
para construir el cuerpo de fracciones de un dominio de integridad, del que el dominio de integridad será un subanillo.

\begin{ndef}[Relación binaria]
  Una relación binaria en un conjunto $A$ es un subconjunto de $A^2$.
  Si $\sim \ \subseteq A^2$ es una relación binaria y $(a,b) \in \sim$, se nota $a \sim b$
  y se dice que ``$a$ está relacionado con $b$ mediante $\sim$''.
\end{ndef}

\begin{ndef}[Relación de equivalencia]
  Una relación de equivalencia en un conjunto $A$ es una relación binaria $\sim$ en $A$ que verifica las siguientes propiedades:

  \begin{nlist}
  \item \textbf{\emph{Reflexividad}}. $a \sim a \ \ \forall a \in A$.
  \item \textbf{\emph{Transitividad}}. Si $a \sim b$ y $b \sim c$, entonces $a \sim c \ \ \forall a,b,c \in A$.
    \item \textbf{\emph{Simetría}}. $a \sim b \implies b \sim a \ \ \forall a,b \in A$.
  \end{nlist}
\end{ndef}

\begin{ndef}[Clase de equivalencia]
  Sean $A$ un conjunto, $a\in A$ y $\sim$ una relación de equivalencia en $A$. Se define la clase de equivalencia
  de $a$ respecto de $\sim$:

  \[
     [a] = \{ b \in A : a \sim b \}
  \]
\end{ndef}

\begin{ndef}[Cociente de un conjunto por una relación de equivalencia]
  Sean $A$ un conjunto y $\sim$ una relación de equivalencia en $A$. Entonces se define el cociente de $A$ por $\sim$:

  \[
     A/\sim \ = \{ [a] : a \in A\}
  \]
\end{ndef}

\begin{ndef}[Cuerpo de fracciones de un D.I.]
  Sea $A$ un dominio de integridad con $|A| \geq 2$. Consideramos

  $$A\times (A-\{0\}) = \{(a,b), a,b \in A\ |\ b\neq 0\}$$

  y en este conjunto la relación de equivalencia:

  \[
     (a,b) \sim (c,d) \iff ad = bc
  \]

  Definimos el cuerpo de fracciones de $A$ como

  $$\mathbb{Q}(A) = [A\times (A-\{0\})]/\sim$$

    Se nota $[(a,b)] = \displaystyle\frac{a}{b}$ y se le llama ``fracción $a$ entre $b$''.

  Sobre él, definimos unas operaciones con las que será un cuerpo:
\begin{nlist}
	\item Suma: \[\frac{a}{b}+ \frac{c}{d} = \frac{ad + cb}{bd}\]

Ahora, como la fracción $\dfrac{a}{b}$ es un conjunto, hay que probar que el resultado es único, es decir:
\[
\frac{a}{b}= \frac{a'}{b'} \text{ y } \frac{c}{d}= \frac{c'}{d'}  \Rightarrow ab' = a'b \text{ y } cd' = c'd
\]
Hay que probar que se cumple:
\[
\frac{ad+cb}{bd} = \frac{a'd'+c'b'}{b'd'}
\]
Equivalementente, tenemos que probar que se cumple:
\[
b'd'(ad+cb) = bd(a'd'+c'b')
\]

Desarrollamos en la izquierda:
\[
b'd'(ad+cb) = b'd'ad + b'd'cb =^{(1)} a'bd'd + b'bc'd
\]
Donde en (1) hemos usado la equivalencia que habíamos dado de $ab' = a'b$ y $cd' = c'd$. Ahora, desarollamos el producto de la derecha y veremos que es igual al resultado obtenido

\[
bd(a'd'+c'b') = bda'd' + bdc'b' = a'bdd' + bb'c'd
\]

Probando la unicidad.

    \item Producto: \[\frac{a}{b} \frac{c}{d} = \frac{ac}{bd}\]

La unicidad del producto se hace desarrollando de la misma manera.
\end{nlist}



Para finalizar, se puede probar que es un cuerpo probando las propiedades de anillo conmutativo y que existe inverso para todo $\dfrac{a}{b}$.
\end{ndef}

\begin{ncor}
	\[\frac{a}{b}= \frac{u}{v} \iff av = bu \iff (a,b) \sim (u,v)\]
\end{ncor}

\begin{nprop}[Fracciones de denominador 1]
  Si $A$ es un dominio de integridad, existe un homomorfismo

  \[
  \begin{array}{llll}
    i :& A & \longrightarrow & \mathbb{Q}(A) \\
       &a & \longmapsto &  i(a) = \dfrac{a}{1}
  \end{array}
  \]

  Que cumple que $i(a+b) = i(a)+i(b)$ y que $i(ab)=i(a)i(b)$, y además es un monomorfismo.
  Así, $A \stackrel{i}{\cong} \img(i)$ es un isomorfismo y $A\le\mathbb{Q}(A) \text{ con } a=\dfrac{a}{1}$.
  Con esta identificación $\dfrac{a}{b} = \dfrac{a}{1}\dfrac{1}{b} = ab^{-1}$.
\end{nprop}

A partir de ahora usaremos la identificación $a = \dfrac{a}{1}$.

\begin{nprop}
  Sea $K$ un cuerpo y $A\le K$, $a,b\in A\ (b\not=0)$.
  \begin{align}
    \implies& a\in K \text{ y } b^{-1}\in K \implies ab^{-1}\in K\notag\\
    \implies&\mathbb{Q}(A)\le K\notag
  \end{align}
\end{nprop}

\begin{nota}
 Sea $K$ un cuerpo. Entonces  $\mathbb{Q}(K)$ es el cuerpo más pequeño que contiene a $K$.
\end{nota}

\begin{nota}
  $A \subseteq \mathbb{Q}(A), A=\text{D.I.} \implies \rac(\rac(A)) = \rac(A)$
\end{nota}


\begin{nprop}
  Sea K un cuerpo, $A\le K$. Si $\forall\alpha\in K\quad \exists a\in A,\;\; a \not= 0 : a\alpha\in A \implies \rac(A)=K$
\end{nprop}

\begin{proof}
  $\alpha\in K,\; \exists a \ne 0, a\in A : a\alpha=b\in A \implies \alpha = ba^{-1} = \frac{b}{a}\in \rac(A)$
\end{proof}

\begin{ejemplo}
  $\ent[i] = \{a+bi : a,b\in \ent\} \le \rac[i] = \{a+bi : a,b\in \rac\} \implies \rac(\ent[i]) = \rac[i]$
  \[
    \alpha\in\rac[i]\implies\alpha=\frac{m}{n}+\frac{m'}{n'}i\; \implies \ent[i] \ni nn'\alpha = n'm+nm'i \in \ent[i]\notag
  \]
\end{ejemplo}

\begin{nprop}
	Si $A$ es un D.I. $\implies$ $A[x]$ es un D.I.
\end{nprop}

\begin{ndef}[Grado de un polinomio]

  Si $f=\sum a_i x^i \neq 0$, entonces

  $$gr(f)=n\in \mathbb{N} : a_n\neq 0 \text{ y } a_m=0 \hspace{0.3cm} \forall m>n$$

El coeficiente $a_n$ se denomina coeficiente líder.

\end{ndef}

\begin{nprop}
	Si A es D.I., $f,g \in A[x]\implies gr(fg) = gr(f)+gr(g)$.
	En caso contrario, tenemos que $gr(fg) \le gr(f)+gr(g)$.
\end{nprop}

\begin{ndef}[Divisibilidad en D.I.]

Sea $A$ un D.I. Sean $a,b\in A$. Decimos entonces que $a$ divide a $b$ ($a$ es un divisor de $b$, $b$ es un múltiplo de $a$) si
\[
	\exists c \in A : b = ac
\]
\end{ndef}

\begin{nprop}
  En las condiciones anteriores:

\begin{align}
	&\exists c \in A : b = ac
\\
	\iff& \text{La ecuación } ax=b \text{ tiene solución}
\\
\iff& \frac{b}{a} \in A
\end{align}
\end{nprop}

\begin{proof}\hfill\\
  \begin{itemize}
  \item $\boxed{(1) \iff (2)}$ Es la definición de que la ecuación $ax = b$ tenga solución.
  \item $\boxed{(1) \iff (3)}$

    $\boxed{\Longrightarrow}$ Si $a$ divide a $b$ $\implies$ $\exists c: b=ac \implies \dfrac{b}{a} = \dfrac{ac}{a} = \dfrac{c}{1} = c \in A$\\
    $\boxed{\Longleftarrow}$ si $\dfrac{b}{a}\in A \implies \dfrac{b}{a} = \dfrac{c}{1} \implies b = ac$
  \end{itemize}

\end{proof}

\textbf{Notación:} Si $a$ divide a $b$, escribiremos $a/b$.

\begin{ndef}[Asociados de un elemento de un D.I.]
  Sean $A$ un dominio de integridad y $a\in A$. Entonces los elementos $ua$ con $u\in U(A)$ se llaman \emph{asociados de a}.
\end{ndef}

\begin{nprop}\hfill
  \begin{nlist}
  \item Los divisores de 1 son las unidades del anillo, los elementos del grupo $U(A)$.
  \item Las unidades son divisores de todos los elementos del anillo.
  \item Si $u\in U(A)$, $\forall a \in A$, $ua/a$.

  \end{nlist}
\end{nprop}

\begin{ndef}
	Los divisores triviales de un número son las unidades y sus asociados.
\end{ndef}

\begin{nprop}
	Sean $a,b\neq 0$. Son equivalentes:
	\begin{nlist}
	\item $a$ es asociado de $b$.
	\item $b$ es asociado de $a$.
	\item $a/b \wedge b/a$, los asociados son los elementos que se dividen mutuamente.

\end{nlist}
\end{nprop}

\begin{nprop} La divisibilidad verifica algunas propiedades:

  \begin{nlist}
  \item Reflexión: $a/a$
  \item Transitividad: $a/b \wedge b/c \implies a/c$
  \item Si $a/b \wedge a/c \implies a/bx +cy\hspace{0.2cm} \forall x,y \in A$
  \item Si $a/b \implies \hspace{0.2cm} \forall c\in A \hspace{0.2cm} a/bc$
  \item Si $c\neq 0$ entonces $a/b \iff ac/bc$
  \end{nlist}
\end{nprop}

\textbf{Notación:} si $a$ y $b$ son asociados, escribiremos $a \sim b$.

\begin{ndef}[Irreducible]

	Sean $A$ un dominio de integridad, $a\in A$, $a\neq 0, a \notin U(A)$. Se dice que $a$ es irreducible si sus únicos divisores son los triviales

\begin{align*}
	\iff& \text{si } b/a \implies b\in U(A) \vee b \sim a
\\
	\iff& \text{si } a=bc \implies b\in U(A) \vee c \in U(A)
\\
	\iff& \text{si } a=bc \implies a\sim b \vee c \sim a
\\
	\iff& \text{si } a=bc \wedge b \notin U(A) \implies c \in U(A)
\end{align*}
\end{ndef}


\section{Dominios euclídeos}
\begin{ndef}[Dominios euclídeos]

	Un dominio euclídeo es un dominio de integridad, $A$, tal que se pueda definir una función $\varphi: A- \{0\} \to \mathbb{N} $ que verifique:

\begin{nlist}
	\item $\varphi(ab) \ge \varphi(a)$
	\item $\forall a,b \in A,\ b \neq 0 \quad \exists q,r \in A : a = bq + r $ con $r=0 \vee \varphi(r) < \varphi(b)$
	%%\item $\forall a,b \in A,\ b \neq 0 \quad \exists q \in A : a - bq = 0 \vee  \varphi(a-bq) < \varphi(b)$
\end{nlist}
\end{ndef}

\begin{nota}
  En algunos textos, se omite la primera propiedad, dado que a partir de una función que verifique la segunda, es posible definir una que verifique ambas.
\end{nota}


\begin{nprop}
  Si A es dominio euclídeo, entonces:

  \[
  \begin{array}{lll}
    & b/a & (1) \\
    \iff & \text{un resto de dividir $a$ entre $b$ es cero} & (2) \\
    %%\iff & \text{cualquier resto de dividir $a$ entre $b$ es cero} & (3)
  \end{array}
  \]
\end{nprop}

\begin{proof}\hfill\\
  %%\begin{itemize}
  %%\item $\boxed{ (1) \iff (2) }$

    $\boxed{\Longleftarrow}$ Trivial.

  $\boxed{\Longrightarrow}$ Por definición de $b/a$, $\implies \exists c \in A$ tal que $a=bc$ y por ser $A$ un dominio euclídeo, $\exists q,r \in A : a = bq + r $
  con $r=0 \vee \varphi(r) < \varphi(b)$. La solución 	es evidentemente correcta para $r = 0$, veamos que sucede para $r \neq 	0$.

    Supongamos $r \neq 0$, entonces $\varphi(r) < \varphi(b)$.
    $$r = a - bq = bc - bq = b(c-q) \quad \quad c-q \neq 0$$
    $$\varphi(r) = \varphi(b(c-q)) \ge \varphi(b)$$
    Llegando a una contradicción.
 %% \end{itemize}

\end{proof}

\begin{nth}[Teorema de Euclides]

	$\forall a,b \in \ent, b \neq 0, \exists !q,r \in \ent$ tales que $a = bq + r $ con $0 \le r < |b|$
\end{nth}


\begin{proof}
Probaremos primero la unicidad. Supongamos
$$a=bq+r \quad \quad 0\leq r <\abs{b}$$
$$a=bq'+r' \quad \quad 0 \leq r' < \abs{b}$$

distintos. Vamos a ver que $r=r'$ y $q=q'$

\begin{itemize}
	\item Si $r \neq r'$ ($\implies q\ne q'$), supongamos $r > r' \implies 0 < r - r' < \abs{b}$
Ahora:
$$ r-r' = a-bq - a +bq' = b(q' -q )$$
$$r-r' > 0 \implies r - r' = \abs{b(q' -q)} = \abs{b}\abs{q' -q}$$

Pero, como $q\neq q' \implies q'-q \neq 0$ y $q,q' \in \ent \implies \abs{q' -q} \geq 1 $, luego:

$$ r-r' = \abs{b}\abs{q' -q} \geq \abs{b}$$

Por lo que tenemos una contradicción con el comienzo de la suposición.

	\item Ahora, si $r=r'\implies b(q' -q) = 0$ y $b\neq 0 \implies q' -q=0 \implies q' =q$
\end{itemize}

Probamos ahora la existencia. Observemos que el caso general se puede reducir al caso $a \ge 0, b > 0$, de la
siguiente forma:

\begin{itemize}
\item si $b<0$, entonces se toman $b' = -b, q' = -q$, y se escribe $a = b'q' +r, 0 \le r < |b'|$.
\item si $a<0$ y $b>0$ (podemos asegurar que $b>0$ por el caso anterior), se toman $a' = -a$, $q' = -q -1, r' = b-r$
  y se escribe $a = bq'+r', 0 \le r' < b$.
\end{itemize}

Por tanto, sean $a,b \geq 0$:

\begin{itemize}
	\item Si $ a < b \implies a = b*0 + a$, luego $q=0$ y $r=a$, ya los tenemos.
	\item Si $a \geq b$, llamamos $R=\{a-bx: x\in \nat \quad a \geq bx\}\subseteq \nat$ que es no vacío, pues está al menos $ a-b \ (\text{para } x=1)$.

	Ahora, por el principio de buena ordenación, $R$ tiene mínimo. Tomo $r=\min(R)$.

$r=a-bq$ para cierto $q \in \nat$ y $ r \geq 0$.

Veremos que $r < b$, por contradicción.

Supongamos $ r \geq b\implies r' = r-b \geq 0 \implies r' = a-bq -b = a -b(q+1)\implies r' \in R$.

Podemos ver que $r' < r$ (pues $r'=r-b$) $\implies  r'$ está en $R$ y es menor que el mínimo, luego es una contradicción y tenemos que $r < b$.

\end{itemize}
%% Por último, vamos a probar que $0 \leq r < \abs{b}$

%% Supongamos:
%% \[
%% r=0 \implies a=bq \begin{cases}
%% 	-a = b(-q)\\
%% 	-a = (-b)q \\
%% 	 a = (-b)(-q)\\
%% \end{cases}
%% \]

%% Ahora, supongamos $r > 0$:

%% \begin{itemize}
%% 	\item $-a = b(-q) -r = b(-q) -b + b -r = b(-q-1) +(b-r)$ y como $0 < r < b \implies b > b-r > 0$
%% 	\item $-a = (-b)q -r = (-b)q + b - b -r = -b(q+1) + (b-r)$ y por el mismo motivo, $b > b-r > 0$
%% 	\item $a = (-b)(-q) +r \implies 0 < r < b = \abs{-b}$
%% \end{itemize}

%% De esta forma, hemos cubierto todos los casos y hemos acabado la demostración
\end{proof}

\begin{ncor}
$\ent$ es un dominio de Euclides con
  \[
  \begin{array}{lll}
  \varphi & = & \abs{\;\cdot\;}: \ent \to \nat \\
  \restr{\varphi}{\nat} & = & 1_{\nat} \\
  \restr{\varphi}{\ent \setminus \nat} & = & -1_{\ent \setminus \nat}
\end{array}
\]
\end{ncor}

\begin{nth}[Teorema de Euclides para polinomios]
		$\forall f,g \in A[x]$ donde $g \neq 0 $ y su coeficiente líder es una unidad de $A$, existen polinomios:

	\[
	q,r \in A[x] : f = gq + r \quad con \quad
	 \begin{cases}
		r = 0\\
		o\\
		gr(r) < gr(g)
	\end{cases}
	\]
	y que son únicos.
\end{nth}

%\newpage
\begin{proof}

	Sean $f = \sum_{i=0}^n a_i x^i$ y $g = \sum_{i=0}^m b_i x^i$ con $b_m \in U(A)$

	\begin{itemize}

	\item Si $n< m \implies f = g * 0 + f \implies \exists q,r\in A[x] : f = gq+r$ con $q=0$ y $r = f$.

	\item Si $n\geq m$, razonamos por inducción en $n=gr(f)$
	\begin{itemize}

	\item Si $n=0 \implies m = 0$ por tanto $f=a_0$ y $g=b_0$ con $b_0 \in U(A)$

	De esta forma:

	\[
	f = a_0 = \frac{a_0}{b_0}b_0 = \frac{a_0}{b_0}b_0 + 0 = g\frac{a_0}{b_0}
	\]

	Podemos tomar como hemos visto $q=\frac{a_0}{b_0}$ y $r=0$ y tenemos el $q$ y $r$ que buscábamos.

	\item Si $n> 0$, haremos la inducción.

	  Consideramos $x^{n-m}g(x)$ y establecemos

          $$f_1 = f - \frac{a_n}{b_m}x^{n-m}g\ \ (1)$$

	  Entonces, podemos ver que $gr(f_1) < n$. Por hipótesis de inducción

          $$ \exists q,r \in A[x] : f_1 = gq_1 + r \ \ (2)$$

	Ahora, utilizando (1) y (2):

	 \[
	 \implies f = f_1 +\frac{a_n}{b_m}x^{n-m}g = gq_1 + \frac{a_n}{b_m}x^{n-m}g +r =
	 \]

	 \[
	 g\left(q_1 + \frac{a_n}{b_m}x^{n-m}\right) + r
	 \]

	Encontramos así el $q$ y el $r$ que queríamos, probando la existencia.

	Vamos a probar ahora la unicidad.

	Sea $f=gq+r$ y $f= gq' +r'$ con
	\[
    \begin{cases}
		 r,r' \neq 0\\
		  \quad o \\
 gr(r) < m
 \\ gr(r') < m

\end{cases}
	\quad
	\]

	Ahora, si $r\neq r'\implies r-r' \neq 0 \implies r-r' = g(q-q') \neq 0$.
	Vemos que $gr(r-r') = gr(g)+ gr(q-q')$.

 Como $q-q' \neq 0 \implies gr(q-q') \geq 0 $ y de esta forma: $gr(g)+ gr(q-q') \geq gr(g) = m$.

 Sin embargo, habíamos dicho que $r,r'$ eran ambas de grado menor que m luego $gr(r-r') < m$, llegando a una contradicción y probando así el resultado.
\end{itemize}


\end{itemize}
\end{proof}

\begin{ncor}
	Si $K$ es un cuerpo, entonces $K[x]$ es un D.E. con función euclídea:

	\[
	gr: K[x]-\{0\} \to \nat
	\]

(función que asigna a cada polinomio su grado)
\end{ncor}

\begin{nota}
	Hacemos el ejercicio de ver si $3x^2 +1$ es divisor de $2x^3 + 4x^2 +4x +3$ en $\ent _5[x]$. (Solución: El resto de la división es 0, con resultado de la división $= 2/3 x + 4/3$)
\end{nota}


\begin{nth}[Teorema de Euclides, enteros cuadráticos]
	Los anillos $\ent[\sqrt{n}]$ para $n=2,3,-1,-2$ son D.E. con función euclídea:
	\[
	\varphi: \ent[\sqrt{n}] \to \nat : \varphi(a+b\sqrt{n}) = \abs{N(a+b\sqrt{n})} = \abs{a^2 - nb^2}
	\]


\end{nth}
\begin{proof}
	Probaremos que $\forall \alpha, \beta \in \ent[\sqrt{n}]$ con $\beta \neq 0 \quad \quad \exists q,r \in \ent[\sqrt{n}] : \alpha = \beta q + r$ con $r=0$ ó $\abs{N(r)} < \abs{N(\beta)}$:

	\begin{itemize}
	\item Si $\abs{N(\alpha)} < \abs{N(\beta)}$ Basta tomar $\alpha  =  \beta * 0 + \alpha$

	\item Si $\abs{N(\alpha)} \geq \abs{N(\beta)}$ consideramos entonces  $\frac{\alpha}{\beta} \in \rac[\sqrt{n}]$.

	Ahora, $\frac{\alpha}{\beta} = a_1 + a_2\sqrt{n}$ con $a_1, a_2 \in \rac $. Esos $a_1,a_2$ se obtienen usando el conjugado de $\beta$.

	Sean $q_1,q_2 \in \ent : \abs{a_1-q_1}\leq 1/2$ y $\abs{a_2-q_2}\leq 1/2$. Esto quiere decir que $q_1$ y $q_2$ son los enteros más cercanos a $a_1,a_2$ respectivamente.

	Sea $q=q_1+q_2\sqrt{n}$  y $r = \alpha - \beta q$.

	Tomo $\abs{N(r)} = \abs{N(\alpha - \beta q)} = \abs{N(\beta(\frac{\alpha}{\beta}+ q))} = \abs{N(\beta)}\abs{N(\frac{\alpha}{\beta}+ q)}$

	Queremos probar que: $\abs{N(\beta)}\abs{N(\frac{\alpha}{\beta}+ q)} < \abs{N(\beta)}$.

	Equivalentemente, queremos probar que:
	\[
	  \abs{N(\frac{\alpha}{\beta}+ q)} < 1 \implies \abs{N(a_1+a_2\sqrt{n}-q_1 -q_2\sqrt{n})} = \abs{N((a_1-q_1) + (a_2-q_2)\sqrt{n})} = \]
	  \[
	 = \abs{(a_1-q_1)^2 - n(a_2-q_2)^2} = m \in \rac
	 \]

Vamos a probarlo para los casos que habíamos anunciado en el teorema, $n = -1,-2,2,3$

\begin{itemize}
	\item $n=-1 \implies m= (a_1-q_1)^2 +(a_2-q_2)^2 \leq 1/4 + 1/4 = 1/2 \implies \abs{m} < 1$

	\item $n=-2 \implies m = (a_1-q_1)^2 +2(a_2-q_2)^2 \leq 1/4 + 1/2 = 3/4 \implies \abs{m} < 1$

	\item $n=2 \implies m = |(a_1-q_1)^2 -2(a_2-q_2)^2 |\implies -1/2 \leq m \leq 1/4 \implies \abs{m} < 1$

	\item $n=3 \implies m = |(a_1-q_1)^2 -3(a_2-q_2)^2| \implies -3/4 \leq m \leq 1/4 \implies \abs{m} < 1$
\end{itemize}

Por lo que queda probado el resultado para esos casos.

\end{itemize}
\end{proof}

\begin{ejemplo}
Vamos a tratar de dividir $\alpha=6+10i$ entre $\beta=1+2i$ en el anillo $\ent[i]$. Tenemos que saber si se puede hacer dicha división o no ($\varphi(ab) \ge \varphi(a)$) y para ello averiguaremos la norma de ambos números.

$$\abs{N(6+10i)}=36+100=136$$
$$\abs{N(1+2i)}=1+4=5$$

Como $1+2i$ tiene una norma menor que la norma $6+10i$ podemos hacer la división, primero dividiremos como si fuesen numeros complejos normales para hallar nuestro número cociente que será de la forma $q=q_1+q_2i$:

\[
	\frac{6+10i}{1+2i}=\frac{(6+10i)(1-2i)}{(1+2i)(1-2i)}=\frac{6-12i+10i+20}{5}=\frac{26-2i}{5}=\frac{26}{5}-\frac{2}{5}i
\]

Tenemos que $5<\frac{26}{5}<6$ y $5$ es más cercano a $\frac{26}{5}$ que $6$ escogemos $q_1=5$ y por el mismo razonamiento $q_2=0$, de forma que $q=5+0i=5$. A continuación, para hallar el resto $r$ hacemos la siguiente operación:

\[
	r=\alpha-\beta \cdot q = 6+10i - (1+2i)(5)=6+10i-5-10i=1
\]

Finalmente, comprobamos que no nos hemos equivocado:
\[
	(6+10i)=5(1+2i)+1\\
	\abs{N(1)}<\abs{1+2i} \implies 1<5
\]

Viéndose así que el ejemplo está correcto.

\end{ejemplo}

\section{Máximo común divisor. Dominios de ideales principales. Ecuaciones diofánticas en D.I.P.}


\begin{ndef}[Máximo común divisor]

Sea $A$ un anillo conmutativo. Dados $a,b \in A$ decimos que un elemento $d\in A$ es un \textit{mcd} de $a$ y $b$ ($d=(a,b)$) si el conjunto de los divisores comunes a $a$ y a $b$ coinciden con el conjunto de los divisores de $d$. Esto es:

\begin{nlist}
	\item $d/a$ y $d/b$
	\item Si $c/a$ y $c/b \implies c/d$
\end{nlist}

\end{ndef}

\begin{ndef}[Primos relativos]
  Si $(a,b) = 1$, $a$ y $b$ se dicen \emph{primos relativos}.
\end{ndef}

\begin{ndef}[Ideal]
  En un anillo conmutativo se llama ideal a un subconjunto suyo no vacío que es cerrado para la suma y que ``absorbe'' el producto en todo
  el anillo. Formalmente, si $A$ es un anillo conmutativo, un subconjunto $\emptyset \not= I\subseteq A$ es un ideal si
	\begin{nlist}
		\item $a,b\in I \implies a+b\in I$
		\item $a\in I, x\in A \implies ax \in I$
	\end{nlist}
\end{ndef}

\begin{ndef}[Ideal principal]
  Sea $A$ un anillo conmutativo. Si $a\in A$,

  $$aA=(a)=\{ax: x\in A\}$$

  es el ideal principal generado por $a$.
\end{ndef}

\begin{ndef}[DIP: dominio de ideales principales]
	Un dominio de ideales principales es un dominio de integridad en el cual todo ideal es principal.
\end{ndef}



En un DIP, el máximo común divisor es único salvo asociados, y verifica algunas propiedades:
\begin{nlist}

	\item $(a,b) = (b,a)$

	\item Si $a \sim a'$ asociados y $b \sim b'$ también $\implies (a,b) = (a',b')$

	\item $(a,b) = a \iff a/b$. En particular, $(a,0) = a, \quad  \quad (a,1) = 1, \quad (a,u) = 1 \iff u \in U(A)$

	\item $((a,b),c) = (a,(b,c)) = (a,b,c)$
	\item $(ac,bc) = c(a,b)$\\
	\begin{proof}
	Primero, llamamos $(ac,bc) = e$ y $(a,b) = d$.\\
	Si a,b o c son 0, se verifica trivialmente. Si no lo son:\\
	\[
	\begin{rcases}
	d/a \implies dc/ac\\
	d/b \implies dc/bc\\
\end{rcases} \implies dc/e \implies \exists u \in A : e = dcu
	\]

	$\begin{rcases}e / ac \implies \exists x \in A : ac = ex \implies ac = dcux \implies a = dux\\
	e/bc\implies \exists y \in A : bc = ey \implies bc = dcuy \implies b = duy\end{rcases} \implies \begin{rcases}du/a\\du/b\end{rcases} du/d \\\implies \exists v \in A : d = duv \stackrel{d\ne 0}{\implies} 1=uv \implies u\in U(A) \implies e \sim  dc$


\end{proof}

	\item Si $c/a$ y $c/b\implies \left(\frac{a}{c},\frac{b}{c}\right) =  \frac{(a,b)}{c}$

	\item
	$\left(\frac{a}{(a,b)}, \frac{b}{(a,b)}\right) = 1$

	\item Si $a/bc \implies a/(a,b)c$\\
	\begin{proof}
	Supongamos que $\exists x \in A : bc = ax  \implies (a,b)c = (ac,bc) = (ac,ax) = a(c,x) \implies a/(a,b)c$

\end{proof}

\item Si $a/bc$ y $(a,b) = 1 \implies a/c$

\item Si $a/c$ y $b/c$ y $(a,b) = 1 \implies ab/c$
	\item Si $(a,b)=1 $ y $a/bc \implies a/c$
	\item Si $a/c$ , $b/c$ y $(a,b)=1 \implies ab/c$
	\item Si $a/c \implies \exists x$ : $c=ax$. Y $b/c \implies b/ax$ con $(a,b)=1 \implies b/x \implies \exists y$ : $x=by$
	Entonces:
	\[
		\begin{cases}
			c=ax \\
			x=by \\
		\end{cases} \implies c=aby \implies ab/c
	\]
	\item Si $(a,b)=1$ y $(a,c)=1 \iff (a,bc)=1$ \\
	\begin{proof}
	$\boxed{\implies}$ Sabiendo que: $(ac,bc)=c(a,b)=c$

	Tenemos que: $1=(a,c)=(a,(ac,bc))=((a,ac),bc)=(a(1,c),bc)=(a,bc)$, por tanto: $1=(a,bc)$

	$\boxed{\Longleftarrow}$ $1=(a,bc)=(a(1,c),bc)=((a,ac),bc)=(a,(ac,bc))=(a,c(a,b))=(\frac{a}{(a,b)}(a,b),c(a,b))=(a,b)(\frac{a}{(a,b)},c)=1 \implies (a,b)\in U(A)\implies (a,b)=1\implies (a,c)\in U(A)\implies (a,c)=1$
\end{proof}
	\item $(a,b)=(a-kb,b)$ $\forall k \in A$
	\item Si $d/b$, $d/a \iff d/(a-kb)$ \\
	\begin{proof}
	$\boxed{\implies}$ Por la propiedad de combinación lineal se confirma.

	$\boxed{\Longleftarrow}$ Igual que la otra implicación pero tomando $a=(a-kb)+kb$
	\end{proof}
\end{nlist}

\begin{nota}
En $\mathbb{Z}[\sqrt{n}]$ si $\alpha$ es un divisor propio de $\beta \implies N(\alpha)$ es un divisor propio de $N(\beta)$ en $\mathbb{Z}$.
\end{nota}

\begin{ejemplo}
Realizamos un ejemplo en el que se puede probar que, usando la Nota anterior, 3 y $(1+\sqrt{5})$ son irreducibles
\end{ejemplo}


\begin{nth}
	Todo dominio euclídeo es un dominio de ideales principales: $DE \implies DIP$
\end{nth}
\begin{proof}
	Sea A un DE con función euclídea $\varphi: A-\{ 0 \} \longrightarrow \mathbb{N}$ y $I\subseteq A$ un ideal:
	\begin{itemize}
		\item Caso $I=\{ 0 \} = (0)= 0A \implies $ trivial
		\item Consideremos $I\neq \{ 0\}$, $\emptyset \neq \{ \varphi (x): x\in I, x\neq 0\} \subseteq \mathbb{N}$, sea $\varphi (b)$ el mínimo de este conjunto, donde $b\in I, b\neq 0 \implies I=(b)$. Probamos esto con la doble inclusión:
		  $\boxed{\subseteq}$ $b\in I \implies (b)\subseteq I$.

		$\boxed{\supseteq}$ Tomamos $a\in I$, por tanto $\exists q,r\in A: a=bq+r$. Supongamos que $r\neq 0 \implies r=a-bq\in I$ con $\varphi (r) < \varphi (b)$, esto es imposible puesto que $\varphi(b)$ es el mínimo, luego $r=0 \implies a\in (b)$.
	\end{itemize}
\end{proof}


\begin{nth}
	Si $A$ es un DIP, $\forall a,b \in A \quad \exists d = (a,b)$. Además, $\exists u,v \in A : d = au+bv$. A esta igualdad se le llama Identidad de Bezout, y u y v son los coeficientes de Bezout, que no son únicos.

\end{nth}

\begin{proof}
	Sea $\emptyset \ne I(a,b) = \{ax + by : x,y \in A\} \subseteq A$

	Vemos que:\\
$(ax+by)+(ax'+by') = a(x+x') + b(y+y') \in I(a,b) \implies$ cerrado para la suma.\\
$(ax+by)z = a(xz)+b(xz) \in I(a,b) \implies$ cerrado para el producto

Ahora, como es un ideal $\implies \exists d\in A : I(a,b) = (d)$ con $(d)=\{dx: x \in A\}$.
$d\in I(a,b) \implies \exists u,v \in A : d=au+bv$.

Ahora, veamos que $d$ es \textit{mcd} de $a$ y $b$.\\
$a\in I(a,b) \implies a \in (d) \implies d/a$\\
$b\in I(a,b) \implies b \in (d) \implies b/d$

Por lo que $d$ es divisor común. Ahora, sea $c: c/a$ y $c/b \implies c/(au+bv=d) \implies c/d$.
Hemos encontrado así un divisor común que es dividido por cualquier divisor común, por tanto es el \textit{mcd}.
\end{proof}

\subsection{Ecuaciones diofánticas en D.I.P.}

En cualquier anillo $A$, llamamos ecuaciones diofánticas a aquellas que son de la forma:
\[
ax +by  = c \quad a,b,c \in A
\]

\begin{nprop}[Lema de Euclides]
  Sean $a,b,c \in A$, con $A$ un D.I.P., tales que $c/ab$ y $(c,a)=1$. Entonces $c/b$.
\end{nprop}

A continuación vamos a estudiar las soluciones de las ecuaciones diofánticas en un D.I.P., en tres pasos.
Comenzaremos por dar una condición suficiente y necesaria para que una ecuación tenga solución,
fácilmente verificable; en ese caso, daremos una solución particular y a partir de ella, la general.

\begin{nlist}
	\item Sea $d=(a,b)\implies$ entonces la ecuación tiene solución $\iff d/c$
	\item Supongamos que tiene solución. Supongamos también que $d= au+bv$
\[
\frac{a}{d}= a',\quad \frac{b}{d}=b',\quad \frac{c}{d}=c'\implies da'x + db'y = dc' \implies d(a'x+b'y)=dc'\]
\[ \implies a'x+b'y = c'
\]
Esta ecuación tiene las mismas soluciones que la ecuación diofántica inicial. Llamaremos a esta la ecuación ``reducida''.\\


Como $c'= a'(c'u)+b'(c'v)$ y ahí tenemos una solución particular. Conociendo esta, podemos hallar todas las soluciones. Si llamamos $x_0 =c'u$ e $y_0 = c'v$

\item Solución general \[\begin{cases}
	x = x_0 +kb'\\
	y= y_0 -ka'
\end{cases} k\in A\]
Si $(x_0,y_0)$ es la solución particular, entonces la solución general es el conjunto de los $(x,y)$ que hemos dado arriba.
\\
\begin{proof}[Demostración de iii)]

\[
a'x + b'y = a'(x_0 +kb') + b'(y_0 - ka') = a'x_0 + a'kb' + b'y_0 - a'kb' =
\]
\[
a'x_0 +b'y_0 = c'
\]

Suponer ahora que $(x,y)$ es cualquier solución:$\implies a'x +b'y = c'$. Por hipótesis: $a'x_0 +b'y_0 = c'$. Si restamos esas dos ecuaciones queda: $a'(x-x_0) +b(y-y_0) = 0 \implies a'(x-x_0) = b(y_0-y)$. Denotamos a esta ecuación como 3.

Ahora, $b'/(a'(x-x_0))$ pero $b'$ y $a'$ son primos entre sí, luego $b'/(x-x_0)\implies \quad \exists k \in A : (x-x_0) = kb'$. Llamamos a esta ecuación 1, y además despejando en ella vemos $x = x_0 + kb'$, una solución de x.

Análogamente, podemos ver que $a/(b(y_0-y)) \implies a/(y_0-y) \implies \exists h \in A: y_0 -y = a'h \implies y = y_0 - ha'$, solución de y. Llamamos a esa ecuación la 2.

Falta probar que $k = h$, pero sustituyendo las ecuaciones 1 y 2 en 3, vemos que $a'kb'  = b'ha' \implies k = h$

\end{proof}

\end{nlist}

\begin{nprop}[Algoritmo de Euclides para el cálculo del MCD]
Supongamos que tenemos dos elementos $a,b$ y queremos hallar su mcd.

\begin{itemize}

	\item Si $b=0 \implies (a,b) = (a,0) = a$. Igual si $a=0$

	\item Si $a \ne 0 \ne b$

	Construimos una sucesión: $r_1,r_2,\cdots, r_n, \cdots , r_m, r_{m+1} = 0$.\\
	Recordamos que A es un D.E. con función euclídea $\varphi : A-\{0\} \to \nat $\\
	Si $\varphi(a) \geq \varphi(b)\implies r_1 = a$ y $r_2 = b$. En el otro caso, lo hacemos al revés, es decir  $r_1 = b $ y $r_2 = a$.\\
	Si $r_{n-1} \ne 0 \implies r_n= $ resto de dividir $r_{n-2}$ entre $r_{n-1} \implies$

	\[
	r_{n-2} = r_{n-1}q_{n-2} + r_n \begin{cases}
	 r_n = 0 \\
	  \varphi(r_n) \leq \varphi(r_{n-1})
\end{cases}
	 \]
La idea es ir reduciendo de la forma:
\[ (a,b) = (r_1,r_2) = \cdots = (r_n,r_{n+1}) = \cdots = (r_m, r_{m+1}) = (r_m,0) = r_m\]
Obteniendo los cocientes de la forma:
\[\begin{cases}
	r_{n-2} = au_{n-2} + bv_{n-2}\\

r_{n-1} = au_{n-1} + bv_{n-1}\\

r_{n-2} - r_{n-1}q_{n-2} = r_n = a(u_{n-2} - q_{n-2}u_{n-1}) + b(r_{n-2}-q_{n-2}v_{n-1})\\
\cdots \\

d = r_m = au + bv
\end{cases}
\quad
\]

	\end{itemize}
\end{nprop}

\begin{ejemplo}
  Un agricultor lleva al mercado 80 sandías y 30 melones. La venta le ha sido rentable, pues ha vendido cada pieza por más de 3 euros, que es lo que le costó producirlos. Vuelve a casa con 600 euros. Calcular precio de sandías y melones.

	(El ejercicio se resuelve resolviendo la ecuación diofántica $80x +30y = 600$ , hallando primero la solución general que viene dada por $x = -60 + 3k$ ; $y = 180 - 8k$ y luego tomando que $x$ e $y$ tienen que ser mayores que 3, viendo que la solución es que $k=22$).
\end{ejemplo}

\section{Mínimo común múltiplo. Ecuaciones en congruencias}
\begin{ndef}[Mínimo común múltiplo]
	Sean $A$ un D.I., y $a, b \in A$. Entonces $m$ es un mínimo común múltiplo de $a$ y $b$,y lo notamos $m = mcm(a,b) = [a,b]$,
si se verifica que $m$ es un múltiplo común de de $a$ y $b$ y que el conjunto de los múltiplos comunes a ambos es igual al conjunto de múltiplos de $m$, formalmente:
\begin{enumerate}
\item $a/m \text{ y } b/m$.
\item Si $a/c$ y $b/c \Rightarrow m/c$.
\end{enumerate}

\textbf{Propiedades. }
\begin{nlist}
\item Si $a\sim a' \ y \ b\sim b' \Rightarrow [a,b]=[a',b']$
\item $[a,b] = [b,a]$
\item $[a,0] = 0$
\item $[a,1] = a$
\item $[a,[c,b]]= [[a,c],b] = [a,b,c]$
\item $[ac,bc] =[a,b]c$\\

\begin{proof}[Demostración del último.]

	Supongamos que $c\ne 0$, pues si no es trivial.\\
	Como $c/ab\implies c/[ca,cb] \implies \exists q \in A : [ac,bc] = cq\quad (1)$\\
	Por otro lado, sea $m = [a,b]; \implies a/m \quad y \quad b/m \implies ac/mc \quad y \quad bc/mc \implies cq/mc$.\\
	Como $c\ne 0 \implies q/m$.\\
	Por otro lado, $ca/cq \quad y \quad cb/cq \implies$ como $c\ne 0 \implies a/q \quad y \quad b/q \implies m/q.$\\
	Hemos llegado a que $q/m \quad y \quad m/q \implies $ son asociados $\implies q=[a,b]$.\\
	Ahora, basta llevarnos esto a $(1)$ en esta demostración para ver que:\\ $[ac,bc] = c[a,b]$
\end{proof}
\end{nlist}
\end{ndef}


\begin{nprop}
	Si A es un DIP $\implies \forall a,b \in A \quad \exists [a,b]$
\end{nprop}
\begin{proof}
	Consideramos $aA = (a)$, el ideal principal generado por a. De la misma forma, consideramos $bA = (b)$, el ideal principal generado por b.\\
	Ahora, tomamos $aA \cap bA \implies$ los números que están simultáneamente en los múltiplos de ambos.
	Ahora, esto es cerrado para sumas y para productos, por tanto también es un ideal.\\
	Por último, por estar en un DIP $\implies$ el ideal es principal y por tanto:\\ $\implies aA\cap bA = mA \implies m = [a,b]$
\end{proof}


\begin{nth}
	Sea A un DI en el cual $\exists(a,b) \quad \forall a,b \in A.$ Entonces, $\exists [a,b] \quad \forall a,b \in A $ y se verifica que: $[a,b](a,b) = ab$
\end{nth}
\begin{proof}
	Sean $0 \ne a,b \in A$. Llamamos $d=(a,b) \implies \begin{cases}
	a = a_1 d\\
	b = b_1 d
\end{cases}$

Podemos observar que:

\[
m = \frac{ab}{d} = a_1b = ab_1
\]
De esta forma, nuestra prueba termina si comprobamos que $m = [a,b]$. Tenemos ya que claramente $a/m \quad y \quad b/m$.\\

Sea $m_1 = a/m_1 $ y $b/m_1$, tenemos que probar que $m/m_1$. Para esto, lo que hay que probar es que $(m,m_1) = m$.\\
Para ello, vamos a llamarlo $k = (m,m_1) \implies k/m$. Llamo $d_1 = \frac{m}{k} \implies m =_{(1)} d_1k$ para un cierto $d_1$. Guardamos la igualdad de $(1)$ para usarla después.\\
Ahora, lo que bastaría probar es que $d_1 \in U(A)$:\\
Tenemos que $a/m$ y $a/m_1 \implies a /k \implies k = au$. Podemos hacer lo mismo con b para ver que $k = bv$. Esto ocurre para ciertos $u$ y $v$.\\
Ahora, usando la igualdad del principio ($m = a_1b = ab_1$) y el $(1)$ podemos ver que $\begin{rcases}
	m = a_1b = kd_1 = bvd_1 \implies a_1 = vd_1\\
m = ab_1 = kd_1 = aub_1  \implies  b_1 = ud_1
\end{rcases} \implies \begin{rcases}
	a = a_1d = vd_1 d\\
b = b_1d = ud_1d
\end{rcases} \implies d_1d/a \quad d_1d/b$


$\implies d_1d /d \implies \exists x \in A : d = dd_1 x \implies 1 = d_1 x \implies d_1 \in U(A)$.\\
$\implies m,k$ son asociados y como k era $mcd(m,m1) \implies$ m también lo es.
\end{proof}

\subsection{Congruencias}
Sean $A$ un anillo, $I \subset A$ un ideal. $a,b \in A$ son ``congruentes módulo I'' si $a-b \in I$. Equivalentemente, si $\exists x \in I : a = b+x$. La notaremos:
\[
a \equiv b \mod(I) \quad \text{ o } \quad a \equiv_I b
\]
Otra notación. En un DIP $I = (m) = mA$
\[
a \equiv b \mod(mA) \to^{\text{notación}} a \equiv b \mod(m)\ (\iff m/a-b \iff a-b = qm)
\]
Para algún $q$ en el último paso, y en ese caso $\iff a = b+qm$ para algún $q$.

\textbf{Propiedades}

\begin{nlist}
	\item $\equiv$ es una relación de equivalencia.
	\begin{itemize}
		\item $a \equiv a$
		\item $a \equiv b \iff b \equiv a$ (dem.:\ $a-b = (-1)(b-a) \in I$)
		\item $a \equiv b$ y $b \equiv c \implies a \equiv c$ (dem.:\ $a-b\in I, b-c \in I \implies a-c \in I$)
	\end{itemize}

	\item $a\equiv b \iff \forall c : a+ c \equiv b + c$
	\item $a\equiv b$ y $c \equiv d \implies a+c \equiv b + d$ (dem.: usando $(ii)$ y $(i)$)
	\item $a \equiv 0 \iff a \in I$
	\item $a \equiv b \implies \forall c\in A : ac  \equiv bc$
	\item $a \equiv b$, $c \equiv d \implies ac \equiv bd$ (dem.:\ $(v)$ y luego uso\ $(i)$)
	\item $ac \equiv bc \mod(mc) $ y $c \ne 0 \implies a\equiv b \mod(m)$\\

\begin{proof}
	$ac \equiv bc \mod(mc) \iff mc/(a-b)c \iff^{c\ne 0} m/a-b \iff a\equiv b \mod(m)$
\end{proof}
	\item Si $(c,m) = 1$, entonces: $ac\equiv bc\mod(m) \iff a \equiv b\mod(m)$\\
	\begin{proof}
	$ac\equiv bc \implies m/(a-b)c \implies$, como $(c,m) = 1 \implies m/a-b$
\end{proof}


\end{nlist}

\subsection{Ecuaciones en congruencias}
\begin{nprop}[Ecuaciones en congruencias]

  Estudiaremos la ecuación

  $$ax \equiv b \mod(m) \  \ (1)$$

\begin{itemize}
	\item Si $m = 0 \implies$ la ecuación es $ax = b$.
	\item Si $a = 0\implies $ la ecuación es $0x \equiv 0 \mod(m) \implies$ tiene solución: todo el anillo.
	\item $a,b \ne 0$.
	\begin{enumerate}

	\item Si $d = (a,m)$ la ecuación tiene solución $\iff  d/b$. \\
        \begin{proof}
	(1) tiene solución $\iff \exists x \in A : ax \equiv b\mod(m) \iff \exists x \in A : m/ax-b \iff \exists x,y \in A : (ax-b) = my \iff \exists x,y \in A: ax -my  = b$, que es una ecuación diofántica, que sabemos ya que tiene solución $\iff d =(a,m) \ y  \ d/b$.
\end{proof}
	\item  Supongamos que tiene solución. Consideramos $a'= \frac{a}{d}$, $b'= \frac{b}{d}$ y $m' = \frac{m}{d}$.

	Ahora, usando (1), $da'x \equiv db'\mod(dm')$, esta es equivalente a $a'x \equiv b' \mod(m')$ a la que llamaremos (2). Esta es su reducida. Tiene las mismas soluciones pero $(a',m') = 1$.

	Podemos hallar los coeficientes de Bezout: $u,v \in A :  1 = a'u + b'v$. Esto nos lleva a ver que:
	\[
	a'u \equiv 1 \mod(m') \implies a'ub' \equiv b' \mod(m')
	\]
	Y así tenemos que $x_0 = ub'$ es una solución particular.

	\item La solución general es de la forma: $x = x_0 +km' \quad k \in A$. Equivalentemente, es de la forma $x \equiv x_0 \mod(m')$\\
	\begin{proof}
	Si $x_0$ es una solución particular $\implies a'x_0 \equiv b'\mod(m')$\\
	Si sustituimos $x_0$ por $x$ pues son congruentes obtenemos: $a'x \equiv b' \mod(m')$.\\
	Vamos a suponer que:
	\[
	\begin{rcases}
	a'x \equiv b' \mod(m')\\
	a' x_0 \equiv b' \mod(m')
\end{rcases} \implies a'x \equiv a'x_0 \mod(m')
	\]
Por la transitividad. Pero $a'$ y $m'$ son primos entre sí, luego $x \equiv x_0 \mod(m')$
\end{proof}

\item Diremos que una solución particular $x_1$ es óptima si $x_1 = 0$ ó $\varphi(x_1) < \varphi(m')$ siendo $\varphi$ la función euclídea de A.\\
Si $x_0$ es cualquier solución particular, entonces:
\[
x_0 = m'q + x_1 \begin{cases}
	x_1 = 0\\
	o\\
	\varphi(x_1) < \varphi(m')

\end{cases}
\]
Y $x_1$ es una solución parcial óptima. En este caso, la solución general óptima es: $x\equiv x_1 \mod(m')$
\end{enumerate}
\end{itemize}

\end{nprop}

\subsection{Sistemas de ecuaciones en congruencias}

En este caso, vamos a abordar un problema en el que tenemos un sistema de ecuaciones en congruencias, que sabemos que se puede expresar de la forma:

\[
\begin{rcases}
	a_1x \equiv b_1 \mod(m1)\\
	a_2x \equiv b_2 \mod(m2)
\end{rcases}(1) \begin{rcases}
	x \equiv a \mod(m)\\
	x \equiv b \mod(n)
\end{rcases}(2)
\]

\begin{nth}[Teorema chino del resto]
	El sistema tiene solución $\iff a \equiv b \mod((m,n))$.
\end{nth}
\begin{proof}
	Sea $d=(m,n)$. Si tomamos $x = a+km; \quad \exists k: a+km \equiv b \mod(n) \iff km \equiv b-a \mod(n) \iff d/b-a \iff b\equiv a\mod(d)$.
\end{proof}

Ahora, supuesto que tiene solución, vamos a hallar las soluciones particular y general del problema.\\
Si $y_0$ es una solución particular de $my \equiv b-a \mod(n)$, entonces su solución general es: $$y = y_0 + k\frac{n}{(m,n)} \quad k \in A$$
Entonces $x_0 = a + my_0$ es una solución particular del sistema dado en (2) y por tanto la solución general de 2 viene dada por:

\[
x = a +m(y_0 +  k\frac{n}{(m,n)}) \quad k \in A
\]
\[
 = a+my_0 + k \frac{mn}{(m,n)} = x_0 +k[m,n] \quad k \in A
\]
\[
\implies x \equiv x_0 \mod [m,n]
\]
Pero si $x_0 = [m,n]q + x_1$ con $x_1 = 0 $ ó $\varphi(x_1) < \varphi([m,n])$ entonces tenemos que $$x_0 \equiv x_1 \mod([m,n])$$

Y obtenemos que la solución general óptima de nuestro sistema es:
\[
x \equiv x_1 \mod([m,n])
\]


\begin{nth}[Teorema de Ruffini]
	Si $f(x) \in A[x], \ a \in A$ entonces $f(a)$ es el resto de dividir $f$ entre $x-a$.
	Equivalentemente, $f = (x-a)q+r$ donde $r\in A$. Así, $f(a) = r$.

	En forma de congruencias: $f \equiv f(a)\mod(x-a)$.
\end{nth}
\section{Anillos de congruencias. Anillo cociente.}


\begin{ndef}[Clase de congruencia respecto de un ideal]
  Sean $A$ un anillo conmutativo, $I \subseteq A$ y $a \in A$. Si definimos la relación de equivalencia

  \[
    a \circeq b \iff a - b \in I
  \]

  es decir, ``ser congruente módulo $I$'', entonces la clase de congruencia de $a$ respecto de $I$ es la clase de equivalencia de $a$ respecto de $\circeq$. La notamos

  \[
    [a] = \bar{a} = a+I
  \]
\end{ndef}



\begin{ndef}[Anillo cociente]
  Sean $A$ un anillo conmutativo e $I$ un ideal suyo. Entonces se define el anillo cociente de $A$ por $I$, $(A/I, \oplus, \odot)$:

  \[
  \begin{array}{l}
    A/I = { [a] = a+I : a \in A} \\
    %% Estas llaves son necesarias, sin ellas se invoca un error oscuro de LaTeX: https://www.sharelatex.com/learn/Errors/Missing_number,_treated_as_zero
    {[a]} \oplus [b] = [a+b] \\
    {[a]} \odot [b] = [ab]
  \end{array}
  \]
\end{ndef}

Se puede comprobar que estas operaciones están bien definidas y que con ellas $A/I$ es un anillo conmutativo. Probaremos que están bien definidas:

Si $[a] = [a']$ y $[b] = [b'] \implies \begin{cases}
	[a+b] = [a' +b']\\
	[ab] = [a'b']
\end{cases}$

\begin{proof}
  \[
  a\equiv_I a' \text{ y } b \equiv_I b' \implies \begin{cases}
	a+b \equiv_I a'+b'\\
	ab \equiv_I a'b'
        \end{cases}\implies \begin{cases}
	[a+b] = [a'+b']\\
	[ab] = [a'b']
  \end{cases}
  \]
\end{proof}


\begin{nprop}
  Si $f:A \to B$ es un homomorfismo de anillos, $\ker(f) = \{a \in A: f(a) = 0\}$ es un ideal de $A$. \\

  \begin{proof}
   Vamos a probar que $\ker(f)$ es cerrado para sumas y para múltiplos. Para ello, en ambos casos usaremos que $f$ es un homomorfismo.

  \[
     \text{Si } a,b \in \ker(f) \implies
     \begin{cases}
       f(a+b) = f(a)+f(b) = 0 + 0 = 0 \\
       f(ab) = f(a)f(b) = 0*0 = 0
     \end{cases}
  \]
\end{proof}

Además, $f$ es un monomorfismo $\iff \ker(f) = 0$
\begin{proof}\hfill\\
	\boxed{\Rightarrow} Trivial\\
	\boxed{\Leftarrow} Si $f(a) = f(b) \implies f(a-b) = 0 \implies a-b \in \ker(f)$ pero hemos dicho que $\ker(f) = 0 \implies a-b = 0 \implies a = b$.
\end{proof}
\end{nprop}

\begin{nth}[Teorema de isomorfía]
  Si $f:A \to B$ es un homomorfismo, se induce un isomorfismo de anillos:

  \[
  \begin{array}{lrll}
    F: &  A/\ker(f) & \cong & \img(f) \\
       & [a] & \mapsto & f(a)
  \end{array}
  \]


\begin{proof}
  Para probar el resultado tenemos que probar tres cosas:

  \begin{nlist}
  \item $F$ está bien definida.
  \item $F$ es biyectiva.
  \item $F$ es un homomorfismo de anillos.
  \end{nlist}

  Vamos a probar que está bien definida. Es decir, que sin importar el representante de la clase de congruencia que tomemos,
  el valor de la aplicación para esa clase es el mismo. Esto es, $[a] = [b] \implies F([a]) = F([b])$.

  Si $[a] = [b] \implies a-b \in \ker(f) \implies f(a-b) = 0 \implies F([a]) = f(a) = f(b) = F([b])$.

  Vamos a ver ahora que es un homomorfismo
  \begin{itemize}
		\item $F([a]+[b]) = F([a+b]) = f(a+b) = f(a)+f(b) = F([a]) + F([b])$
		\item $F([a][b]) = F([ab]) = f(ab) = f(a)f(b) = F([a])F([b])$
		\item $F([1]) = f(1) = 1$
  \end{itemize}

Probamos la inyectividad:

$F([a]) = F([b]) \implies f(a) = f(b) \implies f(a-b) = 0 \implies a-b \in \ker(f) \implies a \equiv b\mod(\ker(f)) \implies [a] = [b]$.
%Tomo $[a] \in \ker(F) \implies F[a] = 0 \implies f(a) = a \implies a \in \ker(f) \implies [a] = [0] = 0$

La sobreyectividad es trivial por la sobreyectividad de $f$ en su imagen.

\end{proof}

\end{nth}

\begin{nprop}
  Sea $A$ un dominio euclídeo con función euclídea $\varphi :A-\{0\} \to \nat$ tal que en $A$ hay unicidad de cocientes y restos. Esto es:

  $$\forall a,b \in A: b\ne 0 \implies \exists ! q,r\in A : a=bq+r \text{ y } \begin{cases}
	r = 0 \text{ o} \\
	\varphi(r) < \varphi(b)
  \end{cases}$$

  Si seleccionamos un $b \in A, b \ne 0$ tal que $\varphi(1) < \varphi(b)$ entonces:\\

  $$\forall a \in A, R_b(a) = r \iff \begin{cases}
	a \equiv r \mod(b) \text{ y} \\
	r = 0 \quad o \quad \varphi(r) < \varphi(b)
  \end{cases}$$

Ahora, llamaremos: $$ A_b = \{ R_b(a): a \in A\} \subseteq A$$ que cumple:

\begin{enumerate}
	\item Si $r\in A_b \implies R_b(r) = r$
	\item $R_b(a+a') = R_b(R_b(a) + R_b(a')) \ \ \forall a,a' \in A$ \\
	\begin{proof}
	$R_b(a+a') \equiv_b a+a' \equiv_b R_b(a) +R_b(a') \equiv_b R_b(R_b(a) + R_b(a'))$.
\end{proof}
	\item $R_b(aa') = R_b(R_b(a)R_b(a')) \ \ \forall a,a' \in A$
\end{enumerate}

Además, se define la suma y el producto de $r,r' \in A_b$ de la forma:
\begin{itemize}
	\item $r+r' = R_b(r+r')$
	\item $rr' = R_b(rr')$
\end{itemize}

Con estas operaciones, $A_b$ es un anillo.
\end{nprop}

\begin{nprop}[Isomorfismo entre unidades]
  Sean $A$ y $B$ anillos conmutativos, y $f: A \to B$ un isomorfismo. Entonces $\restr{f}{U(A)} : U(A) \to U(B)$
  es un isomorfismo.
\end{nprop}

\begin{proof}
  $a \in U(A) \implies \exists a^{-1} : aa^{-1} = 1 \implies f(a)f(a^{-1}) = 1 \implies f(a) \in U(B) \text{ y } f(a)^{-1} = f(a^{-1})$.
  El resto de propiedades de isomorfismo se deducen de que $f$ es un isomorfismo.
\end{proof}

\begin{nprop}[Isomorfismo entre divisores de cero]
  Sean $A$ y $B$ anillos conmutativos, y $f: A \to B$ un isomorfismo. Entonces $\restr{f}{\rdiv_0(A)} : \rdiv_0(A) \to \rdiv_0(B)$
  es un isomorfismo.
\end{nprop}

\begin{proof}
  Análoga a la anterior: $a \in \rdiv_0(A) \implies \exists a' \in A, a' \ne 0 : aa' = 0 \implies f(a)f(a') = 0 \text{ y } f(a') \ne 0$.
\end{proof}


\begin{nprop}
  Sean $A$ un D.E. con función euclídea $\varphi$, en el que hay unicidad en cocientes y restos; y $m\in A: m\ne 0$ con $ \varphi(1) < \varphi(m)$.
	
  Consideramos $A_m = \{R_m(a): a \in A\}$ donde, como ya sabemos,

  $$\begin{array}{l}
	r+r' = R_m(r+r')\\
	rr' = R_m(rr')
  \end{array}$$

  Veamos que $R_m$ es un homomorfismo.\\
  \begin{proof}
    Notando por $\oplus$ la suma en $A_m$,

    \[
    \begin{array}{l}
      R_m(a+b) = R_m(R_m(a)+R_m(b)) = R_m(a) \oplus R_m(b) \\
      R_m(ab) = R_m(a)R_m(b) = R_m(R_m(a)R_m(b))
    \end{array}
    \]

    Por último: $R_m(1) = 1$ por $\varphi(1) < \varphi(m)$.
\end{proof}
\end{nprop}

Además, $\img(R_m) = A_m$ y $\ker(R_m) = (m) = mA = \{mx : x\in A\}$.

\begin{ncor}
  Sean $A$ un DE y $m \in A$. Entonces $A/(m)$ y $A_m$ son isomorfos con

  \[
  \begin{array}{lrll}
    F : & A/(m) & \cong & A_m \\
        & [a] & \mapsto & R_m(a)
  \end{array}
  \]
\end{ncor}



\begin{nprop} Sea $a \in A$
	\begin{nlist}
	\item $[a] \in U(A/(m)) \iff (a,m) = 1 $
	\item $a \in A_m$, entonces $a \in U(A_m) \iff (a,m) = 1$
	\item Todo elemento de $A/(m)$ es una unidad o divisor de cero
	\item Todo elemento de $A_m$ es una unidad o divisor de cero
\end{nlist}
\end{nprop}
\begin{proof}
	Vamos a probar i) y iii). Luego ii) y iv) son consecuencia del isomorfismo entre $A/(m)$ y $A_m$.\\

        $\boxed{i)}$ Sea $[a] \in U(A/(m)) \iff \exists x \in A: [a][x] = [1] \iff \exists x \in A : ax \equiv 1\mod(m) \iff (a,m) = 1$.
	
	$\boxed{iii)}$ Sea $a \in A$. Entonces la aplicación dada por $x \mapsto ax$ puede ser

        \begin{itemize}
        \item Inyectiva: entonces es sobreyectiva por ser $A/(m)$ finito. Por tanto, existe $x \in A$ tal que $ax = 1$, luego $a$ es una unidad.
        \item No inyectiva: entonces existen $u,v \in A, \ u \ne v$ tales que $au = av$, luego $a(u-v) = 0$ y $u-v \ne 0$.
        \end{itemize}

        %% TODO: poner que unidades y div. de cero son disjuntos en una proposición.
        Ahora, si esta aplicación no es sobreyectiva, entonces no puede ser inyectiva (de nuevo porque $A/(m)$ es finito), por tanto $a$ es divisor de cero.
        Si es sobreyectiva, $a$ es una unidad. Además, sabemos que los divisores de cero y las unidades de un anillo son disjuntos.
\end{proof}

\begin{ncor}
	En las mismas condiciones, son equivalentes:

	\begin{nlist}
	\item m es irreducible
	\item A/(m) ó $(A_m)$ es DI
	\item A/(m) ó $(A_m)$ es un cuerpo
\end{nlist}
\end{ncor}
\begin{proof}
	$\boxed{iii) \implies ii)}$ Todo cuerpo es un dominio de integridad.\\
	$\boxed{ii) \implies iii)}$ Supongamos $[a] \in A/(m) $ si $[a] \ne [0]$ entonces, por la proposición anterior, $a$ es una unidad.

	$\boxed{i) \implies iii)}$ Sea $[a] \in A/(m), [a] \ne [0] = 0 \implies m$ no divide a $a \implies (a,m) = 1$, como m es irreducible, sus únicos divisores son m y 1 salvo asociados $\implies [a] \in U(A/(m))$

	$\boxed{ii)\implies i)}$ Supongamos que m no es irreducible $\implies m = ab$ con a y b divisores propios $\implies m$ no divide a $a$ y m no divide a $b$ $\implies [a] \ne 0 $ y $[b] \ne 0$ en $A/(m)$ pero $[a][b] = [ab] = [m] = [0] = 0$ y como [a] y [b] son distintos de cero $\implies A/(m)$ no es DI, en contradicción con la hipótesis.
\end{proof}

En $\ent_n$ si $p \geq 2$ es un irreducible y $\ent_p $ es un cuerpo. En general, si $K$ es un cuerpo y $f(x) = \sum a_i x^i \in K[x]$ es de grado $n \implies K[x]_{f(x)}$
 es un cuerpo $\iff f(x)$ es irreducible, con $K[x]_{f(x)} = \{b_0 +b_1 x + ... + b_{n-1}x^{n-1} :  b_i \in K \}$


En particular, si $p$ es un irreducible de $\ent$ y $f(x)$ es un irreducible de $\ent_p[x]$ de grado $n \implies \ent_p[x]_{f(x)}$ es un cuerpo con $p^n$ elementos. Lo notamos $F_pn = \ent_p[x]_{f(x)}$

Salvo isomorfismos es el único cuerpo con $p^n$ elementos , al variar p y n obtenemos todos los cuerpos finitos que existen.
\subsection{Ecuaciones en $\ent_n$}

Vamos a intentar ahora encontrar una solución para una ecuación $ax = b$ en $\ent_n$ con $a \ne 0$.

\begin{enumerate}
	\item Tiene solución $\iff d =(a,n)/b$

	\item Si tiene solución, tiene exactamente $d$ soluciones distintas.
\end{enumerate}

\begin{proof}
Utilizaremos el siguiente isomorfismo para simplificar esta prueba: $\ent_n \cong \ent/(n)$.\\
\begin{enumerate}
	\item $\exists x \in Z_n :  ax = b \iff \exists[x] \in \ent/(n) : [a][x] = [b] \iff \exists x \in \ent /(n) : ax \equiv_n b \iff d/b$. Quedando probado 1.
	\item  Para demostrar 2., suponemos que d/b.\\
	 Sean $a' = \frac{a}{b}, b' = \frac{b}{d}, n' = \frac{n}{d}$. Recuperamos la propiedad anterior, $a'x \equiv b'\mod(n')$. De este expresión obtenemos la solución óptima: $x_0 : a'x_0 \equiv b'\mod(n'), 0 \le x_0 \textless n'$. Siendo la solución general, $x \equiv x_0 \mod(n')$.\\
Ahora, si $x = x_0 + kn' \quad k \in \ent$, los $x$ que satisfacen nuestro problema original son los restos de estos elementos:\quad $\{x_0,\ x_0 + n',\ x_0 +2n',...,\ x_0 + (d-1)n'\}$, si $k \in \ent, 0 \le k \textless d \implies x_0 + kn' \textless \frac{n}{d} + (d-1)\frac{n}{d} = \frac{dn}{d} = n$. Por tanto, estas son las únicas soluciones.

Podemos expresar las soluciones como $\{ [x] \in \ent /(n) : x = x_0 + kn', k \in \ent\}$. Si $k \in \ent \ y \ k = qd + r \ con \ 0 \le r \textless d, \ x_0 + kn' = x_0 + (qd + r)n' = x_0 + rn' + qdn'= x_0 + rn' +qn \implies [x_0 + kn'] = [x_0 + rn']$. Las soluciones de la ecuación original serán $\{[x_0],\ [x_0 + n'],\ [x_0 +2n'],...,\ [x_0 + (d-1)n']\} \subseteq \ent /(n) \cong \ent _n, \forall r, 0 \le r \le d - 1.$
$\{x_0,\ x_0 + n',\ x_0 +2n',...,\ x_0 + (d-1)n'\} \subseteq \ent _n$
\end{enumerate}



\end{proof}

\section{Función de Euler.}
\[
  \varphi : \mathbb{N}-\{0\} \longrightarrow \mathbb{N}
\]

definida de la siguiente forma $\forall n\ge 1$:
\[
  \varphi(n) = |\{m\in \mathbb{N} : 1\le m \le n \text{ y } (m,n)=1\}|
\]

Que es igual al número de naturales menores que $n$ y primos con él.
%\vspace{5mm}

\begin{nprop}
  Si $(m,n) = 1 \implies \varphi(mn) = \varphi(m)\varphi(n)$
\end{nprop}

Necesitamos algunos resultados para probar esto:

\begin{ndef}[Anillo producto.]
Si $A$ y $B$ son anillos:
\[
  A\times B = \{(a,b) : a\in A, b\in B\}
\]

Donde se definen:
\begin{align*}
  (a_1,b_1)+(a_2,b_2) &= (a_1+a_2, b_1+b_2)\\
  (a_1,b_1)(a_2,b_2) &= (a_1a_2, b_1b_2)
\end{align*}

$U(A\times B) = U(A)\times U(B)$

\end{ndef}

\begin{nota} Un anillo producto nunca es un cuerpo.

\end{nota}

\begin{nth}[Versión clásica del teorema chino del resto.]
Si $(m,n)=1 \implies \mathbb{Z}_{mn} \cong \mathbb{Z}_m \times \mathbb{Z}_n \iff \mathbb{Z}_{|(mn)} = \mathbb{Z}_{|(m)} \times \mathbb{Z}_{|(n)}$
\end{nth}

\begin{proof}
\begin{align*}
f:& \mathbb{Z} \longrightarrow \mathbb{Z}_{|(m)}\times \mathbb{Z}_{|(n)}\text{ es un homomorfismo de anillos.}\\
& a \longmapsto ([a]_m, [a]_n)
\end{align*}

Probaremos que es sobreyectivo:\\\\
$([b]_m,[c]_n)$ ¿$\forall b,c \in \mathbb{Z}\quad \exists a \in \mathbb{Z} : [a]_m = [b]_m \text{ y } [a]_n = [c]_n$?

¿$\forall b,c \in \mathbb{Z}\quad \exists a \in \mathbb{Z} : \begin{cases}
  a \equiv b \mod(m)\\
  a \equiv c \mod(n)
\end{cases}$?$ \iff b \equiv c \mod(m,n) \implies b \equiv_1 c$

Sí es sobreyectiva.

$\ker(f) = (mn)$, pues $m$ y $n$ son primos entre sí. Ahora, como $f$ es sobreyectiva, por el teorema de isomorfía tenemos el resultado.
\end{proof}

\begin{ncor}
  Si $(m,n) = 1 \implies \varphi(mn) = \varphi(m)\varphi(n)$
\end{ncor}

\begin{proof}
  $\varphi(mn) = |U(\mathbb{Z}_{mn})| = |U(\mathbb{Z}_{|m}\times\mathbb{Z}_{|n})| = |U(\mathbb{Z}_{|m}) \times U(\mathbb{Z}_{|n})| = \varphi(m)\varphi(n)$
\end{proof}
\begin{nota}
  Como sabemos (aunque no lo hayamos probado), si $n\in \mathbb{N}, n= p_1^{e_1}\dots p_n^{e_n} \text{ con } p_i \ne p_j,\; p_i\ \text{primo de $\mathbb{Z}$ (irreducible)}$. Así, $\varphi(n) = \varphi(p_1^{e_1})\dots\varphi(p_n^{e_n})$.

  $\varphi(p^e) = p^e(1-\frac{1}{p}) = p^e-p^{e-1}$
\end{nota}



\begin{nth}
	Si $(a,m)=1 \implies a^{\varphi(m)} = 1$ en $ \mathbb Z_{m}$ \hspace{0.5cm} $a \in \mathbb Z_{m} $

Por tanto, $ a^{\varphi(m)-1} = a^{-1} $ en $ \mathbb Z_{m} $

\end{nth}


\begin{nth}[Teorema de Euler]
	$ \forall a \in \mathbb Z $ si $(a,m) = 1 \implies a^{\varphi(m)} \equiv 1 \mod(m)$.
\end{nth}


\begin{nth}[Teorema pequeño de Fermat]
	Si p es irreducible, $ \forall a \in \mathbb Z, $ $ a^{p} \equiv a \cdot \mod(p) $
\end{nth}



\begin{ncor}
	Si p es irreducible, $ \forall a \in \mathbb Z_{p} $	 \hspace{0.25cm} $ a^{p} = a $ en $ \mathbb Z_{p} $
\end{ncor}


\begin{proof}
	Partiendo de la hipótesis del teorema demostraremos el corolario:

$ m = p, $ $ a \neq 0 $ en  $ \mathbb Z_{p}$ (si $a=0$ es trivial) $ \implies (a,p) = 1 \implies a^{\varphi(p)} = 1 $ en $ \mathbb Z_{p} $

$ \varphi(p) = p(1-\frac{1}{p}) = p-1 \implies a^{p-1} = 1 $ en $ \mathbb Z_{p} \implies a^p=a $ en $ \mathbb Z_{p} $

\end{proof}



\section{Dominio de Factorización Única (DFU)}

Un Dominio de Integridad, A, es llamado un DFU si $\forall a \in A, a\ne 0$ y $a\notin U(A)$, entonces $\exists$ irreducibles $q_1 ,...,q_r \in A : a = q_1 ...q_r$ tales que la factorización es esencialmente única en el sentido de que si $q_1' ,...,q_s' \in A $ con $q_j$ irreducible $\implies r = s$ y $\exists \sigma:\{1,...,r\} \cong \{1,...,s\}$ una permutación tal que $q_i'$ es asociado con $q_{\sigma(i)}$

\begin{ndef}[Conjunto representativo de los irreducibles de A]
	Si A es un DFU, vamos a denotar $\mathcal{P}= $ un conjunto representativo de los irreducibles de A.
\end{ndef}

\begin{itemize}
\item $\forall p \in \mathcal{P}$, p es un irreducible
\item $\forall p, q \in \mathcal{P}$, $p$ y $q$ no son asociados entre sí
\item $\forall p $ irreducible de A, $\exists q \in \mathcal{P} : p \sim q$
\end{itemize}

Supongamos ahora que estamos en un DFU y hemos seleccionado un conjunto $\mathcal{P}$.

Si tenemos un $a\in A,\ a \notin U(A), a \ne 0$, por definición existirán $q_1 ,...,q_r \in A $ irreducibles tales que $a = q_1...q_r $. Entonces, $\forall i = 1,...,r\ \exists p_1,...,p_r \in \mathcal{P} : q_i = u_ip_i $ con $u_i \in U(A)$. Así, $a$ se puede expresar como: $a = (u_1...u_r)p_1...p_r$ pero todos los $u_i$ son unidades del anillo, luego $\exists p_1,...,p_r \in \mathcal{P} $ y $u \in U(A) : a = u(p_1...p_r)$. Esta descomposición es esencialmente única pero de forma más fuerte que antes. Además, es única salvo orden de escritura de los $p_i$.

\begin{ejemplo}
	En $\ent$ el $-6$ se puede escribir como $(-1)*2*3$ ó como $(-1)*3*2$
\end{ejemplo}

Estos $p_i$ pueden repetirse, así que si agrupamos en términos obtenemos:

$\forall a \in A,\ a \ne 0\ \exists p_1,...,p_s \in \mathcal{P}$ con $p_i \ne p_j$, $e_1,...,e_s \in \ent$ y $ u \in U(A) : a = u(p_1^{e_1}...p_s^{e^s})$


\begin{ndef}[]
	Si $p\in \mathcal{P}$ y $a \in A,\ a \ne 0$ denotamos $e(p,a)$ como:
	\begin{nlist}
	\item exponente con que $p$ aparece en la factorización de $a$, si aparece. $e(p_i,a) = e_i$ \quad $i = 1,...,s$
	\item 0 en otro caso. $e(p,a) = 0 \quad \forall p \notin \{p_1,...,p_s\}$

\end{nlist}

	Vamos a asumir a partir de ahora que $a^0  = 1 $ en cualquier anillo. Así, podemos ver que:

	$$\forall a \in A, \quad a = u(\prod _{p \in \mathcal{P}}p^{e(p,a)})$$
\end{ndef}

\textbf{Propiedades:}
\begin{nlist}
	\item $e(p,ab) = e(p,a) + e(p,b)$.\\

	\begin{proof}

	Con el $a$ anterior y $b= v(\prod _{p \in \mathcal{P}}p^{e(p,b)})$. Entonces $ab = uv(\prod _{p \in \mathcal{P}}p^{e(p,a)+ e(p,b)})$
\end{proof}

\item $a,c \ne 0$ y $a/c \iff \forall p \in \mathcal{P}, \quad e(p,a) \leq e(p,c)$\\

\begin{proof}
	$\boxed{\Rightarrow} \quad \exists b: c =ab \implies \forall p \in \mathcal{P}, \quad e(p,c) = e(p,ab) = e(p,a)+e(p,b) \geq e(p,a) $

	$\boxed{\Leftarrow} $ ¿Existe un $b$ tal que $ab= c$?

	Si $c$ es: $c = v(\prod _{p \in \mathcal{P}}p^{e(p,c)})$

	Si tomamos $b = (u^{-1}v)(\prod _{p \in \mathcal{P}}p^{e(p,c) - e(p,a)}) $ y multiplicamos por $a$, obtenemos $c$.
\end{proof}
\end{nlist}

\begin{nprop}
En un DFU existen mcd y mcm de cualesquiera elementos. Así:

\[
\forall a,b \ne 0 \quad (a,b) = (\prod _{p \in \mathcal{P}}p^{min\{e(p,a),e(p,b)\}})
\]
\[
[a,b] = (\prod _{p \in \mathcal{P}}p^{max\{e(p,a),e(p,b)\}})
\]

\end{nprop}
\begin{proof}
	Probaremos el caso del mcd, el caso de mcd se hace de la misma forma.\\
	Sea d un elemento del anillo tal que $e(p,d) = min\{e(p,a), e(p,b)\}\ \forall p \in \mathcal{P}$, es evidente que d/a,b. Ahora, si tenemos un divisor común cualquiera , digamos $c \implies c/a$  y $c/b \implies e(p,c) \leq e(p,a),e(p,b) \implies e(p,c) \leq e(p,d) \implies c/d$ luego $d$ es un máximo común divisor.
\end{proof}

\begin{ndef}[Elemento Primo]
	Si $A$ es un D.I. un elemento $p\in A , p \notin U(A), p \ne 0$ es llamado “primo" si se verifica la siguiente propiedad:

	Si $p$ no divide a un elemento $a$ ni a un elemento $b \implies p$ no divide a su producto.\\
	Equivalentemente: si $p/ab \implies p/a$ o $p/b$

\end{ndef}

\begin{nprop}
\begin{nlist}
	\item Todo primo es irreducible en cualquier anillo A
	\item Si A es un DFU, entonces todo irreducible es primo.
\end{nlist}
\end{nprop}
\begin{proof}
	\begin{nlist}
	\item Sea $p$ un elemento primo. Supongamos que $p=ab$, producto de dos elementos, bastaría ver que uno de ellos es un asociado solo. Ahora, como $p/p \implies p/ab \implies  p/a $ o $p/b \implies a \sim p$ o $b \sim p$

	\item $p \in \mathcal{P}$, veamos que $p$ es primo.

Supongamos que $p/ab$. Veamos que $p$ divide a $a$ o a $b$.
Si $p/ab\implies e(p,ab) \geq 1$ pero sabemos que $e(p,ab) = e(p,a)+e(p,b) \implies e(p,a) \geq 1$ o $e(p,b) \geq 1$. Si ocurre lo primero, $p/a$ y si ocurre lo segundo $p/b$ luego si $p$ divide a un producto, entonces $p$ divide a uno de los dos elementos del producto.
\end{nlist}
\end{proof}



\begin{nth}
	Sea A un D.I. Entonces, son equivalentes:
	\begin{nlist}
	\item A es un DFU
	\item \begin{enumerate}
	\item Todo elemento no nulo ni unidad de A factoriza como producto de irreducibles
	\item Todo irreducible de A es primo
\end{enumerate}
	\item \begin{enumerate}
	\item Idem
	\item $\forall a,b \in A \ , \exists \ mcd(a,b)$
\end{enumerate}
\end{nlist}
\end{nth}
\begin{proof}
	Que $(i)\implies (ii)$ es trivial. Veamos que $(ii) \implies (i)$

	Lo único que falta para probar que es un DFU es probar que las factorizaciones son únicas. Sea $a= p_1...p_r = q_1...q_s$ con $p_i,q_j$ irreducibles. Vamos a ver que $r=s$. Para ello, vamos a hacer una inducción en $r$.

	\begin{itemize}
	\item Caso $r=1 \implies p_1 = q_1...q_s$. Ahora, ¿puede ser $s>1$? Los $q$ no son unidades, pues son irreducibles, por tanto, si $s$ fuese mayor que 1 serían los divisores propios de $p_1$, pero eso no puede ocurrir porque $p_1$ es irreducible. Como no se puede dar que $s>1$ $\implies$ $s=1=r \implies p_1 = q_1$

	\item Si $r> 1$ y usando la hipótesis de inducción, entonces $s > 1$.

	Nos fijamos en $p_1$, que es claro que divide a $a\implies p/(q_1...q_s)\implies^{p_1 \ primo} \exists j : p_1 /q_j$ y reordenando podemos suponer que $p_1/q_1$.

	Esto implica que $p_1 \sim q_1 \implies \exists u \in U(A): q_1 = up_1$. Ahora nos podemos llevar la expresión a la igualdad de a($a= p_1...p_r = q_1...q_s$ ) $\implies p_1...p_r = up_1q_2...q_s $ y podemos reducir dividiendo por $p_1$ y nos queda $p_2...p_r = uq_2...q_s$.

	Ahora, usando la hipótesis de inducción, nos queda en cada lado $r-1$ elementos y $s-1$ elementos y por tanto $r-1 =s-1\implies r=s$
\end{itemize}

Ahora, que $(i) \implies (iii)$ es trivial. Como (i) y (ii) son equivalentes, basta probar que $(iii) \implies (ii)$

Queremos probar que todo irreducible es primo. Sea $p$ un irreducible. Supongamos que $p$ no divide ni a $a$ ni a $b$. Probaremos que entonces, no divide al producto $ab$

Es fácil ver que $(p,a) = 1$ y que $(p,b) = 1$. Ahora, por la propiedad del mcd que asegura que:
\[
(a,b) = 1 \ y \ (a,c) = 1 \iff (a,bc) = 1
\]

Entonces, $(p,ab) = 1 \implies$ $p$ es primo relativo con el producto, por tanto, $p$ no divide al producto y así $p$ es primo.
\end{proof}

	\textbf{Lema previo:} En un DIP, toda cadena ascendente de ideales es estacionaria. En otras palabras, si A es un DIP, $I_1 \subseteq I_2 \subseteq ... \subseteq I_{n}$ es una sucesión de ideales creciente respecto a la inclusión (cada uno está incluido en el siguiente). $\implies \exists m : I_m = I_{m+1} = ... = I_{m+k} \quad k \geq 1$.

	\begin{proof}[Demostración del lema.]
	Podemos ver que:
	\[
	I = U_{n\geq 1} I_n = \{a \in A : \exists n \ con \ a \in I_n\} \quad \forall a,b \in I \implies \exists n: a,b \in I_n\]\[ \implies a+b \in I_n \implies a+b \in I
	\]
	Realizando la prueba análoga para el producto, I es un ideal y por estar en un DIP, es principal $\implies \exists a \in A : I = (a) = \{ax : x \in A\}$. Que es no vacío, pues $a \in I$.

	Si $a\in I$, en particular estará en alguno de los $I_i$ de la unión $\implies \exists m : a \in I_m \implies (a) \subseteq I_m$, pero  $I = (a) \subseteq I_m \subseteq I_{m+k} \subseteq I \implies I_i = I_j$ para todo $i$ y $j$.

\end{proof}

\begin{nth}
	Todo DIP es un DFU (Lo cual implica que todo DE es un DFU)
\end{nth}
\begin{proof}
	Tenemos que probar que en un DIP todo elemento se puede descomponer como producto de irreducibles. Para ello, vamos a negar la tesis. Supongamos que estamos en un DIP y que en ese anillo existen elementos distintos de cero que no son unidades y no se pueden descomponer como producto de irreducibles.

	Supongamos que $a$ es un elemento de esa “clase”. Entonces $\exists  \ a'$ divisor propio de $a$ que también es de esa “clase de elementos”.

	Podemos asegurar que $a$ no es un irreducible, pues no admite factorización en irreducibles y si fuera irreducible, él mismo sería una factorización como irreducibles. $\implies \exists b,c : a = bc$ con b y c divisores propios. Entonces, uno de los dos (b ó c) no puede admitir una factorización como producto de irreducibles, pues si no, $a$ admitiría esa factorización. Entonces, llamamos $a'$ a $b$ o a $c$ según sea el que no admita esa factorización.

	Ahora, vamos a construir una sucesión $\{a_n\} \in A$ con $a_1 = a,\ a_{n+1}=a'_n$. Cada elemento siguiente, es un divisor propio del anterior y es de la “clase” que establecimos al principio (no es cero, ni una unidad, ni se puede factorizar en producto de irreducibles). Esto implica que $a_{n+1} = a_n'$ y $a_{n+1}$ no es sociado con $a_n$.

	Si consideramos los ideales principales generados por los elementos de esta sucesión, podemos ver que (como $a_{n+1}$ es divisor de $a_n$):

	\[
	\implies (a_1) \subset (a_2) \subset ... \subset (a_n) \subset (a_{n+1}) \subset ...	\]

	Y en esta cadena no hay igualdades, pues si $(a_n) = (a_{n+1}) \implies a_{n+1} \in (a_n) \implies a_n/ a_{n+1}$ y esto no puede ocurrir.

	Pero esto contradice el lema que hemos visto anteriormente, por tanto hemos probado así que todos los elementos deben tener una factorización y por tanto estamos en un DFU.
\end{proof}


\begin{nprop}
	Si $\alpha \in \ent[\sqrt n]$ es un divisor propio de $\beta$ en $\ent[\sqrt n]$ entonces $N(\alpha)$ es un divisor propio de $N(\beta)$ en $\ent$
\end{nprop}
\begin{proof}
	Como $\alpha$ es un divisor de $\beta$ entonces $\exists \gamma \in \ent[\sqrt n] : \beta = \alpha \gamma \implies N(\beta) = N(\alpha)N(\gamma) \implies N(\alpha)/N(\beta)$.

	Ahora, la norma de $\alpha$ no puede ser ni 1, ni -1 pues si no sería una unidad y, por tanto, no sería divisor propio; $\alpha$ no puede ser un asociado pues si no, $\gamma$ sería un divisor propio también, luego $N(\alpha)$ tiene que ser un divisor propio de $N(\beta)$

\end{proof}
\begin{ncor}
	Si $N(\alpha) = \pm p$ con $p $ un primo de $\ent, p\geq 2\implies \alpha$ es irreducible en $\ent[\sqrt n]$
\end{ncor}
\begin{ncor}
	Si $\alpha$ es primo en $\ent[\sqrt n]\implies N(\alpha) = \pm p $ ó $\pm p^2$ con $p \geq 2 $ un primo de $\ent$. Además, si $N(\alpha) = \pm p^2 \implies \alpha \ y \ p$ son asociados en $\ent[\sqrt n]$
\end{ncor}
\begin{proof}
	Supongamos que $\alpha \in \ent[\sqrt n]$, primo. Consideramos su norma: $N(\alpha)$ que no es ni 1 ni -1 pues si no sería una unidad. Así: $N(\alpha) = p_1...p_r$ con $p_i\in \ent$ primos, lo que implica que $\alpha \bar{\alpha} = p_1...p_r \implies \alpha/p_1...p_r $ pero $\alpha$ es primo, luego $\exists i \in \{1,...,r\}: \alpha / p_i \implies \exists p \geq 2 $ primo de $\ent$ tal que $\alpha / p \ en \ \ent[\sqrt n]$.

	 Esto implica $p = \alpha \beta \ con \ \beta \in \ent[\sqrt n] \implies p^2 = N(\alpha) N(\beta) \implies N(\alpha)/p^2 \implies N(\alpha) = \pm p$ ó $\pm p^2$, como queríamos.\\
	 Si $N(\alpha) = p^2 \implies N(\beta) = 1 \implies \beta$ es una unidad $\implies \alpha$ y $p$ son asociados
\end{proof}
\begin{nota}
	Si estuviéramos en un DFU, ser irreducible y ser primo son equivalentes, luego estos enunciados valdrían igual para elementos primos.
\end{nota}

\begin{ejemplo}
	Factorizar $2i$ y $11+7i$ en producto de irreducibles(primos por estar en un DFU).

	\begin{enumerate}
	\item Primero, calcularemos su norma. $N(11+7i) = 11^2 + 7^2 = 170$
	\item Factorizamos la norma en $\ent$. $170 = 2*85 = 2*5 *17$.
	\item Ahora, los factores irreducibles serán los enteros de Gauss cuya norma sea un primo o el cuadrado de un primo. Por tanto, un divisor de este número será un entero de Gauss $Z[i]$ cuya norma sea un divisor de la norma de $11+7i$, por tanto su norma será 2, 5 ó 17.
	\item  $N(a+bi) = a^2 +b^2 = 2 \iff a=\pm 1$ y $b= \pm 1.$ Los enteros de Gauss de norma 2 son: $1+i,1-i,-1+i,-1-i$, es decir $1+i$ y sus 3 asociados.
	\item Ahora, tenemos que plantearnos si $1+i / 11+7i$, vemos que la división es: $11+7i/1+i = 9-2i\in \ent[i]$. Además, como $1+i$ tiene norma 2, que es un primo de $\ent$ luego ya tenemos un irreducible por el corolario 8.
	\item Tenemos que repetir el proceso para $9-2i$.
	\item Su norma es $N(9-2i) = 5*17$ pues es el de antes quitándole el irreducible cuya norma vale 2.
	\item Buscamos los enteros de Gauss cuya norma valga 5. $N(a+bi) = a^2 +b^2 = 5 \iff a = \pm 1$ y $b =\pm 2$ ó $a \pm 2$ y $b = \pm 1$.

	Estos son: $1+2i$ y sus asociados para el primer caso y $2+i$ y sus asociados para el segundo caso.

	\item Ahora, tenemos que ver si estos dividen a $9-2i$.
	\begin{itemize}
	\item $9-2i/2+i  = \frac{16}{5} + \frac{13}{5}i \notin \ent[i]$
	\item $9-2i /1+2i = 1-4i \in \ent[i]$
\end{itemize}

Por lo que tenemos que $11+7i = (1+i)(1+2i)(1-4i)$ y ahora tenemos justo 3 irreducibles con las normas que buscábamos, luego tenemos hecha la factorización en irreducibles.

\end{enumerate}

Ahora, haciendo lo mismo para $2i$ vemos que $2i = (1+i)^2$.


\end{ejemplo}
\begin{ejemplo}[2]
	Vamos a factorizar 180 en $\ent[i\sqrt2]$. Para ello, vemos que $180 = 2^2 * 3^2 * 5$. Recordamos que en este anillo, $N(a+b\sqrt2) = a^2 +2b^2$ y $U(\ent[i\sqrt2]) = \pm 1$

	Ahora, como $N(2) = 4$, un divisor propio del 2 tendrá por norma un divisor propio del 4 en $\ent$. En $\ent$, sólo el 2 es divisor propio del 4. Por ello, tenemos que plantearnos la ecuación $a^2+2b^2 = 2$. Entonces, los únicos elementos que hay que tienen son $\sqrt{-2}$ y $-\sqrt{-2}$.
	Ahora vemos si alguno de estos divide a 2:
	\[
	\frac{2}{\sqrt{-2}} =  \frac{2*(-\sqrt{-2})}{\sqrt{-2}*(-\sqrt{-2})}= \frac{2(-\sqrt{-2})}{2} = -\sqrt{-2} \ \in \ent[i\sqrt2]
	\]
	Ahora, como $N(\sqrt{-2})=2$ que es un primo en $\ent \implies \sqrt{-2}$ es un primo de $\ent[\sqrt{-2}] = \ent[i\sqrt2]$ y por ello $2=-(\sqrt{-2})^2$.

	Seguimos, haciendo lo mismo con el 3. $N(3) = 9$. ¿Existen a y b : $a^2 + 2b^2 = 3$? Vemos que tomando $a=\pm 1$ y $b=\pm1$ se puede llegar a la igualdad. Es decir, tenemos los elementos: $\{1+\sqrt{-2}, \ 1-\sqrt{-2}, \ -1-\sqrt{-2}, \ - 1 + \sqrt{-2}\}$.

	Probamos dividiendo $3/(1+\sqrt{-2}) = 1 - \sqrt{-2}\implies 3 = (1+\sqrt{-2})(1-\sqrt{-2})$ y ambos son irreducibles.

	Hacemos lo mismo con el 5. Tenemos que discutir la ecuación: $a^2+2b^2 = 5$. Sin embargo, en este caso no hay ningún elemento que tenga solución luego 5 es primo en $\ent[i\sqrt2]$.

	Por tanto, la factorización de 180 en $\ent[i\sqrt2]$ es: $180 = (\sqrt{-2})^4*(1+\sqrt{-2})^2*(1-\sqrt{-2})^2*5$
\end{ejemplo}

\begin{ejemplo}[Ejemplo de anillo que no es un DFU] $\quad$

Como ejemplo, vamos a probar que \textbf{\underline{$\ent[\sqrt{-5}]$ no es un DFU.}} En este anillo, $N(a+bi\sqrt5) = a^2 + 5b^2$ y $U(\ent[\sqrt{-5}])= \{\pm 1\}$

Vamos a considerar el elemento $1+i\sqrt5$. Su norma es: $N(1+i\sqrt5) = (1+i\sqrt5)(1-i\sqrt5) = 6 = 2 * 3$. ¿Es este elemento irreducible?

Vamos a plantearnos qué elementos del anillo $\ent$ tienen norma 2 o norma 3.
\begin{itemize}
	\item En la ecuación $a^2+5b^2= 2$ no hay soluciones en $\ent[i\sqrt5]$
	\item En la ecuación $a^2+5b^2 = 3$ tampoco hay soluciones en este anillo.
\end{itemize}
	Como no tiene divisores propios, entoces este elemento es irreducible. Su conjugado, por el mismo motivo, también es un irreducible.

	Ahora, por la norma de $1+i\sqrt5$ hemos obtenido una factorización del 6 en producto de irreducibles. Pero el 6 también es $2*3$ en este anillo. Esta podría ser otra factorización en irreducibles de 6 en $\ent [i\sqrt5]$. El 2 no tiene ningún divisor propio en este anillo, pues es el único divisor de 4 en $\ent$ luego el 2 es irreducible en este anillo. Lo mismo ocurre con el 3.

	Por tanto, tenemos dos descomposiciones del 6 en producto de irreducibles, que no son iguales ni asociados luego este anillo no puede ser un DFU.

	Ahora, el 2 es un irreducible, \underline{¿es 2 primo?} Vemos que $2/6$, y $6 = (1+i\sqrt5)(1-i\sqrt5)$. Si fuese primo, necesitaríamos $2/1+i\sqrt5$ ó $2/1-i\sqrt5$ y eso no ocurre pues sus normas no se dividen en $\ent$ por tanto el 2 no es primo.
\end{ejemplo}

\subsection{$\ent[x]$ es un DFU y no es un DIP}

Vamos a estudiar ahora que este anillo es un DFU sin ser un DIP ni un DE.

Bastaría de hecho tomar los elementos 2 y x para ver que $(2,x) = 1$ y no existen los coeficientes de Bezout para estos elementos, es decir:
\[
\nexists f,g \in \ent[x] : 1 = 2f(x)+ xg(x)
\]

Vamos ahora a enunciar y a demostrar el Teorema de Gauss sobre los DFU. Sin embargo, antes debemos aclarar algunos conceptos.
\begin{ndef}
Si $f \in A[x],\ gr(f) \geq 1$, se define su contenido como el m.c.d. de sus coeficientes. Lo denotamos por $c(f)$. (Si $f= \sum a_i x^i \implies c(f) = mcd(a_0,...,a_n)$	)\\
Se dice que $f$ es primitivo si $c(f) = 1$. El contenido es único salvo asociados.
\end{ndef}
\textbf{Lema.} $ \ c(af) = ac(f)$. Esta propiedad es consecuencia directa de que $(ab,ac) = a(b,c)$
\textbf{Lema.} Todo $f: gr(f) \geq 1$ se puede factorizar de la forma $f = af'$ con $f'$ primitivo. Además, esta factorización es esencialmente única.
\begin{proof}
	Sea $f$ un polinomio y $a = c(f)$. Entonces, $a/a_i \ \forall i$. Tomamos $a_i' = \frac{a_i}{a} \in A$. Sea también $f' = \sum a_i'x^i$. Ahora, si tomáramos $af'= \sum aa_i' x^i = f \implies a = c(f) = c(af') = ac(f') \implies c(f') = 1$ simplificando por a, luego hemos encontrado una factorización $f=af'$.

	Además, podemos ver que la factorización es única, pues si $f=bg$ con $c(g)=1$ y $f=af'$, entonces $a = c(f) = b c(g) = b \implies a=b$ y $g=f'$
\end{proof}

\textbf{Lema.} Todo $\phi,\ gr(\phi) \geq 1$ se puede factorizar de forma esencialmente única como $\phi = \frac{a}{b} f$ con $f$ primitivo.
\begin{proof}
	Sea $\phi = \sum \frac{a_i}{b_i} x^i$. Tomamos $b= \prod b_i \implies b_i/b \ \forall i \implies b_i /ba_i \ \forall i \implies b \phi = \sum \frac{ba_i}{b_i}x^i$ Pero el numerador es un múltiplo del denominador, luego $g = \sum \frac{ba_i}{b_i}x^i \in A[x]$.

	De esta forma, $b\phi = af \implies \phi = \frac{a}{b} f$.

	Además, veamos que es única. Sea ahora $\phi = \frac{c}{d} f'$ con $f'$ primitivo$\implies \frac{a}{b} f = \frac{c}{d} f' \implies daf=bcf' \implies da = bc$ y $f=f'$
\end{proof}

Enunciaremos también un lema muy importante, el lema de Gauss.\\

\textbf{Lema de Gauss.} Sea A un DFU. Si tenemos $f,g \in A[x]\implies c(fg) = c(f)c(g)$. En particular, el producto de polinomios primitivos es primitivo.\\
\begin{proof}
	Probaremos en primer lugar que el producto de polinomios primitivos es primitivo. Sean $f= \sum a_i x^i $ y $g= \sum b_j x^j$ ambos en $A[x]$ y primitivos. Consideramos también su producto $fg = \sum c_k x^k$. Sabemos que  $c_k = \sum_{i+j=k}a_ib_j$. Negaremos la tesis e intentaremos llegar a una contradicción.

	Supongamos que $fg$ no es primitivo. Entonces, $c(fg) \ne 1$. Como A es un DFU, el contenido se podrá expresar como un producto de primos (irreducibles). Existirá al menos un primo que lo divida. $\exists p $ primo de $A$ con $p/c(fg) \implies p/c_k \ \forall k$. Sin embargo, ese p no puede dividir a $a_i \ \forall i$ pues f es primitivo. Por la misma razón, no puede dividir a todos los $b_j$.

	Por tanto, sea $r$ el primer índice tal que $p$ no divide a $a_r$. Tomemos también $s$ el primer índice tal que $s$ no divide a $b_s$. Ahora, vamos a fijarnos en el coeficiente $c_{r+s}$ de $fg$. Este coeficiente se expresa como:
	 \[c_{r+s} = \sum _{i+j=r+s}a_ib_j = \sum _{i+j =r+s \ \ i < r}a_ib_j + a_rb_s + \sum_{i+j=r \ \ i > r}a_ib_j\]
	 De esta expresión, podemos ver que $p/c_{r+s}$, que $\forall i < r$ entonces $p/a_ib_j \implies p/\sum_{i+j =r+s}a_ib_j$ por dividir a los $a_i$. Ocurre lo mismo con el mismo en el término de $i > r $ pues p divide a $b_j$. Si despejáramos $a_rb_s$ veríamos que $p/a_rb_s$ pues divide a todos los sumandos, pero p es un primo, por tanto si divide a un producto tiene que dividir a alguno de los factores, pero como habíamos dicho que no divide a ninguno de los dos, hemos llegado a una contradicción y, por tanto, $fg$ es primitivo.

	 Ahora demostraremos c(fg) = c(f)c(g). Podemos poner $f=af'$ con $f'$ primitivo e igual con $g=bg'$. Entonces $fg= abf'g'\implies c(fg)=c(abf'g') = abc(f'g')$ pero $f'$ y $g'$ son primitivos luego su contenido es 1. Así, el contenido de f es $a$ y el de g es $b$ luego $c(fg) = c(f)c(g)$ probando así el lema.
\end{proof}
\newpage
\textbf{Lema.}
\begin{nlist}
	\item Si $a \in A$, $a$ es irreducible en A[x] $\Leftrightarrow$ $a$ es irreducible en A.
	\item Si $gr(f) \geq 1$, f es irreducible en A[x] $\Leftrightarrow$ f es primitivo y es irreducible en K[x].
 \end{nlist}

\begin{proof}
$(i)$ Una factorización de un polinomio de grado 0 será irreducible si es irreducible en A (porque U(A[x]) = U(A)).\\
$(ii)$ $\displaystyle \boxed{\Rightarrow}\ $ Si f es irreducible, el contenido tiene que ser o una unidad o un asociado. No puede ser asociado porque...
Como c(f)/f, si f es irreducible $\Rightarrow$ c(f) = 1 $\Rightarrow$ f es primitivo. Vamos a probar que es irreducible en K[x] por contradicción. Supongamos que f no es irreducible en K[x] $\Rightarrow$ f = $\phi \psi$ con $gr(\phi) \geq 1$ y $gr(\psi) \geq 1$
$$
 \left .
      \begin{matrix}
        \phi =  \frac{a}{b}g : g\ es\ primitivo\\
        \psi =  \frac{c}{d}h : h\ es\ primitivo
      \end{matrix}
   \right \}
   \Longrightarrow \quad f = \frac{a}{b} \frac{c}{d} gh \Rightarrow bdf = acgh
  $$
  Aplicando contenidos sobre esta igualdad nos queda, $c(bdf) = c(acgh) \Rightarrow bdc(f) = acc(gh) \Rightarrow bd = ac \Rightarrow f = gh $ CONTRADICCIÓN f es irreducible y ninguno es unidad ya que el grado es mayor o igual que 1.

$\displaystyle \boxed{\Leftarrow}\ $Suponemos que $f=gh$ en A[x], $g$ y $h$ no unidades. Entonces, $gr(g)$ y $gr(h)$ es mayor o igual que 1. (Suponemos $gr(g)=0$, $g = a \in A; f=ah \Rightarrow 1 = ac(h) \Rightarrow g=a \in U(A)=U(A[x])$, llegamos a una contradicción con la suposición inicial).  \\
$f=gh$ se da en K[x] y muestra que f no es irreducible en K[x], contradicción con la hipótesis.
En K[x] no hay irreducibles de grado 0, las constantes tienen inverso.
\end{proof}

\textit{Observación:} Sea $\phi \in$ K[x], $gr(\phi) \geq$ 1, $\phi=\frac{a}{b}f$ con $f$ primitivo. $\phi$ es irreducible en K[x] $\Leftrightarrow$ f es irreducible en A[x].

\begin{nth}[Teorema de Gauss]
	Si A es un DFU $\implies A[x]$ es también un DFU.
\end{nth}

\begin{proof}
Sea $f \in$ A[x], $f \neq 0$ y $f \notin$ U(A[x])=U(A)
\begin{itemize}
	\item Caso $gr(f) = 0$. $f = a \in$ A, $a \neq 0$ y $a \notin$ U(A). Como A es un DFU, existen $p_1,...,p_r$ irreducibles de A (también de A[x]) tales que $a = p_1...p_r$
	\item Caso $gr(f) \geq 1$ y $f$ primitivo. Existen $\phi_1,...,\phi_r \in$ K[x] irreducibles tales que $f=\phi_1...\phi_r$. $\phi_i = \frac{a_i}{b_i}f_i : f_i \in$ A[x], primitivo $\implies f = \frac{a}{b}f_1...f_r$ donde $a=\prod a_i,\ b=\prod b_i \implies bf = af_1...f_r$ $\implies^{(1)} b=a \implies f = f_1...f_r$. Por el lema anterior, como $f_i$ son irreducibles en K[x] y, además, son primitivos, $f_i$ son irreducibles en A[x].
	\item Caso general, $gr(f) \geq 1$ y $f$ no primitivo. Tomamos $f=af'$
con $f'$ primitivo. $c(f) = a \neq 0,\ a \notin$ U(A)\\
$^{(1)}$ Como $f_i$ son primitivos
\end{itemize}
\end{proof}


En lo sucesivo, vamos a notar $A\subseteq K = Q(A)$. De esta forma, también sucede $A[x] \subseteq K[x]$. También vamos a notar:
\begin{itemize}
	\item $a,b,c.... \in A$
	\item $f,g,h... \in A[x]$
	\item $\phi,\psi, ... \in K[x]$
\end{itemize}

\begin{ncor}
	Si $f\in A[x]$ es irreducible en $A[x]$ con $gr(f) \geq 1\implies$ $f$ es primo en $A[x]$.
\end{ncor}
\begin{proof}
	Supongamos que $f/gh$ en $A[x]\subseteq K[x]$. Como $f$ es irreducible en $A[x]$, lo  es en $K[x]$ y $K[x]$ es un DFU por ser K un cuerpo, entonces $f$ es primo en $K[x]$. Entonces, podemos asegurar que o bien $f/g$ o bien $f/h$ en $K[x]$.

	Para lo que sigue, supongamos que $f/g$ en $K[x]$ (Si quisiéramos para hacerlo para $h$, sólo habría que cambiar las letras). Entonces, $f\phi = g $ y $\phi = \frac{a}{b}f'$ con $f'$ un polinomio de $K[x]$ primitivo. Esto implica: $$\frac{a}{b}ff'=g \implies aff' = bg(1)$$

	Ahora, calcularemos los contenidos aplicando el lema de Gauss. Como f y $f'$ son primitivos,

	$$ a c(f)c(f') = bc(g) \implies a = bc(g)$$

	Ahora, si nos llevamos esta igualdad a (1) obtenemos:
	\[
	bc(g)ff'=bg \implies f(c(g)f')=g \implies f/g \ en \ A[x]
	\]
	Y por tanto, obtenemos que $f$ es primo.
\end{proof}

\begin{nota}
	De esta forma, podemos ver por inducción que cualquier anillo de la forma $A[x_1,...,x_r] $ es un DFU si A es un DFU.
\end{nota}
\begin{nota}[2]
	En $\ent$ hay infinitos primos.

\end{nota}

Vamos a intentar ahora a buscar una factorización de un $f \in A[x]$. ¿Cuándo es f irreducible?

 Sabemos que $gr(f) = 0 \iff f = p \in A$ que sería irreducible si $p$ lo es en A.

 Supongamos ahora que $gr(f) \geq 1$, $f$ es irreducible en $A[x] \iff$ f es primitivo e irreducible en $K[x]$

\begin{nota}
	Si K es un cuerpo, todo polinomio de grado 1 es irreducible.
	\begin{proof}
	Supongamos que $\phi = \phi_1 \phi_2 \implies 1 = gr(\phi) =  gr(\phi_1)+gr(\phi_2) \implies \phi_1$ es una unidad o lo es $\phi_2$, luego necesariamente $\phi$ es irreducible.
\end{proof}
\end{nota}


\textbf{Ejercicio.} Factorizar $(120x+100)^2$ en $\ent[x]$.

\[
120x+100 = 20(6x+5) = 2^2*5*(6x+5)
\]
Pero, $6x+5$ es un polinomio de grado 1, primitivo en $\ent[x]$. Además, 2 y 5 son irreducibles en $\ent$, por tanto la factorización del polinomio al cuadrado es:
\[
2^4*5^2(6x+5)^2
\]

\begin{nota}
	Si $K$ es un cuerpo, $\phi \in K[x]$ tiene un factor de grado 1 $\iff \phi$ tiene una raíz en $K$
	\begin{proof}
	$$\alpha x + \beta = \alpha(x-(\frac{-\beta}{\alpha}) = \alpha(x - \gamma) \implies \alpha x+b/\phi \iff x-\gamma / \phi \iff^{Rufinni} \phi(\gamma) = 0$$
	Donde $\gamma = \frac{-\beta}{\alpha}$
\end{proof}
\end{nota}

\begin{nprop}
	Sea $f=a_0+a_1x+...+a_nx^n \in A[x]$. Entonces $f$ tiene un factor irreducible de grado 1 en $K[x]$ si y solamente si tiene una raíz en $K$.

	Equivalentemente, $(a,b)=1$ y $f(\frac{a}{b})=0 \iff bx-a/f$ en $A[x]$. En tal caso, esta posible raíz verifica que $a/a_0$ y $b/a_n$ en $A$.
\end{nprop}
\begin{proof}
	Si tenemos que $f(\frac{a}{b})=0 \iff x-\frac{a}{b}/f$ en $K[x]$.

	$\boxed{\Rightarrow}$ Esto implica que $f=(x-\frac{a}{b})\phi$ con $\phi=\frac{c}{d}g$ con $g \in$ A primitivo, $$\implies f= \frac{c}{d}(x-\frac{a}{b})g = \frac{c}{bd}(bx-a)g \implies bdf=c(bx-a)g$$
	Si aplicamos contenidos:
	\[
	bdc(f)=c
	\]
	Volviendo a la igualdad anterior:
	\[
	f=(bx-a)(c(f)g)
	\]
	Por tanto, $bx-a/f$ en $A[x]$.

	$\boxed{\Leftarrow}$ Si $f=(bx-a)g \implies f(\frac{a}{b}) = (b\frac{a}{b}-a)g(\frac{a}{b})=0$
\end{proof}
\textbf{Criterio de la raíz:}
En $K[x]$ todo polinomio de grado 1 es irreducible y es asociado a uno de la forma $x-\alpha$ con $x-\alpha/\phi(x)\iff \phi(\alpha) = 0$.

Así, también en $A[x]$ los polinomios irreducibles de grado 1 son los de la forma $bx-a$ con $(a,b)=1$ por lo que $(a,b)=1$ y $bx-a/f(x)\iff f(\frac{a}{b})= 0$.

Si $f(x)= a_0+a_1x+...+a_nx^n$ y si $f(\frac{a}{b}) \implies a / a_0$ y $b/a_n$ en $A$.

Esto lo podríamos ver como una condición de irreducibilidad, pues si un polinomio fuera irreducible, no podría tener ninguna raíz, pues carecen de factores de grado 1.

\begin{ejemplo}
	Factorizar en producto de irreducibles el polinomio $f=6x^4+3x^3-18x^2+33x+21$ en $\ent[x]$.

Lo primero que deberíamos hacer es buscar el contenido de este polinomio:
\[
f=3(2x^4+x^3-6x^2+11x+7)
\]
El contenido de $f$ es 3 y $f'=2x^4+x^3-6x^2+11x+7$. 3 es un primo en $\ent $ y por tanto lo es en $\ent[x]$ y por tanto nos centramos en el primitivo asociado. Aplicamos el criterio de la raíz. Las posibles raíces que el polinomio tuviera en $\rac$ serían $\{\pm 1, \pm \frac{1}{2}, \pm 7, \pm \frac{7}{2}\}$. Calculamos las imágenes de estos puntos para ver si alguno es cero:
\[
f'(1) = 2+1-6+11 + 7 \ne 0; \quad f'(-1) = 2-1-6-11+7 \ne 0;\]
\[f'(\frac{-1}{2}) = 2*(1/2^4)-(1/2^3)-6*(1/4)-11*(1/2) + 7 = 0
\]
Y por tanto hemos encontrado una raíz, es decir:
\[
2x+1/f'(x) \ en \ \ent[x]
\]
Ahora, dividiremos $f'(x)$ entre $2x+1$ para reducirlo. Al dividirlo nos queda cociente $x^3-3x+7$ y resto 0. Por tanto, f nos queda factorizado como:

\[
f(x)=3*(2x+1)*(x^3-3x+7)
\]
Tenemos ya uno de grado 3, por tanto es irreducible si y solo si carece de raíces en $\rac$. Para ello, debemos buscar fracciones que sean divisores del 7, que son $\{\pm 1, \pm 7\}$ y haríamos sus imágenes por el polinomio $x^3-3x+7$ y vemos que ninguna da de resultado cero, por tanto estas no son raíces suyas y por tanto es irreducible en $\ent[x]$ y la factorización del polinomio en producto de primos sería justo la que acabamos de obtener.

Si nos planteamos la discusión en $\rac[x]$: 3 y $x^3-3x+7$ son mónicos, $2x+1$ no, pero podemos sacar factor común para hacerlo mónico y nos quedaría como resultado final:
\[
f(x)= 6(x+\frac{1}{2})(x^3-3x+7) \quad \in \rac[x]
\]
\end{ejemplo}
\begin{ejemplo}[2]
	Factorizaremos el polinomio: $20x^3 +10x^2-80x+30 \in \ent[x]$. Para resolverlo, había que observar que su contenido es 10 y que por tanto:
	\[
	f=2*5*(2x^3+x^2-8x+3)
	\]
	Y por tanto habría que buscar sus raíces entre $\{\pm 1, \pm \frac{1}{2}, \pm 3 \pm \frac{3}{2}\}$ y veríamos que en $\frac{3}{2}$ el polinomio evaluado da cero, por tanto factoriza como:
	\[
	2*5*(2x-3)(x^2+2x+2)
	\]
	Y podemos ver que $x^2+2x+2$ no tiene raíces y por tanto ese es el polinomio descompuesto en primos.
\end{ejemplo}
\begin{ejemplo}[3]
	Consideremos el polinomio: $x^4+x^3+x^2+x+1$ en $\ent_2[x]$. ¿Este polinomio es irreducible?. \\Estamos en $\ent_2$, luego sustituyendo el 0 y el 1 vemos que no tiene de raíces. Como no tiene raíces no tiene factores de grado 1.

	Ahora, si no fuese irreducible habría que descomponerlo como 2 factores de grado 2. Si dividimos por el polinomio $x^2+x+1$, que es irreducible en $\ent_2[x]$ el resto no es 0, luego no es divisible y por tanto el polinomio inicial ya es irreducible en $\ent_2[x]$.
\end{ejemplo}

\begin{nprop}[Criterio de Eisenstein]
	Sea $f(x)=a_0+a_1x+...+a_nx^n \in A[x]$ primitivo y siendo $A$ un DFU.
Si existe un primo $p\in A$ cumpliendo alguna de estas dos:
	\begin{enumerate}
	\item $p/a_i \ \ \forall i = 0,...,n-1$ y $p^2$ no divide a $a_0$
	\item $p/a_i \ \ \forall i=1,...,n$ y $p^2$ no divide a $a_n$
\end{enumerate}
Entonces $f$ es irreducible.

\end{nprop}
\begin{proof}
	Vamos a demostrarlo para el primer caso, la demostración para el segundo es equivalente. Vamos a negar la tesis y veremos qué ocurre.

	Supongamos que $f=(b_0+...+b_mx^m)(c_0+...+c_rx^r)$ con $r+m=n$ y $r,m \geq 1$. Entonces, $a_0=b_0c_0$ para que fuera el producto. Como $p$ es primo, tiene que dividir a $b_0$ o a $c_0$. Además, no puede dividir a ambos pues si no $p^2$ dividiría a $a_0$. Supongamos que $p/b_0$ y por tanto que no divide a $c_0$.

	Supongamos también que $p$ no divide a $a_n$ pues si no, formaría parte del contenido. Además, como hemos obtenido una factorización, $a_n = b_mc_r$ y como $p$ no divide a $a_n$ entonces no divide ni a $b_m$ ni a $c_r$.

	Sea $i$ el primer índice tal que $p$ no divide a $b_i$. Es claro que $0<i \leq m < n$, que existirá pues al menos es el último. Como lo habíamos factorizado, entonces $a_i = b_ic_0+(b_{i-1}c_1+...+b_0c_i)$. Entonces, $p/b_j \ \ \forall j < i \implies p/b_ic_0$ pero $p$ es primo luego divide o a $b_i$ o a $c_0$, pero ya habíamos dicho que no dividía a ninguno de los dos anteriormente en la demostración, luego hemos llegado a una contradición. Por tanto, f es irreducible.
\end{proof}

\begin{ejemplo}
	Podemos ver que $3x^7-70x^3+140 \in \ent[x]$ por el criterio de Eisenstein para el primo $p=5$, que divide a todos los coeficientes menos al coeficiente líder pero su cuadrado $p^2 = 25$ no divide a ninguno de los coeficientes.
\end{ejemplo}
\begin{ejemplo}[2]
	El polinomio : $y^3 + x^2y^2 +xy + x \in \ent[x,y]$, pues el primo $p=x$ divide a todos los coeficientes menos al líder y su cuadrado no divide a todos los coeficientes.
\end{ejemplo}

\begin{nprop}[Cambio de Anillo]
	Sean A y B dominios de integridad (así, el grado del producto es el producto de los grados). Sea $\phi : A[x] \to B[x]$ un homomorfismo en el que $gr(\phi(f))  \leq gr(f)$.

	Si $f\in A[x] : gr(\phi(f))= gr(f)$ y $\phi(f)$ no tiene divisores de grado $r \implies f$ tampoco tiene divisores de grado $r$
\end{nprop}
\begin{proof}
	Supongamos que $f=gh$ donde $gr(f) = n$, $gr(g) = r$ y $gr(h) = s$ con $r+s = n$.
	Aplicamos entonces el homomorfismo:
	\[
	\phi(f) = \phi(gh) = \phi(g)\phi(h)
	\]
	Y ahora, $gr(\phi(f)) = n$, $gr(\phi(g)) \leq r$ y $gr(\phi(h)) \leq s$ y entonces implica que $gr(\phi(g)) = r$ (Si $r < gr(\phi(g)) \implies r+s \neq n$, contradicción)
\end{proof}
\begin{ejemplo}
	$x^4+1$ es irreducible en $\ent[x]$.

	Si tomamos $\phi:\ent[x] \to \ent[x]$ que lleva $x \mapsto x+1$ el homomorfismo de evaluación en $x+1$. ($\phi(\sum a_i x^i) = \sum a_i (x+1)^i$).

	Así, $\phi(x^4+1) = (x+1)^4+1 = x^4 +4x^3 +6x^2 +4x + 2 \in \ent[x]$, que es irreducible por el criterio de Eisenstein para el primo $p=2$. Ahora, por el criterio del cambio de anillo, entonces $x^4+1$ es también irreducible en $\ent[x]$
\end{ejemplo}

\begin{nprop}[Criterio de reducción módulo un primo]
	En las condiciones anteriores, consiste en aplicar el homomorfismo:
	\[
	R_p:\ent[x] \to \ent_p[x] : \quad R_p(\sum a_i x^i) =  \sum R_p(a_i) x^i \in \ent_p[x]
	\]
	Para utilizar este criterio, es muy útil conocer los irreducibles mónicos de los anillos $Z_p[x]$
\end{nprop}

\textbf{Irreducibles mónicos de $\ent_p[x]$:}
\begin{itemize}
	\item de grado 2 en $\ent_2[x]$: $x^2+x+1$
	\item de grado 3 en $\ent_2[x]: x^3+x+1,\ x^3+x^2+1$
	\item de grado 2 en $\ent_3[x]: x^2+1,\ x^2+x+2,\ x^2+2x+2$
\end{itemize}
\begin{ejemplo}
	Sea $f = x^4 +3x^2 -2x +5 \in \ent[x]$. Si le aplicamos el criterio de reducción módulo un primo, en este caso el 2, el polinomio queda: $R_2(f) = x^4 + x^2 +1 \in \ent_2[x]$ y por el criterio de reducción, como este polinomio no tiene divisores de grado ni 1 ni 3 entonces podemos asegurar que $f$ tampoco los tiene.

	Si lo volvemos a hacer para el primo $p=3$, queda $R_3(f) = x^4 + x + 2 \in \ent_3[x]$ , que nos da exactamente la misma información que la reducción anterior. Como conocemos los mónicos de $\ent_3[x]$, dividimos este polinomio entre los 3 irreducibles mónicos de $\ent_3[x]$ y vemos que los restos son no nulos. Así, este polinomio es irreducible en $\ent_3[x]$ y por el criterio de reducción módulo un primo, el polinomio primero es irreducible en $\ent[x]$.
\end{ejemplo}



\section{Matrices sobre Anillos conmutativos}
Sea $K$ un cuerpo, y consideremos $\mathcal{M}_{mxn}(K)$ y  $\mathcal{M}_{n}(K)$. Imaginemos ahora que $A\leq K$ es un subanillo. De esta forma, podemos concebir las matrices que están en este conjunto: $\mathcal{M}_{mxn}(A)$ ó  $\mathcal{M}_{n}(A)$ .

Así, si sumáramos dos matrices que estuvieran en  $\mathcal{M}_{mxn}(A)$, nos quedaría una matriz cuyas entradas están en el subanillo A. Lo mismo ocurre con la multiplicación.


Podemos considerar también:
\[
GL_n(K) = U (\mathcal{M}_n(K)) = \{ M \in \mathcal{M}_n(K) :  \ \exists M^{-1} \in \mathcal{M}_n(K) \ \ con \ \ MM^{-1} = M^{-1}M = I_n\}
\]

Que sabemos que es igual al conjunto $\{M \in \mathcal{M}_n(K) : | M | \ne 0\}$. Ahora, podemos considerar, como hemos hecho anteriormente, lo mismo pero en el subanillo A.

Dada $M \in \mathcal{M}_n(A)$, si $M \in GL_n(A)$, cumplirá que $MM^{-1} = I$, tomando determinantes vemos que: $|I|=|M||M^{-1}|$. El determinante de M y el determinante de $M^{-1}$ están en A, y es un producto que da como resultado 1, por lo que podemos asegurar que $|M| \in U(A)$.

Veámoslo al revés. Si $M\in \mathcal{M}_n(A)$ con $|M| \in U(A) \implies 1/|M| \in A$. Si consideramos la adjunta de M, esta también pertenece a $\mathcal{M}_n$ y, por tanto, $M \in GL_n(A)$ pues tiene una inversa cuyas entradas están todas en el subanillo A.
Con todo esto, hemos deducido la siguiente proposición:

\begin{nprop}
	Sea $M \in \mathcal{M}_n(A)$, entonces $M \in GL_n(A) \iff |M| \in U(A)$
\end{nprop}

\begin{ejemplo}
	El grupo lineal de $\ent$ es: $GL_n(Z) = \{M \in \mathcal{M}_n(A) :  |M| = \pm 1\}$
\end{ejemplo}

\begin{ejemplo}[2]
	El grupo lineal de $K[x]$ es: $GL_n(K[x]) = \{M \in \mathcal{M}_n(K[x]) :  |M| = k - \{0\} \},$ donde k es un polinomio de grado 0 no nulo.
\end{ejemplo}

\begin{ndef}[Matrices equivalentes]
	Dos matrices M y N $\in \mathcal{M}_{mxn}(A)$ son equivalentes si $\exists P\in GL_m(A), Q \in GL_n(A)$ tal que
	\[
	M = PNQ
	\]
Estando en un cuerpo, dos matrices serán equivalentes $\iff$ tienen el mismo rango.
\end{ndef}

\begin{nth}[Teorema de la forma Normal de Smith]
	Si A es un Dominio Euclídeo, toda matriz  $M \in \mathcal{M}_{mxn}(A)$ es equivalente a una de la forma:
		\[ \begin{pmatrix}
			 d_{1}&  &  & & \\
			 &  & \ddots &  &  \\
			 & & &d_{r} &\\
			 & & & & 0 \\
			 & & & & & \ddots& \\
			 & & & &  & & 0 \\
		\end{pmatrix} \]
con $d_{1},...,d_{r} \in A$ ubicados en la diagonal, donde $ 0 \leq r \leq min\{m,n\} $ y cada $d_{i}\ne 0$ y $d_{i}/ d_{i+1}$.  Además, estos $d_{1}, ...,d_{r}$ son únicos para la clase de equivalencia de M y se llaman los \textbf{factores invariantes} de la matriz.
\end{nth}

\begin{ncor}
	Dos matrices $M,N \in \mathcal{M}_{mxn}(A)$ son equivalentes $\iff$ tienen los mismos factores invariantes.
\end{ncor}

Sea esta matriz

	\[P= \left( \begin{array}{ccccc}
 1&  &  & & \\
 &   \ddots & & a \\
 &  &  &  &  \\
 & & & & 1
 \end{array} \right)\]
Que tiene en la posición (i,j) un elemento $a\in A$. Entonces, esta matriz es invertible y además su inversa es igual pero cambiando $a$ por $-a$.

Si:
	\[ N =  \left( \begin{array}{c}
 F_1 \\
 \vdots\\
 F_n
 \end{array} \right)   \]
 Y multiplicáramos $P_1N$, entonces quedaría $F_k' = F_k $ si $k \ne i$ con $i$ el número de fila que contenía el elemento $a$ y $F_i'=F_i+aF_j$. Lo mismo ocurriría con las columnas si...\\
\textbf{FALTA CONTENIDO DE MATRICES, añadirla.}



\begin{nota}
	Si multiplicáramos en una matriz todos los elementos de una fila o los de una columna por una unidad del anillo, entonces la matriz inicial y la resultante son equivalentes.
\end{nota}
\begin{nota}
	Si permutamos dos filas, o dos columnas, las matrices inicial y final son equivalentes.
\end{nota}

\begin{ndef}
	Una matriz $M \in \mathcal{M}_{mxn}$ que sea:
		\[ \begin{pmatrix}
 d_{1}&  &  & & \\
 &  & \ddots &  &  \\
 & & &d_{r} &\\
 & & & & 0 \\
 & & & & & \ddots& \\
 & & & &  & & 0 \\
\end{pmatrix} \]
 con $d_1,...,d_r \in A, \ \ d_i \ne 0$ y $d_i / d_{i+1}$, en las posiciones (i,i) con $1 \leq i \leq r$ y 0 en todas las demás posiciones, se dice matriz Normalizada (Smith).
\end{ndef}

Vamos a explicar ahora un método a partir del cual, teniendo una matriz cualquiera podemos obtener una matriz normalizada que es equivalente a la primera. No olvidemos que estamos trabajando en Dominios euclídeos.

Sea $A \in \mathcal{M}_{mxn}$. $A \ne 0$. Sea $\varphi:A -\{0\} \to \nat$ la función euclídea.

Vamos a llamar $\varphi(M) = min\{\varphi(a_{ij}) : a_{ij} \ne 0\}$
Primero, tomaremos un $a_{ij}$ donde la función euclídea tome el valor más pequeño que tome en toda la matriz. Ahora, por operaciones elementales podemos llevarlo a la posición $(1,1)$.

Supongamos ahora que $a_{11}$ no divide a $a_{1k}$, como estamos en un dominio euclídeo, podemos expresar el cociente entre ellos como:
\[
a_{1k} = a_{11}q+ b_{1k}
\]
donde $b_{1k} \ne 0$ y $\varphi(b_{1k}) < \varphi(a_{11})$.

Ahora, realizamos la operación:
\[
C_k \mapsto C_k - q C_1
\]

		\[ \begin{pmatrix}
			 a_{11} \hdots  a_{1k} \hdots \\
			   \hdots &  \\
		\end{pmatrix}
		\xrightarrow{C_k - qC_1}
		\begin{pmatrix}
			 a_{11} \hdots  b_{1k} \hdots \\
			   \hdots &  \\
		\end{pmatrix}
\]
Reiteramos el proceso las veces que haga falta, hasta encontrarnos una matriz que sea de la forma $(b_{ij})$. Ahora, si $b_{11}$ no divide a $b_{k1}$, entonces este último se expresará como $b_{k1} = b_{11}q' + c_{k1}$ con $c_{k1}\ne 0$ y $\varphi(c_{k1}) < \varphi(b_{11})$

Entonces, realizamos la operación:
\[
F_k \mapsto F_k - q'F_1
\]
Volvemos a reiterar el procedimiento cuantas veces haga falta.

Tras una serie de pasos, llegaremos a una matriz equivalente a la original en la que $b_{11}/b_{k1}$ y $b_{11}/b_{1k} \ \ \forall k$. Llegados a esta situación, por transformaciones elementales podemos obtener una matriz que sea del tipo:
\[
\begin{pmatrix}
 b_{11} & 0 & \cdots & 0 \\
 0 & c_{22} & \cdots & c_{2n} \\
 \vdots &  & \ddots& \\
 0 & c_{m2} & \cdots & c_{mn} \\
\end{pmatrix}
\]

Donde $b_{11}/ c_{ij}$ para cada $i$ y $j$. Ahora, lo que haremos es tomar la submatriz que empieza en $c_{22}$ y volveremos a realizar el mismo proceso.

\begin{ejemplo}
	Sea $A \in \mathcal{M}_{3x2}$ tal que:
	\[
	A =
	\begin{pmatrix}
 14 & 10 \\
 22 & 14\\
 10 & 10\\
\end{pmatrix}
	\]
	Ahora, tomamos el número que tenga $\varphi$ menor, en este caso el 10 y hacemos transformaciones elementales para llevarlo a la primera posición.
	\[\begin{pmatrix}
 14 & 10 \\
 22 & 14\\
 10 & 10\\
\end{pmatrix} \xrightarrow{C_2,\ C_1}
	\begin{pmatrix}
 10 & 14 \\
 14 & 22\\
 10 & 10\\
\end{pmatrix} \xrightarrow{C_2\ -\ C_1} \begin{pmatrix}
 10 & 4 \\
 14 & 8\\
 10 & 0\\
\end{pmatrix}\xrightarrow{C_2,\ C_1}
\begin{pmatrix}
 4 & 10 \\
 8 & 14\\
 0 & 10\\
\end{pmatrix}\xrightarrow{C_2\ -\ 2C_1}  \begin{pmatrix}
 4 & 2 \\
 8 & -2\\
 0 & 10\\
\end{pmatrix}
	\]
	\[
	\xrightarrow{C_2\ -\ C_1}
	\begin{pmatrix}
 2 & 4 \\
 -2 & 8\\
 10 & 0\\
\end{pmatrix} \xrightarrow{C_2\ -\ 2C_1}
\begin{pmatrix}
 2 & 0 \\
 -2 & 12\\
 10 & -20\\
\end{pmatrix}\xrightarrow{F_2\ +\ F_1}
\begin{pmatrix}
 2 & 0 \\
 0 & 12\\
 10 & -20\\
\end{pmatrix} \xrightarrow{F_3\ -\ F_1}
\begin{pmatrix}
 2 & 0 \\
 0 & 12\\
 0 & -20\\
\end{pmatrix}
	\]
	Y ahora, tendríamos que hacer lo mismo atendiendo a la submatriz resultante:
	\[
	\begin{pmatrix}
 2 & 0 \\
 0 & 12 \\
 0 & -20 \\
\end{pmatrix}\xrightarrow{F_3\ +\ 2F_2}
\begin{pmatrix}
  2 & 0 \\
  0 & 4 \\
  0 & 12 \\
\end{pmatrix}\xrightarrow{F_3\ -\ 3F_2}
\begin{pmatrix}
  2 & 0 \\
  0 & 4 \\
  0 & 0 \\
\end{pmatrix}
	\]
	Y nos queda la matriz:
	\[
	\begin{pmatrix}
 2 & 0 \\
 0 & 4\\
 0 & 0\\
\end{pmatrix}
	\]
	Que es equivalente a la inicial.
\end{ejemplo}

\begin{ndef}[Rango de la matriz]
	El entero $r$ tal que $M$ tiene al menos un menor de orden $r$ no nulo y todos los de orden $r+1$ son nulos se llama Rango de la matriz.
\end{ndef}

Vamos a definir ahora $\Delta_s(M) = mcd\{$ menores de orden s de M no nulos$\}$
Así, $\Delta_1(M)$ sería el máximo común divisor de todos los $m_{ij}$ (si $M$ está normalizada será $d_1$). Ahora, con esta definición, podemos decir que:
\[
rg(M) = max\{s : \Delta_s(M) \ne 0\}
\]
Lo que implica que si $rg(M) = r \implies \Delta_1(M),...,\Delta_r(M)$ son $\ne 0$ y $\Delta_{r+1}(M),...$ son nulos.

Sea $P \in \mathcal{M}_m(A)$. Construimos la matriz $PM\in \mathcal{M}_{mxn}(A)$ (tenemos la matriz $M \in \mathcal{M}_{mxn}(A)$). Si miramos el elemento de PM en la posición $(k,j)$, este es: $\sum_{i=1}^m p_{ki}a_{ij}$. Así, la fila $k-$ésima de $PM$ es:
\[
(\sum_{i=1}^m p_{ki}a_{i1}, \cdots , \sum_{i=1}^m p_{ki}a_{in}) = \sum_{i=1}^m p_{ki}(a_{i1},\cdots , a_{in})
\]
Por tanto, las filas de la matriz $PM$ son combinación lineal de la matriz M. También podemos afirmar que los menores de orden $s$ de la matriz PM son combinación lineal de los menores de orden $s$ de la matriz M. Entonces, $\Delta_s(M)/\Delta_s(PM)$.

Si $P\in GL_m(A)\implies \Delta_s(PM) / \Delta_s(P^{-1}PM) = \Delta_s(M) \implies \Delta_s(M) = \Delta_s(PM)$

Si $Q \in GL_n(A) \implies \Delta_s (M) = \Delta(MQ)$

\begin{nota}
	$\Delta(M) = \Delta(M^t)$. Del mismo modo, $\Delta_s(MQ) = \Delta_s((MQ)^t) =  \Delta_s (Q^tM^t) = \Delta_s (M^t) =  \Delta_s(M)$
\end{nota}

De las implicaciones anteriores, obtenemos como conclusión:

$\forall \ P,\ Q \ \ \ \Delta_s(PMQ) = \Delta_s(M) \ \ \forall s \implies$
si $M$ y $M'$ son equivalentes, entonces $\Delta_s(M) = \Delta_s(M') \ \forall s$\\
Si $M \sim M' \sim \begin{pmatrix}
 d_{1}'&  &  & & \\
 &  & \ddots &  &  \\
 & & &d_{r}' &\\
 & & & & 0 \\
 & & & & & \ddots& \\
 & & & &  & & 0 \\
\end{pmatrix}  \implies r = r'$

Tenemos que $\Delta_1(M) = d_1 = \Delta_1 (M') = d_1' \implies d_1 = d_1'$ ; y si lo hacemos en general vemos que:
\[
\Delta_s(M) = d_1...d_s = \Delta_s(M') = d_1'...d_s' \implies d_s =  \frac{\Delta_s(M)}{\Delta_{s-1}(M)} =  \frac{\Delta_s(M')}{\Delta_{s-1}(M')} = d_s'
\]

\begin{nota}
	Si el anillo es un cuerpo $\implies d_i = 1 \ \ \forall i$
\end{nota}

\section{A-módulo}
\begin{ndef}[A-módulo]
	Es un conjunto no vacío M dotado de dos operacions:
	\begin{itemize}
	\item $MxM \to M$ tal que $(x,y) \mapsto x+y$ que es interna
	\item $Ax M \to M$ tal que $(a,x) \mapsto ax$
\end{itemize}
que verifican las propiedades asociativa, conmutativa, elemento neutro y elemento opuesto. También verifican la distributiva respecto de la suma y respecto del producto, asociatividad del producto y neutro del producto.

\end{ndef}
Durante todo el apartado nos referiremos a $x, y,...$ como elementos de $M$ y a $a, b, ...$ como elementos de A.\\
\textbf{Propiedades:}
\begin{nlist}
	\item Unicidad del 0 y del opuesto.
	\item $0x = 0$ y $a0 = 0$
	\item $(-a)x = -(ax) = a(-x)$
	\item Si para una lista de elementos $(x_1,...,x_n)$ de M con al menos 2 elementos definimos su suma simultánea $\sum_{i=1}^n x_i$ por inducción:$ \sum_{i=1}^n x_i = (\sum_{i=1}^{n-1}x_i) +x_n$
	\item $\forall i < r < n,$ entonces $\sum_{i=1}^n x_i = (\sum_{i=1}^{r}x_i) + (\sum_{i=r+1}^{n}x_i) $
	\item $ (\sum_{i=1}^n a_i)(\sum_{j=1}^m x_i) = \sum_{i=1}^n \sum_{j=1}^m a_ix_j$
\end{nlist}

Si $I \leq A$ es un ideal, si tomamos $M= A/I$, entonces se verfican:
\begin{itemize}
	\item $[x]+[y] = [x+y]$
	\item $a[x] = [ax] = [a][x]$
\end{itemize}

\begin{ndef}[Submódulos]
	Si M es un A-Módulo, un submódulo es $\emptyset \ne N \subseteq M$ que es cerrado para sumas ($x,y\in N \implies x+y \in N$) y para múltiplos($a\in A,\ x\in N \implies ax\in N$).

	Con estas operaciones de suma y producto restringidas, $N$ también es un A-módulo. Por tanto, todo submódulo es un módulo.
\end{ndef}
\begin{nprop}
	Si $M_1,...,M_r \subseteq M$ son submódulos, existe siempre un módulo que los contiene. Se le llama submódulo suma y se representa como:
	\[
	\sum_{i=1}^r M_i =\{\sum_{i=1}^r x_i : x_i \in M_i\}
	\]
	Además, se puede usar la conmutatividad de la suma y la asociatividad generalizada.

\end{nprop}
\begin{ndef}[Suma directa]
	Se dice que la suma de $M_1,...,M_r$ es directa y se representa:
	\[
	\sum_{i=1}^r = \oplus_{i=1}^r M_i \ si \ (\sum_{i=1}^r x_i =  \sum_{i=1}^r y_i) \ con \ x_i,y_i \in M_i \iff x_i = y_i\ \forall i	\]
\end{ndef}
\begin{ndef}[Submódulo cíclico]
	Si $x\in M$, siempre existe un submódulo de $M$ que es el menor submódulo que lo contiene. Más aún, que contiene al elemento $x$ y está contenido en cualquier otro submódulo que contenga a $x$. Se denota como $Ax = \{ax : a \in A\}$, pues es donde están todos los múltiplos del elemento. Se le llama el submódulo cíclico de x. También se suele notar “(x)”.
\end{ndef}
\begin{ndef}
	Si $x_1,...,x_r \in M$, existe un módulo $M$ que es el más pequeño módulo que los contiene. Este es:
	\[
	\sum_{i=1}^r Ax_i = \{\sum_{i=1}^r a_ix_i : a_i \in A\}
	\]
	Estos son combinaciones lineales de $x_1,...,x_r$, se le llama submódulo generado por $x_1,...,x_r$ y se representa como $(x_1,...,x_r)$.
\end{ndef}
\begin{nota}
	Si $N=(x_1,...,x_r)$ y $N'= (y_1,...,y_s)$, entonces
	\[
	N \subseteq N' \iff todo\ x_i \ es \ c.l de \ y_1,...,y_s \ \forall i
	\]
	Y $N=N'\iff$ todo $x_i$ es combinación lineal de $y_1,...,y_s$ y todo $y_j$ es combinación lineal de $x_1,...,x_r$
\end{nota}
\begin{ndef}[Sistema de generadores]
	Si $x_1,...,x_r\in M : M = (x_1,...,x_r)$ se dice que $x_1,...,x_r$ es un sistema de generadores de M  y se dice que M es finitamente generado.
\end{ndef}
\begin{nprop}
	Si A es un dominio euclídeo, todo submódulo de un módulo finitamente generado es finitamente generado.
\end{nprop}
\begin{ndef}
	Un sistema de generadores $x_1,...,x_r$ de un módulo $M$ es una base de $M$ si se verifica:
	\[
	\sum_{i=1}^r a_i x_i = \sum_{i=1}^r b_i x_i \iff a_i = b_i \ \forall i
	\]
	O, equivalentemente:
	\[
	\sum_{i=1}^r a_ix_i = 0 \iff a_i  = 0  \ \forall i
	\]
\end{ndef}
\begin{ndef}[Linealmente independientes]
	Un conjunto de elementos de un A-Módulo $x_1,...,x_r$ se dice que son linealmente independientes si la única combinación lineal de ellos que da como resultado 0 es aquella en la que todos los coeficientes son cero.
\end{ndef}
\begin{nota}
	En los A-módulos, es falso decir que todo A-módulo tiene una base o que todo sistema de generadores contiene a una base. Tampoco podemos siempre extender un conjunto linealmente independiente a una base.
\end{nota}
\begin{ndef}[Módulo Libre]
	Un módulo que contenga una base es un módulo libre.
\end{ndef}

\begin{ndef}
	Si M, N son dos módulos, y $\varphi:M \to N$ una aplicación, la llamamos lineal u homomorfismo entre módulos si:
	\[
	\begin{rcases}
	\varphi(x+y) = \varphi(x) + \varphi(y)\\
	\varphi(ax) = a\varphi(x)
\end{rcases} \iff \varphi(\sum_{i=1}^n a_ix_i) = \sum_{i=1}^n a_i \varphi(x_i)
	\]
	También puede ser un monomorfismo (si y solo si su núcleo es 0), epimorfismo (es sobreyectiva) o un isomorfismo (es biyectiva).
\end{ndef}
\begin{ejemplo}
	Si $m\in A$, $A/(m)$ es un A-módulo ( $[a]+[b] = [a+b]$ y $a[b] = [ab]$. Si A tiene unicidad de cocientes y restos ($\ent,K[x]$), entonces $A_m$ es un A-módulo con $r+r' = R_m(r+r')$ y $ar = R_m(ar)$.

	Recordemos que existe un isomorfismo entre $A/(m) \cong A_m$ que lleva $[a] \mapsto R_m(a)$
\end{ejemplo}

\begin{ejemplo}[2]
	Si $M_1,...,M_n$ son A-módulos, su “producto" $M_1 \times ... \times M_n =  \prod M_i$ es un A-módulo con $(x_1,...,x_n) + (y_1,...,y_n) = (x_1+y_1,...,x_n+y_n)$  y $a(x_1,...,x_n)= (ax_1,...,ax_n)$

		Supongamos que $M_1,...,M_n$ son submódulos de M y que $M = M_1 \oplus ... \oplus M_n = \oplus_i M_i$
	Entonces, existe:
	\begin{align}
	\varphi:\prod_i M_i \to M=\oplus_i M_i \\
	(x_1,...,x_n) \mapsto x_1+...+x_n = \sum_i x_i
\end{align}
que es un isomorfismo.
\end{ejemplo}

\section{Clasificación de los módulos finitamente generados sobre un Dominio Euclídeo.}

Para comprenderlo mejor, vamos a ver primero un caso particular sencillo.
\subsection{Módulos cíclicos}
\begin{ndef}[Anulador minimal de M]
	Si M es cualquier módulo, entonces
	\[
	\{a: ax = 0 \ \forall x \in M\}\subseteq A \ \  \text{es un ideal}
	\]
	a cuyo generador llamamos anulador minimal de M, $\mu(M)$, que es único salvo asociados.
\end{ndef}

Dado un cierto $x\in M$, podemos ver el conjunto $\{a : ax = 0\} \subseteq A$. Este también es un ideal y, a su vez, es principal. Al generador de este se le llama el anulador minimal de x y se le denota:
\[
\mu(x)
\]

\begin{nprop}
	Si expresamos $M= (x_1,...,x_r) = Ax_1 + ...+ Ax_r \implies \mu(M) = mcm(\mu(x_1),...,\mu(x_r))$
\end{nprop}
\begin{proof}
	$\forall i \ \mu(M)x_i = 0 \implies \forall i=1,...,r, \ \ \mu(x_i) / \mu(M)$.

	Por tanto, si $a\in A: \mu(x_i)/a \ \forall i \implies ax_i = 0 \ \forall i \implies a\sum_i a_ix_i = \sum_i a_i(ax_i) = 0$
	$\implies ax = 0 \ \forall x \in M \implies \mu(M) / a$
\end{proof}

En particular, si estamos en un módulo cíclico (generado por un único elemento), entonces $\mu(M) = \mu(x)$ salvo asociados.

\begin{nota}
	En estos casos, si $\mu(x) = 0$, la aplicación $\varphi: A \to M$ que lleva a cada elemento de A en M de forma $a \mapsto ax$, es lineal, sobreyectiva y es inyectiva, pues $\ker(\varphi) = \{0\}$ (el anulador minimal es el cero), por tanto, es un isomorfismo, $M \cong A$.
\end{nota}

\begin{nota}
	Si $\mu(M) = \mu(x) \ne 0$, la aplicación $\varphi :A/\mu(x) \to M$ tal que $[a]\mapsto ax$ está bien definida ( si $a \equiv b \mod(\mu(x)) \implies a = b+q\mu(x) \implies ax = bx +q\mu(x)x = bx$), es lineal, sobreyectiva e inyectiva (si $\varphi([a]) = 0 \implies ax = 0 \implies  \mu(x)/a \implies [a] = 0$) , por lo que
	\[
	M  \cong A/(\mu(M))
	\]
\end{nota}

\begin{ndef}[Presentacion de M por generadores y relaciones]
	Sea M un A-módulo finitamente generado. Se llama entonces “Presentación de M por generadores y relaciones” a toda expresión de la forma
	\[
	\langle x_1,...,x_n : \sum_{i=1}^n a_{i1}x_i = 0,...,\sum_{i=1}^n a_{im}x_i = 0 \rangle
	\]
	Donde $(x_1,...,x_n)$ es un sistema de generadores de $M$ y las expresiones de la derecha son relaciones de dependencia lineal entre ellos verificadas en el módulo tales que cualquier relación de dependencia entre $x_1,...,x_n$ es combinación lineal de las que aparecen en la presentación

\end{ndef}

\begin{nprop}
		Si
	\[
	\langle x_1,...,x_n : \sum_{i=1}^n a_{i1}x_i = 0,...,\sum_{i=1}^n a_{im}x_i = 0 \rangle
	\]
	es una presentación del módulo M, a la matriz:
	\[
	\mathcal{M}_R = \begin{pmatrix}
 a_{11} & \cdots & a_{1n}\\
 \vdots & & \vdots\\
 a_{m1} & \cdots & a_{mn}
\end{pmatrix}
	\]
	cuyas filas son exactamente los coeficientes que aparecen en la presentación la llamamos “matriz de relaciones de los generadores de M".

	Decir entonces que las combinaciones lineales de la presentación son cero, es lo mismo que decir:
	\[
	\begin{pmatrix}
 a_{11} & \cdots & a_{1n}\\
 \vdots & & \vdots\\
 a_{m1} & \cdots & a_{mn}
\end{pmatrix} \begin{pmatrix}
 x_1  \\
 \vdots\\
 x_n
\end{pmatrix} =  \begin{pmatrix}
 0  \\
 \vdots\\
 0
\end{pmatrix}
	\]
\end{nprop}


	Si tenemos un módulo $A^n$, podemos ver que este es libre y finitamente generado, además de que tiene una base (la canónica de $\mathbb{R}^n$). Vamos a considerar una aplicación $\varphi: A^n \to M$ tal que $\varphi:(a_1,...a_n) = a_1x_1+...+a_nx_n$. Esta aplicación es un epimorfirmo.

\begin{nprop}
	\[
	\sum_{i=1}^n a_i x_i  = 0 \iff (a_1,...,a_n) \in \ker(\varphi)
	\]
	En particular, que se verifiquen las relaciones de dependencia que aparecen en la presentación, significa que $F_1,...,F_m \in \ker(\varphi)$.
	Y decir que cualquier relación de dependencia entre $x_1,...,x_n$ es combinación lineal de las que aparecen en la presentación significa que si $F=(a_1,...a_n) \in \ker(\varphi) \implies F$ es combinación lineal de $F_1,...,F_m$, por tanto, estamos diciendo que
	\[
	\ker(\varphi) = (F_1,...,F_m)
	\]
\end{nprop}

\begin{nota}
	Es importante recordar que estos sistemas de generadores pueden ser sometidos a operaciones elementales sin dejar de ser sistemas de generadores. Las operaciones pueden ser:
	\begin{itemize}
	\item Permutar dos de ellos
	\item Multiplicar uno por una unidad del anillo
	\item Sumar a uno de ellos un múltiplo de otro
\end{itemize}
Aún así, la matriz de relaciones se verá alterada y tendremos “otra presentación distinta".
\end{nota}

Realizando estas operaciones, siempre podemos llegar a obtener la forma normal de Smith de una matriz.

\begin{nprop}
	$\forall M $, existe una presentación cuya matriz de relaciones está normalizada (tiene la forma de una matriz de Smith). Es decir, es de la forma:
	\[
	M : \langle \ \  y_1,...,y_n : \ d_1y_1 = 0 \ ,..., \ d_ry_r = 0 \ \  \rangle
	\]

\end{nprop}
Es más, si ahora consideráramos un nuevo epimorfismo $\varphi: A^n \to M$ tal que $\varphi (a_1,...a_n) = \sum a_i y_i$ y
\[
\ker(\varphi)= \ \langle\  d_1e_1 \ \ , .... , \ \ d_r e_r \ \rangle
\]
con $e_i$ el vector de la base canónica (tiene un 1 en la posición i-ésima y lo demás son ceros).

\begin{nprop}
	Se puede comprobar que:
	\[
	ay_i = 0 \implies \varphi (ae_i) = 0 \implies ae_i \in \ker(\varphi) \implies ae_i = b_1d_1e_1+...+b_rd_re_r
	\]
	\[
	\implies b_1d_1e_1+...+b_rd_re_r - ae_i = 0
	\]
	\begin{itemize}
	\item  si $i>r \implies a = 0 \implies \mu(y_i) = 0$.
	\item si $i \leq r \implies a = b_id_i$ pues son todos linealmente independientes $\implies d_i/a \implies \mu(y_i) = d_i$
\end{itemize}
\end{nprop}

\begin{nprop}
	$M \oplus_i Ay_i$, todo elemento de M se expresa de forma única como combinación lineal de elementos de $A$
\end{nprop}
\begin{proof}
	Supongamos que $\sum a_i y_i = 0 \implies \varphi(\sum a_i e_i ) = 0 \implies \sum a_ie_i \in \ker (\varphi) $\\
	$\implies a_1e_1+...+ a_re_r+...+a_ne_n = b_1d_1e_1+...+b_rd_re_r \implies a_i = 0 \ \ \forall i > r$\\
	$\implies a_iy_i = 0$

	Y si $i \leq r\implies a_i = b_id_i \implies a_iy_i = 0$
\end{proof}
También podemos decir, si $d_1 = ... = d_s = 1 \implies y_1 = y_2 = ... = y_s = 0$ y, por tanto, tendríamos que:
\[
M = Ay_{s+1}  \oplus ... \oplus Ay_r \oplus ... \oplus Ay_n
\]
Y tenemos que $\mu(y_i) = d_i \in A$, no es cero ni unidad $\forall i = s+1,...,r$ y cada $d_i/d_{i+1}$ y $\mu(y_{r+1}) = ... =\mu(y_n) = 0$

\begin{nth}[Teorema de estructura de módulos cíclicos]
	En estas condiciones,
	\[
	M \cong Ay_{s+1} \times ... \times Ay_r \times ... \times Ay_n
	\]
	Siendo $Ay_{s+1},..., Ay_r$ módulos cíclicos y desde $r+1$ hasta $n$, $Ay_j$ es isomorfo al propio anillo, por lo que podemos decir que:
	\[
	M \cong Ad_{s+1}\times ... \times Ad_r \times A^{n-r}
	\]
\end{nth}

A $n-r$ se le llama el rango del módulo y $d_{s+1},...,d_r$ es la lista de los factores invariantes del módulo M

La unicidad de estos invariantes es consecuencia de la unicidad de los invariantes de la forma normal de Smith.

\begin{ncor}
	Para conocer un módulo, basta conocer la lista de sus factores invariantes y su rango.
\end{ncor}
\begin{ncor}
	Si el rango de un módulo ($n-r$) es mayor que cero, entonces $\mu(M) = 0$.\\

	Si $n-r$ = 0, entonces $\mu(M) $ es el mayor de los factores invariantes, el $d_r$.
\end{ncor}

A continuación, ilustraremos el teorema que hemos visto con varios ejemplos:

\subsubsection{$\ent$-Módulos. Grupos abelianos.}

Vamos a tratar de llevarnos los conceptos que hemos obtenido con el teorema de estructura al caso de los $\ent-$módulos.\\
Si M es un $\ent$-módulo, $n\in \ent$ con $n>0$ y $x\in M$, entonces $nx = (\sum_{i=1}^n 1)x = \sum_{i=1}^n 1*x = \sum_{i=1}^n x$. Lo mismo podemos ver si lo hacemos con $n$ negativo.

Entonces, un $\ent$-módulo es un grupo abeliano. Hablaremos de aquí en adelante de grupos abelianos, aunque sabremos que es lo mismo que un $\ent-$módulo.

Si $M$ es un grupo abeliano y $x\in M$, entonces $\ent x = (x) = \{nx : n \in \ent  \}$.
Tenemos también que $\mu(x) = d \iff dx = 0$ y $mx = 0 \implies d/m$.

Si $M=\ent x$ y $\mu(x) = 0 \implies M \cong \ent$

\begin{nth}[Teorema de estructura en grupos abelianos]\hfill \\
	Para todo grupo abeliano finitamente generado $M$ existen enteros positivos $d_1,...,d_s$ con $d_i \geq 2$ tales que $d_i/d_{i+1}$ y existe un $r> 0$ tal que
	\[
	M \cong \ent_{d_1} \times ... \times \ent_{d_s} \times \ent^n
	\]
	y a estos $d_i$ se les llama factores invariantes del grupo abeliano y $r$ es su rango.
	A la descomposición que hemos obtenido se le llama descomposición cíclica del grupo abeliano.
\end{nth}

\begin{nprop}
	Dos grupos abelianos son isomorfos $\iff$ tienen los mismos factores invariantes y el mismo rango. Por tanto, son isomorfos $\iff$ tienen la misma descomposición cíclica.
\end{nprop}

\begin{nprop}
	$M$ es finito $\iff$ su rango es 0. En ese caso, el tamaño del grupo es el producto de sus invariantes  $|M| = d_1...d_s$ y en tal caso $\mu(M) = d_s$
\end{nprop}
\begin{nprop}
	Si el rango de un módulo $M$ no es cero, entonces $\mu(M) =0$
\end{nprop}

\begin{ejemplo}
	Un grupo abeliano $M$ está dado por:
	\[
	M: \langle \ x,y,z :  12x+26y+13z = 0  \ , \ 6x+12y+6z = 0 \ , \ 6x+26y+13z = 0 \ \rangle
	\]
	¿Cuál es la descomposición cíclica de este grupo?
\end{ejemplo}
\begin{proof}[Solución]
	Lo primero que habría que hacer es construir la matriz de relaciones.

	\[
	\begin{pmatrix}
 12 & 26 & 13 \\
 6 & 12 & 6 \\
 6 & 26 & 13
\end{pmatrix}
	\]

	Y ahora buscamos su forma normal de Schmidt, para ver cuál es el rango de esta matriz.

	\[
	\begin{pmatrix}
 12 & 26 & 13 \\
 6 & 12 & 6 \\
 6 & 26 & 13
\end{pmatrix} \sim
\begin{pmatrix}
 6 & 12 & 6 \\
 12 & 26 & 13 \\
 6 & 26 & 13
\end{pmatrix} \sim^{C_2 - 2C_1}
\begin{pmatrix}
 6 & 0 & 6 \\
 12 & 2 & 13 \\
 6 & 14& 13
\end{pmatrix}
	\]
	Seguimos haciendo este método a llegar a una de la forma de Schmidt y nos queda la matriz:
	\[
	\begin{pmatrix}
 1 & 0 & 0 \\
 0 & 6 & 0 \\
 0 & 0 & 0
\end{pmatrix}
	\]
	(comprobar si está bien).

	Así, los factores invariantes de la matriz son 1 y 6, y los factores invariantes de M serán , de los obtenidos de la matriz, los que no sean unidades, por lo que en este caso es el 6. Cuidado, porque si aparecieran dos 6, los factores invariantes serían $\{6,6\}$.

	El rango de M, es el número de columnas nulas de esta matriz, que es 1.

	Por tanto, la descomposición cíclica es:
	\[
	M \cong \ent_6 \times \ent
	\]
\end{proof}

\begin{ejemplo}
	Sea
	\[
	M: \langle \ x,y,z :  2x+y+3z = 0  \ , \ 3x+6y-2z = 0 \ , \ 2x+4y+4z = 0 \ \rangle
	\]
	¿Cuál es la descomposición cíclica de este grupo?
\end{ejemplo}
\begin{proof}[Solución]:

De nuevo, hallamos su matriz y le hacemos su forma normal de Schmidt

\[
	\begin{pmatrix}
 2& 1  & 3 \\
 3 & 6 & -2 \\
 2 & 4 & 4
\end{pmatrix}  \sim 	\begin{pmatrix}
 1 & 0 & 0 \\
 0 & 1 & 0 \\
 0 & 0 & 48
\end{pmatrix}
\]
Y hacemos el mismo procedimiento de antes para hallar sus invariantes (solo el 48), su rango (0) y su descomposición cíclica
\[
M \cong \ent_{48}
\]


\end{proof}

\begin{nprop}
	Si $(m,n) = 1 \implies \ent_{mn} \cong \ent_m \times \ent_n$.
Si $d = p_1^{e_1}...p_r^{e_r}$ con $p_i$ primo distintos entre sí y $e_i \geq 1$, entonces:
\[
\ent_d \cong \ent_{p_1^{e_1}} \times ... \times \ent_{p_r^{e_r}}
\]
\end{nprop}

\begin{nprop}
	Si
	\[
	M \cong \ent_{d_1} \times ... \times \ent_{d_s} \times \ent^r
	\]
	con $d_i \geq 2 $ y $d_i / d_{i+1}$, entonces, la descomposición de invariantes en primos estará dada por los mismos primos pero si $d_i = p_1^{e_{i1}}...p_r{e_{ir}}$ y tenedremos $0 \leq e_{ij} \leq e_{i+1j}$

	Así, podemos tener:
	\[
	M \cong \prod_i^s \prod_j ^k \ent_{p_i^{e_{ij}}} \times \ent^r
	\]
	A la que llamamos descomposición cícica Primaria. En ella, la lista ${p_i^{e_{ij}}}$ es la lista de divisores elementales de $M$.
\end{nprop}
\begin{ejemplo}
	Un grupo abeliano tiene como divisores elementales: $\{3,9,5,25,125 y 7\} $ (que son potencias de primos) y rango 0. Así, ese grupo es, en su descomposición primaria:
	\[
	\ent_3 \times \ent_9 \times \ent_5 \times \ent_{25} \times \ent_{125} \ent_7
	\]

Sus invariantes serían $d_1 = 9*125*7 = 7875$, $d_2 = 3*5^2 = 75$, $d_3 = 5$ y ahora los ordenamos, quedando que: $d_3 = 9*125*7 = 7875$, $d_2 = 3*5^2 = 75$, $d_3 = 5$, por lo que su descomposición cíclica sería:
\[
\ent_5 \times \ent_{75} \times \ent_{7875}
\]
\end{ejemplo}

\end{document}


\newpage

\part{Ejercicios}
%%%%%%%%%%%%%%%%%%%%%%%%%%%%%%%%%%%%%%%%%%%%%%%%%%%%%%%%%%%%%%%%
%
% Ejercicios de la asignatura Modelos Matemáticos I.
% Doble Grado de Informática y Matemáticas.
% Universidad de Granada.
% Curso 2016/17.
% 
% 
%
% Agradecimientos:
% Andrés Herrera (@andreshp) y Mario Román (@M42) por
% las plantillas base.
%
% Sitio original:
% https://github.com/libreim/apuntesDGIIM/
%
% Licencia:
% CC BY-NC-SA 4.0 (https://creativecommons.org/licenses/by-nc-sa/4.0/)
%
%%%%%%%%%%%%%%%%%%%%%%%%%%%%%%%%%%%%%%%%%%%%%%%%%%%%%%%%%%%%%%%


%------------------------------------------------------------------------------
%   ACKNOWLEDGMENTS
%------------------------------------------------------------------------------

%%%%%%%%%%%%%%%%%%%%%%%%%%%%%%%%%%%%%%%%%%%%%%%%%%%%%%%%%%%%%%%%%%%%%%%%
% Plantilla básica de Latex en Español.
%
% Autor: Andrés Herrera Poyatos (https://github.com/andreshp) 
%
% Es una plantilla básica para redactar documentos. Utiliza el paquete  fancyhdr para darle un
% estilo moderno pero serio.
%
% La plantilla se encuentra adaptada al español.
%
%%%%%%%%%%%%%%%%%%%%%%%%%%%%%%%%%%%%%%%%%%%%%%%%%%%%%%%%%%%%%%%%%%%%%%%%%

%%%
% Plantilla de Trabajo
% Modificación de una plantilla de Latex de Frits Wenneker para adaptarla 
% al castellano y a las necesidades de escribir informática y matemáticas.
%
% Editada por: Mario Román
%
% License:
% CC BY-NC-SA 3.0 (http://creativecommons.org/licenses/by-nc-sa/3.0/)
%%%

%%%%%%%%%%%%%%%%%%%%%%%%%%%%%%%%%%%%%%%%
% Short Sectioned Assignment
% LaTeX Template
% Version 1.0 (5/5/12)
%
% This template has been downloaded from:
% http://www.LaTeXTemplates.com
%
% Original author:
% Frits Wenneker (http://www.howtotex.com)
%
% License:
% CC BY-NC-SA 3.0 (http://creativecommons.org/licenses/by-nc-sa/3.0/)
%
%%%%%%%%%%%%%%%%%%%%%%%%%%%%%%%%%%%%%%%%%


% Tipo de documento y opciones.
%\documentclass[11pt, a4paper, twoside]{article} % Usar para imprimir
\documentclass[11pt, a4paper]{article}


%---------------------------------------------------------------------------
%   PAQUETES
%---------------------------------------------------------------------------

% Idioma y codificación para Español.
\usepackage[utf8]{inputenc}
\usepackage[spanish, es-tabla, es-lcroman, es-noquoting]{babel}
\selectlanguage{spanish} 
%\usepackage[T1]{fontenc}

% Fuente utilizada.
\usepackage{courier}    % Fuente Courier.
\usepackage{microtype}  % Mejora la letra final de cara al lector.

% Diseño de página.
\usepackage{fancyhdr}   % Utilizado para hacer cabeceras y pies de página.
\usepackage{titlesec} 	% Utilizado para hacer títulos propios.
\usepackage{lastpage}   % Referencia a la última página.
\usepackage{extramarks} % Marcas extras. Utilizado en pie de página y cabecera.
\usepackage[parfill]{parskip}    % Crea una nueva línea entre párrafos.
\usepackage{geometry}            % Geometría de las páginas.

% Símbolos y matemáticas.
\usepackage{amssymb, amsmath, amsthm, amsfonts, amscd}
\usepackage{upgreek}

\usepackage{mdframed}
\usepackage{xcolor}
%\usepackage{enumerate}


% Otros.
\usepackage{enumitem}   % Listas mejoradas.
\usepackage[hidelinks]{hyperref}
\usepackage{pgfplots}
\usepackage{graphicx}            % Gráficos.

% Dibujos Tikz
\usepackage{pgf,tikz}
\usepackage{mathrsfs}
\usetikzlibrary{arrows,calc,intersections,through,backgrounds}

% Fuentes personalizadas

\usepackage[scaled=.85]{newpxtext,newpxmath}
\usepackage[scaled=.85]{FiraSans}
\usepackage[T1]{fontenc}

%   OPCIONES PERSONALIZADAS
%---------------------------------------------------------------------------

\DeclareMathVersion{sans}
\SetSymbolFont{operators}{sans}{OT1}{cmbr}{m}{n}
\SetSymbolFont{letters}{sans}{OML}{cmbr}{m}{it}
\SetSymbolFont{symbols}{sans}{OMS}{cmbrs}{m}{n}
\SetMathAlphabet{\mathit}{sans}{OT1}{cmbr}{m}{sl}
\SetMathAlphabet{\mathbf}{sans}{OT1}{cmbr}{bx}{n}
\SetMathAlphabet{\mathtt}{sans}{OT1}{cmtl}{m}{n}
\SetSymbolFont{largesymbols}{sans}{OMX}{iwona}{m}{n}

\DeclareMathVersion{boldsans}
\SetSymbolFont{operators}{boldsans}{OT1}{cmbr}{b}{n}
\SetSymbolFont{letters}{boldsans}{OML}{cmbrm}{b}{it}
\SetSymbolFont{symbols}{boldsans}{OMS}{cmbrs}{b}{n}
\SetMathAlphabet{\mathit}{boldsans}{OT1}{cmbr}{b}{sl}
\SetMathAlphabet{\mathbf}{boldsans}{OT1}{cmbr}{bx}{n}
\SetMathAlphabet{\mathtt}{boldsans}{OT1}{cmtl}{b}{n}
\SetSymbolFont{largesymbols}{boldsans}{OMX}{iwona}{bx}{n}

\newif\IfInSansMode
\let\oldsf\sffamily
\renewcommand*{\sffamily}{\oldsf\mathversion{sans}\InSansModetrue}
\let\oldmd\mdseries
\renewcommand*{\mdseries}{\oldmd\IfInSansMode\mathversion{sans}\fi\relax}
\let\oldbf\bfseries
\renewcommand*{\bfseries}{\oldbf\IfInSansMode\mathversion{boldsans}\else%
\mathversion{bold}\fi\relax}
\let\oldnorm\normalfont
\renewcommand*{\normalfont}{\oldnorm\InSansModefalse\mathversion{normal}}
\let\oldrm\rmfamily
\renewcommand*{\rmfamily}{\oldrm\InSansModefalse\mathversion{normal}}

% Colores

\definecolor{50}{HTML}{E3F2FD}
\definecolor{300}{HTML}{64B5F6}
\definecolor{500}{HTML}{2196F3}
\definecolor{700}{HTML}{1976D2}
\definecolor{900}{HTML}{0D47A1}

% ---------------------------------------------------------------------------
% OPCIONES PERSONALIZADAS
% ---------------------------------------------------------------------------

% Redefinir letra griega épsilon.
\let\epsilon\upvarepsilon

% Formato de texto.
\linespread{1.3}            % Espaciado entre líneas.
\setlength\parindent{0pt}   % No indentar el texto por defecto.
\setlist{leftmargin=.5in}   % Indentación para las listas.

% Estilo de página.
\pagestyle{fancy}
\fancyhf{}
\geometry{left=3cm,right=3cm,top=3cm,bottom=3cm}   % Márgenes y cabecera.

% Estilo de las cabeceras

\titleformat{\section}
{\Large\bfseries\sffamily}{\thesection}{1em}{}
\titleformat{\subsection}
{\large\sffamily}{\thesubsection}{1em}{}

% Estilo de enlaces
\hypersetup{
    % hidelinks = true,   % Oculta todos los enlaces.
    colorlinks = true,   % Muestra todos los enlaces, sin bordes alrededor.
    linkcolor=black,     % Color de enlaces genéricos
    citecolor={blue!50!black},   % Color de enlaces de referencias
    urlcolor={blue!80!black}     % Color de enlaces de URL
}

% Ruta donde buscar gráficos
\graphicspath{{../Recursos/Plantillas/}, {Recursos/Plantillas/}}

% Redefinir entorno de demostración (reducir espacio superior)
\makeatletter
\renewenvironment{proof}[1][\proofname] {\vspace{-15pt}\par\pushQED{\qed}\normalfont\topsep6\p@\@plus6\p@\relax\trivlist\item[\hskip\labelsep\it#1\@addpunct{.}]\ignorespaces}{\popQED\endtrivlist\@endpefalse}
\makeatother

% Aumentar el tamaño del interlineado
\linespread{1.3}

% Numerar las ecuaciones del entorno {equation} con respecto a las secciones
\numberwithin{equation}{section}

%---------------------------------------------------------------------------
%   COMANDOS PERSONALIZADOS
%---------------------------------------------------------------------------

% Valor absoluto: \abs{}
\providecommand{\abs}[1]{\lvert#1\rvert}    

% Fracción grande: \ddfrac{}{}
\newcommand\ddfrac[2]{\frac{\displaystyle #1}{\displaystyle #2}}

% Texto en negrita en modo matemática: \bm{}
\newcommand{\bm}[1]{\boldsymbol{#1}}

% Línea horizontal.
\newcommand{\horrule}[1]{\rule{\linewidth}{#1}}

% Renombrar la R de los reales

\newcommand{\R}{\mathbb{R}}



%---------------------------------------------------------------------------
%   CABECERA Y PIE DE PÁGINA
%---------------------------------------------------------------------------


\renewcommand{\sectionmark}[1]{%
\markboth{#1}{}}

\renewcommand{\subsectionmark}[1]{%
\markright{#1}{}}

%\addtolength{\headheight}{4ex}
\renewcommand{\headrulewidth}{0pt}

\addtolength{\headwidth}{\marginparsep}
%\addtolength{\headwidth}{\marginparwidth}

\fancypagestyle{section}{%
    \fancyhead{}
    %\addtolength{\headheight}{-10ex}
    %\renewcommand{\headheight}{0pt}%
    %\setlength{\footskip}{-48pt}%
    \fancyfoot[LE,RO]{\Large\sffamily\thepage}
    \renewcommand{\headrulewidth}{0pt}%
    \renewcommand{\footrulewidth}{0pt}%
}

\let\originalsection\section
\RenewDocumentCommand{\section}{som}{%
    \IfBooleanTF{#1}
    {\originalsection*{#3}}
    {\IfNoValueTF{#2}
        {\originalsection{#3}}
        {\originalsection[#2]{#3}}%
    }%
    \thispagestyle{section}%
}

\fancyhead[LE,RO]{\rule[-4ex]{0pt}{2ex}\sffamily\textsl{\rightmark}}
\fancyhead[LO,RE]{\sffamily{\leftmark}}
\fancyfoot[LE,RO]{\Large\sffamily\thepage}

%---------------------------------------------------------------------------
%   ENTORNOS PARA MATEMÁTICAS
%---------------------------------------------------------------------------

% Nuevo estilo para definiciones.
\newtheoremstyle{definition-style} % Nombre del estilo.
{}               % Espacio por encima.
{}               % Espacio por debajo.
{}                   % Fuente del cuerpo.
{}                   % Identación.
{\bf\sffamily}                % Fuente para la cabecera.
{.}                  % Puntuación tras la cabecera.
{.5em}               % Espacio tras la cabecera.
{\thmname{#1}\thmnote{ #3}}     % Especificación de la cabecera (actual: nombre en negrita).

% Nuevo estilo para notas.
\newtheoremstyle{remark-style} 
{10pt}                
{10pt}                
{}                   
{}                   
{\itshape \sffamily}          
{.}                  
{.5em}               
{}                  

% Nuevo estilo para teoremas y proposiciones.
\newtheoremstyle{theorem-style}
{}                
{}                
{}           
{}                  
{\bfseries \sffamily}             
{.}                
{.5em}           
{\thmname{#1}\thmnumber{ #2}\thmnote{ (#3)}}


% Configuración general de mdframe, los estilos de los teoremas, etc
\mdfsetup{skipabove=12pt,skipbelow=12pt,innertopmargin=12pt,innerbottommargin=4pt}

% Creamos los 'marcos' de los estilos

\mdfdefinestyle{nth-frame}{
    linewidth=2pt, %
    linecolor= 500, % 
    %linecolor=black,
    topline=false, %
    bottomline=false, %
    rightline=false,%
    leftmargin=0pt, %
    innerleftmargin=1em, % 
    rightmargin=0pt, %
    innerrightmargin=0pt, % % 
    %splittopskip=\topskip, %
    %splitbottomskip=\topskip, %
}% 

\mdfdefinestyle{nprop-frame}{
    linewidth=2pt, %
    linecolor= 300, %
    %linecolor= gray, 
    topline=false, %
    bottomline=false, %
    rightline=false,%
    leftmargin=0pt, %
    innerleftmargin=1em, %
    innerrightmargin=0em, 
    rightmargin=0pt, %
    %splittopskip=\topskip, %
    %splitbottomskip=\topskip, %
}%       

\mdfdefinestyle{ndef-frame}{
    linewidth=2pt, %
    linecolor= 300, % 
    %linecolor= gray!50,
    backgroundcolor= 50,
    %backgroundcolor= gray!5,
    topline=false, %
    bottomline=false, %
    rightline=false,%
    leftmargin=0pt, %
    innerleftmargin=1em, %
    innerrightmargin=1em, 
    rightmargin=0pt, % 
    innertopmargin=1.5em,%
    innerbottommargin=1em, % 
    splittopskip=\topskip, %
    %splitbottomskip=\topskip, %
}%               

% Asignamos los marcos a los estilos

\surroundwithmdframed[style=nth-frame]{nth}
\surroundwithmdframed[style=nprop-frame]{nprop}
\surroundwithmdframed[style=nprop-frame]{ncor}
\surroundwithmdframed[style=ndef-frame]{ndef}
\surroundwithmdframed[style=ndef-frame]{ejer}

% Nuevo estilo para ejemplos.
\newtheoremstyle{example-style}
{10pt}                
{10pt}                
{}                  
{}                   
{\bf \sffamily}              
{:}                 
{.5em}               
{}                   

% Teoremas, proposiciones y corolarios.
\theoremstyle{theorem-style}
\newtheorem{nth}{Teorema}[section]
\newtheorem{nprop}{Proposición}[section]
\newtheorem{ncor}{Corolario}[section]
\newtheorem{lema}{Lema}[section]
% Definiciones.
\theoremstyle{definition-style}
\newtheorem{ndef}{Definición}[section]
\newtheorem{ejer}{Ejercicio}[section]

% Notas.
\theoremstyle{remark-style}
\newtheorem*{nota}{Nota}
\newtheorem*{sol}{Solución}

% Ejemplos.
\theoremstyle{example-style}
\newtheorem{ejemplo}{Ejemplo}[section]

% Listas ordenadas con números romanos (i), (ii), etc.
\newenvironment{nlist}
{\begin{enumerate}
    \renewcommand\labelenumi{(\emph{\roman{enumi})}}}
{\end{enumerate}}

% División por casos con llave a la derecha.
\newenvironment{rcases}
{\left.\begin{aligned}}
    {\end{aligned}\right\rbrace}

    %---------------------------------------------------------------------------
    %   PÁGINA DE TÍTULO
    %---------------------------------------------------------------------------

    % Título del documento.
    \newcommand{\subject}{Modelos Matemáticos I}

    % Autor del documento.
    \newcommand{\docauthor}{Doble Grado en Ingeniería Informática y Matemáticas}

    % Título
    \title{
        \normalfont \normalsize 
        \textsc{Universidad de Granada} \\ [25pt]    % Texto por encima.
        \horrule{0.5pt} \\[0.4cm] % Línea horizontal fina.
        \huge \subject\\ % Título.
        \horrule{2pt} \\[0.5cm] % Línea horizontal gruesa.
    }

    % Autor.
    \author{\Large\sffamily{\docauthor}}

    % Fecha.
    \date{\vspace{-1.5em} \normalsize \sffamily Curso 2016/17}


    %---------------------------------------------------------------------------
    %   COMIENZO DEL DOCUMENTO
    %---------------------------------------------------------------------------
    \begin{document}

    \maketitle  % Título.
    \tableofcontents    % Índice
    \vfill
    \begin{center}
        \includegraphics{by-nc-sa}  % Licencia.
    \end{center}
    \newpage


    %----------------------------
    %   Introducción.
    %----------------------------

    \section{Relación 1}
    \subsection{Ejercicio 2}

    Sea la ecuación:
    \[
        x _{n+1} =  1/3 x_n +2^n
    \]
    Hallar una solución del tipo $x_n=c2^n$. 

    Llo primero es ver que esta ecuación es una ecuación en diferencias de primer orden lineal no homogénea. Vamos a hallar el $c$ que nos piden. Entonces:
    \[
        c2^{n+1}= \dfrac{1}{3}c2^n+2^n  \implies 2^n[2c-c\dfrac{1}{3}-1] = 0
    \]
    \[
        2c-\dfrac{1}{3}c-1 = 0 \implies c = \dfrac{3}{5}
    \]
    Por tanto, la solución es: $x_n' =\dfrac{3}{5}2^n $, pero esta es una solución particular.

    Ahora, probaremos que $x_n$ es solución de la ecuación inicial $\iff z_n = x_n - x_n'$ donde $x_n'$ es una solución particular, es solución de:
    $z _{n+1} = \dfrac{1}{3} z_n$.

    Lo probamos:
    Que $x_n$ es solución de la inicial implica que $x _{n+1} =  \dfrac{1}{3} x_n +2^n$.
    Que $x_n'$ es solución de la inicial implica que $x _{n+1}' = \dfrac{1}{3}x_n' 2^n$.

    Si restamos estas dos obtenemos que:
    \[
        x _{n+1} - x _{n+1}' = \dfrac{1}{3}(x_n- x_n')
    \]
    Lo que implica que $z _{n+1} = x _{n+1} - x _{n+1}'$ y $z_n =(x_n- x_n') $

    Ahora, continuamos la resolución del ejercicio:
    \begin{enumerate}
        \item Encontrar una solución particular de la ecuación inicial, $x_n' = \dfrac{3}{5}2^n$
        \item Resolvemos la ecuación homogénea asociada
            \[
                z _{n+1} = \dfrac{1}{3}z_n  \implies z_n = K(\dfrac{1}{3})^n
            \]
        \item Usando el resultado:
            \[
                x_n = x_n' + z_n = \dfrac{3}{5}2^n+K(\dfrac{1}{3})^n
            \]
        \item Aplicamos la condición inicial $x_0 = 1$.


    \end{enumerate}
    Así, la solución es $x_n = \dfrac{3}{5}2^n +\dfrac{2}{5}(\dfrac{1}{3})^n$ con $n\geq 0$

    \subsection{Ejercicio 3}
    Primero calculamos $KerL$.
    $$ x \in KerL \Leftrightarrow L(x) = 0$$
    $$x_n^{*} = 0 \Leftrightarrow x_{n+1} - (2 + i ) x_n = 0 \Leftrightarrow x_{n+1} = (2+i)x_n$$
    Solución general: $$ x_n = K (2+i)^{n} \ , \ K \in C$$
    Por lo tanto los elementos de $X$ que pertenecen al $KerL$ son de la forma que acabamos de indicar. Veamos ahora los elementos de $X$ tales que $L(x) = b$ siendo b = \{ 1,1,1,1.. \}.
    $$x_{n+1} -(2+i)x_n = 1 \Leftrightarrow x_{n+1} = (2+i)x_n + 1$$

    \[
        \begin{rcases}
            x _{n+1} = (2+i) x_n + 1 \\
            x _{n+1}^* = (2+i) x_n^* +  1 
        \end{rcases}
    \]
    Restamos y tenemos que:
    $$ x _{n+1}  - x _{n+1}^* = (2 + i) (x_n - x_n^*)$$
    LLamamos $y_n = x_n^* - x_n$, y tenemos que $ y_{n+1} = (2+i) y_n $. Por lo tanto la solucion general de $y_n$ es: $y_n = K (2+i)^n \ , \ k \in C $.
    Ahora calculamos la solucion constante $x^*$.
    $$ x^* = (2+i)x^* +1  \Leftrightarrow x^* (1-2-i) = 1 \Leftrightarrow x^* = \frac{1}{-1-i} \Leftrightarrow x^* = \frac{1(-1+i)}{(-1-i)(-1+i)} \Leftrightarrow x^* = \frac{-1+i}{2} $$
    Y ya podemos calcular la solución general de $x_n$ que es:
    $$ x_n = y_n + x^* = K (2+i)^n + \frac{-1+i}{2}$$
    Es decir los elementos de $X$ que cumplen que $L(x) = b$ son de esa forma.

    \subsection{Ejercicio 4}
    Llamaremos primero $x_n$ el número de clientes de $Paga^+$ en el año n.
    Llamaremos $y_n$ el número de clientes de $Paga^-$ en el año n.
    Ahora, tenemos que:
    \[
        \begin{rcases}
            x _{n+1} =0.5 x_n + 0.25 y_n \\
            y _{n+1} = 0.75 y_n +  0.5 x_n 
        \end{rcases}
    \]
    Sabemos que $x_n+y_n = 1$, que es el $100\%$ de los clientes. Podemos despejar así una en función de la otra y nos queda:
    \[
        x _{n+1} = 0.5 x_n + 0.25(1-x_n)
    \]
    Esta es una ecuación en diferencias de primer orden lineal no homogénea. Despejando, tenemos:
    \[
        x _{n+1} = 0.25 x_n +0.25 \implies x_n =  x^* + z_n \implies x_n =  \dfrac{0.25}{1-0.25} + K(0.25)^n
    \]
    Ahora, tenemos que tomar el límite pues nos piden el mercado a largo plazo. Esto es:
    \[
        \lim_{n\to \infty}x_n = \dfrac{0.25}{0.75} = \dfrac{1}{3}
    \]
    Podemos afirmar ahora que, asintóticamente , $\dfrac{1}{3}$ de los clientes estarán en $Paga^+$ y $\dfrac{2}{3}$ en $Paga^-$

    \subsection{Ejercicio 6}
    Vamos a llamar $A_n$ al número de empleados en el departamento $A$ en el año $n$. Del mismo modo, llamaremos $B_n$ a los del departamento $B$ en el año $n$. Si $0 \leq p \leq 1$  y $0 \leq q \leq 1$, tenemos que:
    \[
        \begin{rcases}
            A _{n+1} = (1-p)A_n + qB_n\\
            B _{n+1}= pA_n + (1-q)B_n
        \end{rcases}
    \]
    Y tenemos ahora que $A_n + B_n = M$, (aunque se podría tomar como 1) por lo que el sistema lo reducimos a una ecuación teniendo:
    \[
        A _{n+1} = (1-p)A_n +q(M-A_n)\implies A _{n+1} =(1-p-1)A_n +qM
    \]
    Teniendo de nuevo una ecuación en diferencias de primer orden lineal no homogénea. Así, su solución es:
    \[
        A_n = \dfrac{qM}{1-(1-p-q)} + K(1-p-q)^n
    \]
    Estudiamos primero la solución constante.
    \[
        A_* = \dfrac{qM}{1-(1-p-q)}= \dfrac{qM}{p+q}
    \]
    Impondremos que $p+q\ne 0$, al menos 1 trabajador cambiará de departamento.
    Nuestro término general es:
    \[
        A_n = \dfrac{qM}{p+q} + K(1-p-q)^n
    \]
    Estudiaremos esto a largo plazo. Tenemos que $-1 \leq 1-p-q < 1$. Tenemos que $|1-p-q| < 1$, por lo que:
    \[ \begin{cases}
            \lim_{n \to \infty} A_n = \dfrac{qM}{p+q} = M \dfrac{q}{p+q}\\
            \lim_{n \to \infty}B_n = \lim_{n\to \infty}(M-A_n) = M- M \dfrac{q}{p+q} = \dfrac{p}{p+q}M
        \end{cases}
    \]

    \subsection{Ejercicio 8}
    \begin{enumerate}
        \item $X_n$ es el número de árboles en el año n.
            \[
                x _{n+1} = 0.9 x_n +K
            \]
            Esto tiende a la solución constante $x_* = \dfrac{K}{0.1}$
    \end{enumerate}

    \subsection{Ejercicio 12}
    La ecuación logística de Pielou es una ecuación en diferencias no lineal de la forma:
    \[
        x_{n+1}= \frac{\alpha x_n}{1+ \beta x_n} \quad \quad \alpha > 1 \quad \beta > 0
    \]

    \begin{enumerate}
        \item Demuestre que posee un punto de equilibrio positivo.
            Tenemos que buscar un $\bar{\alpha}\in \R$ tal que $\bar{\alpha} = f(\bar{\alpha})$ donde  $f(x)= \frac{a x}{1+bx}$. Por tanto, buscamos la solución de:
            \[
                \alpha = \frac{a x}{1+bx} \implies b\alpha^2 - a\alpha + \alpha = 0 \implies \alpha[b\alpha +(1-a)]=0 \implies \alpha = 0
            \]
            Es un punto de equilibrio, así que tenemos que:
            \[
                b\alpha +(1-a) = 0 \implies \alpha =  \frac{a-1}{b} > 0
            \]

        \item Tomando $\alpha = 2$ y $\beta = 1$, probar que el punto de equilibrio es Asintóticamente Estable. Para ello , trabajaremos ahora en $\R^+$.
            Tenemos que $f(x) = \dfrac{2x}{1+x} \in \mathcal{C}^\infty$ y que su derivada es:
            \[
                f'(x) = \frac{2(1+x)-2x}{(1+x)^2} = \frac{2}{(1+x)^2}\implies |f'(1)|= \frac{1}{2} < 1
            \]
            Por lo que sí es A.E.

        \item Demostrar que el cambio de variable $x_n = \dfrac{1}{z_n}$ trasnforma esa ecuación en una ecuación lineal del primer orden.
            Tenemos ahora por tanto:
            \[
                \frac{1}{z_{n+1}} =  \frac{a \dfrac{1}{z_n}}{1+b \dfrac{1}{z_n}} \implies z_{n+1}= \dfrac{1}{a}z_n + \dfrac{b}{a} \quad (2)
            \]
            Que es una ecuación lineal de primer orden

        \item A partir del resultado anterior, determinar el comportamiento asintótico de las soluciones de la ecuación logística.

            Esta es una ecuación en diferencias lineal no homogénea, por tanto tendremos una solución $z_n = z_* +y_n$ donde $z_*$ es la solución constante y $y_n$ es la solución de la homogénea asociada.

            Tenemos que $z_*=\frac{b/a}{1- (1/a)} = \frac{b}{a-1}> 0$. Ahora, calculamos la ecuación homogénea asociada a $(2)$, calculado en el apartado anterior:
            \[
                y_n = \mathcal{C} + (\dfrac{1}{a})^n
            \]
            Ya tenemos la solución, por tanto la solucion de $z_n$ es:
            \[
                z_n =  \dfrac{b}{a-1} + \mathcal{C}(\dfrac{1}{a})^n \quad n \geq 0
            \]
            El comportamiento asintótico de esta es convergente a $\dfrac{b}{a-1}$ pues , como $a > 1$, entonces $(\frac{1}{a})^n \to 0$

            Ahora, volviendo al cambio de variable, como $z_n$ tiene límite distinto de cero, entonces:
            \[
                \lim_{n\to \infty} x_n = \frac{1}{\lim_{n\to \infty}z_n} = \frac{1}{b/(a-1)}= \frac{a-1}{b}
            \]
            Que es el punto de equilibrio de $x_n$, luego $\alpha= \frac{a-1}{b}$ es Asintóticamente Estable.

    \end{enumerate}

    \subsection{Ejercicio 15}

    \begin{enumerate}
        \item Hecho en clase
        \item Para conseguir un equilibrio poblacional asintóticamente estable (a.e.), se propone vender una fracción $\alpha$$(0 < \alpha < 1)$ de la población en cada periodo de tiempo dando lugar al modelo:
            \[
                p_{n+1} = 10(1 - \alpha)p_n e^{(1-\alpha)p_n}
            \]

            i) Encuentre el intervalo abierto (de amplitud máxima) donde elegir $\alpha$ para que esté asegurada la estabilidad
            asintótica del equilibrio positivo.


            Si :
            \[
                \mu = \frac{ln(10(1-\alpha)}{(1-\alpha)}
            \]
            Vamos a comprobar que es L.A.E.
            \[
                f'(x) = 10(1-\alpha)e^{-(1-\alpha)x}(1+x(-(1-\alpha))
            \]
            \[
                f'(\mu) 0 1-ln(10(1-\alpha)) \quad y \quad 0 < ln(10(1-\alpha))< 2
            \]
            Y ahora, si:
            \begin{itemize}
                \item $\alpha < 9/10$
                \item $\alpha > \frac{10-e^2}{10}$
            \end{itemize}
            Por tanto $\mu$ es L.A.E si $\frac{10-e^2}{10} < \alpha < 9/10$

            ii) Calcule el valor de $\alpha$ para el que la población de equilibrio alcanza su valor máximo y es a.e.

            Ahora, si $g(x) = \frac{ln(10(1-x))}{1-x}\implies g'(x)= \frac{-1+ln(10(1-x))}{(1-x)^2}$.

            Ahora, $g'(x) = 0\implies 1-ln(10(1-x)) = 0 \implies x= \frac{10-e}{10}$ es donde la poblaciñon alcanza su máximo.

            Ahora, si comprobamos: $x_0 > \mu \implies x_1 = 10(1-\alpha)x_0 e^{-(1-\alpha)x_0} = x_0 \frac{10(1-\alpha)}{e^{(1-\alpha)x_0}} < x_0 \implies x _{n+1} < x_n$. Y Si miramos ahora si $x_0 < \mu$, veremos que $x _{n+1} > x_n$
    \end{enumerate}

    \subsection{Ejercicio 17}
    En cierto mercado, los precios de determinado producto siguen una din\'amica basada en los postulados del modelo de la telaraña pero se ha observado que las funciones de oferta y demanda vienen dadas por:

    \[ 
        O(p) = 1 + p^{2}, \ D(p) = c - dp
    \]\

    donde c \textgreater 1 y d \textgreater 0. Suponemos, adem\'as, que el equilibrio del mercado se da si la Oferta iguala a la Demanda y que la oferta en el periodo (n+1)-\'esimo depende del precio del periodo n-\'esimo.

    a) Deduzca la ED, $p_{n+1} = F(p_n)$, que describe la din\'amica planteada y calcule el precio de mercado $p^{*}$ (punto de equilibrio econ\'omicamente factible).

    b) Deduzca las condiciones sobre c y d que aseguran la estabilidad asit\'otica de $p^{*}$. ¿Qu\'e ocurre si d = 2 y c = 4 ?

    c) Para c = 3 y d = 2, use el diagrama de Cobweb para trazar los valores $p_1$ y $p_2$ a partir de $p_0$ = 1. A largo plazo, ¿c\'omo se comportar\'an los precios en el caso contemplado?

    a)

    Deducimos la ED, sabemos que el equilibrio del mercado se da si la oferta es igual a la demanda por lo que :

    \[
        O(p_{n-1}) = D (p_n) \Rightarrow
    \]

    \[ 
        \Rightarrow 1 + p_{n-1}^{2} = c - dp_n \Rightarrow 
    \]

    \[ 
        \Rightarrow p_n = \frac{c-1}{d} - \frac{p_{n-1}^{2}}{d} \Rightarrow 
    \]

    puesto que $d>0$

    \[
        \Rightarrow p_{n+1} = \frac{c-1}{d} - \frac{p_{n}^{2}}{d} = F(p_n)
    \]

    Para calcular el punto de equilibrio ($\alpha$) tenemos que resolver la ecuaci\'on 

    \[
        \alpha = F(\alpha) \Rightarrow
    \]

    \[
        \Rightarrow F(\alpha) = \frac{c-1}{d} - \frac{\alpha^{2}}{d} = \alpha \Rightarrow
    \]

    \[
        \Rightarrow \alpha^{2} + d\alpha -c+1 = 0 \Rightarrow
    \]

    \[
        \Rightarrow \alpha = \frac{-d\pm\sqrt{d^{2} - 4(-c+1)}}{2} = \frac{-d\pm\sqrt{d^{2} +4c-4}}{2}
    \]

    Dado que $\alpha$ es el precio de equilibrio y \'este es un n\'umero real entonces se tiene que dar que el discriminante sea mayor que 0 lo cual es cierto ya que c$>$1 y d$>$0. Adem\'as se tiene que el precio de equilibrio es un n\'umero positivo así que debemos escoger aquellos $\alpha$s que sean positivos. As\'i se verifica que :

    \[
        d^{2} +4c-4 > d^{2} \Rightarrow 
    \]

    \[
        -d + \sqrt{d^2 +4(c-1)} > -d + \sqrt{d^{2}} = 0
    \]

    Se puede ver que el precio de equilibrio que buscamos es:

    \[
        \alpha = \frac{-d+\sqrt{d^{2} +4c-4}}{2}
    \]

    b) 
    Recordemos que para que un punto de equilibrio sea localmente asintóticamente estable (L.A.E.) se tienen que dar dos condiciones: la primera es que la funci\'on sea de clase 1 y segundo que el valor absoluto de la derivada evaluada en el punto de equilibrio sea menor que 1. Como podemos ver $F \in C^{\infty}$ luego, en particular, $F \in C^{1}$, basta con ver para que valores de c y d se tiene que la mencionada derivada sea menor que 1.

    \[
        F(x) = \frac{c-1}{d} - \frac{1}{d}x^{2} \Rightarrow
    \]

    \[
        \Rightarrow |F'(x)| = |\frac{-2}{d}x| 
    \]

    \[
        \alpha = \frac{-d+\sqrt{d^{2} +4c-4}}{2} \Rightarrow
    \]

    \[
        \Rightarrow |F'(\alpha)| = |\frac{d-\sqrt{d^{2}+4c-4}}{d}| = | 1 - \frac {\sqrt{d^{2}+4c-4}}{d}| < 1
    \]

    \[
        \Rightarrow -1 <  1 - \frac {\sqrt{d^{2}+4c-4}}{d} < 1 \Rightarrow
    \]

    \[
        \Rightarrow -2 < - \frac {\sqrt{d^{2}+4c-4}}{d} < 0 \Rightarrow
    \]

    \[
        \Rightarrow 2 > \frac {\sqrt{d^{2}+4c-4}}{d} > 0 \Rightarrow
    \]

    \[
        \Rightarrow 2d > \sqrt{d^{2}+4c-4} > 0 \Rightarrow
    \]

    \[
        \Rightarrow 4d^{2} > d^{2}+4c-4 \Rightarrow 
    \]

    \[
        \Rightarrow 3d^{2}-4c+4 > 0
    \]

    Si d = 2 y c = 4

    \[
        \Rightarrow 3*4 - 4*4 +4 = 0 
    \]

    Luego la condici\'on anterior no nos garantiza que ese punto sea L.A.E.
    \

    \

    \

    Sustituyendo en $F(x)$:

    \[
        F(x) = \frac{4-1}{2} - \frac{1}{2}x^{2} = \frac{3}{2} - \frac{1}{2}x^{2}\Rightarrow
    \]

    \[
        |F(x)|' = |-x| = x 
    \]

    Sustituyendo $\alpha$ por x:

    \[
        x = \alpha = \frac{-2+\sqrt{2^{2} +4*4-4}}{2} = 1
    \]

    El resultado como podemos ver es 1, lo cual no nos dice que sea L.A.E. ni inestable. Por tanto tendremos que comprobar que ocurre con la segunda derivada (si es distinta o mayor que 0 entonces es inestable y si es menor que 0 es L.A.E.) :

    \[
        F(x)' = -x \Rightarrow
    \]

    \[
        \Rightarrow F(x)'' = -1 \neq 0
    \]

    \[
        \Rightarrow F(\alpha)'' = -1
    \]

    Por tanto, es inestable.

    c) Si c = 3 , d = 2 y $p_0 = 1$ , entonces usando que : 

    \[
        p_{n+1} = \frac{c-1}{d} - \frac{p_{n}^{2}}{d} = F(p_n) \Rightarrow
    \]

    \[
        \Rightarrow p_{n+1} = 1 - \frac{1}{2}p_n^{2} = F(p_n)
    \]

    \[
        \Rightarrow p_1 = 1 - \frac{p_0^{2}}{2} = \frac{1}{2}
    \]

    \[
        \Rightarrow p_2 = 1 - \frac{p_1^{2}}{2} = \frac{7}{8}
    \]

    %\includegraphics[scale=0.32]{/home/usuario/Escritorio/diagrama(ej17).png}
    \

    Podemos ver que el punto de equilibrio es un atractor y que la funci\'on converge a dicho punto:
    \[
        \alpha = \frac{-d+\sqrt{d^{2} +4c-4}}{2} = \frac{-2+\sqrt{4 +12-4}}{2} = \frac{-2+2\sqrt{3}}{2} = -1+\sqrt{3}  
    \]

    \

    Para realizar este gr\'afico he usado wxmaxima, en concreto, la orden 

    \[
        staircase(1-(1/2)*x^2,1,100,[x,0,1.5],[y,0,1.5]);
    \]

    donde el primer parámetro es la función obtenida antes $(F_n)$ representada por al curva roja, el segundo parámetro es el punto inicial, el tercero es el número de iteraciones o valores que calculamos (n=100) y los últimos parámetros son opciones para representar el gráfico.

    \subsection{Ejercicio 19}

    Determine razonadamente la veracidad o falsedad de las siguientes afirmaciones.
    a) Sea la ecuación en diferencias de primer orden, $x _{n+1} = f (x_n )$, donde $f(x)$ es suficientemente derivable en
    un entorno del punto de equilibrio, s. Si $f'(s) = 1 $y $f'' (s) = 0$, entonces el punto de equilibrio siempre será
    inestable.

    El punto de equilibrio es siempre inestable


    b) Toda ecuación en diferencias $x _{n+1} = f (x_n)$, con un punto de equilibrio inestable en $x = s$ admite, al menos,
    un 2-ciclo o solución 2-periódica.

    Falso. Basta tomar la función $f(x) = ax+b$ tal que si $a>1$, entonces la derivada es mayor que 1 y la función no tiene ningún ciclo.

    c) Todas las soluciones no constantes de la ecuación en diferencias lineal $x_{n+2}+ x_{n-1} = 0$, son 3-periódicas.

    Si despejamos la ecuación, vemos que:
    \[
        x_{n+3} = -x_n
    \]
    Dados 3 datos iniciales, $x_0,x_1,x_2$, podemos escribir algunos términos de la recurrencia en función de los anteriores y escribimos su órbita, quedando:
    \[
        \{x_0,x_1,x_2,-x_0,-x_1,-x_2,x_0,\cdots\}
    \]

    Por lo que el periodo es 6, no 3 y por tanto es FALSO.

    d ) La dinámica de cierta población ovina se rige mediante un modelo logístico discreto. Si la tasa de crecimiento
    para una población de 10 cientos es de $1.5$ y para una de población de 20 cientos es de $1$; entonces, a largo plazo, la población tiende a un valor constante.

    Tenemos que $P_n=10$ y la tasa de $P_n$ es 1.5 . También que si $P_m = 20$, entonces su tasa es de 1.

    Ajustando por un modelo conveniente, tenemos que:
    \[
        tasa(P_n) = \frac{P _{n+1} - P_n}{P_n} = a -bP_n
    \]
    y si despejamos, obtenemos:
    \[
        P _{n+1} = P_n + P_n -b P_n^2 \implies P _{n+1} = (1+a)P_n(1-\frac{b}{1+a}P_n)
    \]
    ASí que si tomamos $\mu = (1+a)$, $x_n = \frac{b}{1+a}P_n$, entonces:
    \[
        x _{n+1} = \mu x_n (1-x_n)
    \]
    Ahora, debemos hallar a y b. planteamos el sistema:
    \[
        \begin{rcases}
            1.5 = a-10b\\
            1 = a-20b
        \end{rcases} 
        \begin{cases}
            a = 2 \\
            b = 0.05
        \end{cases}
    \]
    y planteamos la ecuación:
    \[
        P _{n+1} = 3P_n - 0,05 P_n^2
    \]
    entonces:
    \[
        P_* = 3P_* - 0,05P_*^2
    \]
    Si igualamos $0,05P_*^2 -2P_* = 0 \implies P_* = 0$, pero también tenemos que tener en cuenta el caso $0,05P_* - 2 = 0 \implies P_* = 40$

    Ahora, si $f(x) = 3x - 0,05x^2$, entonces $f'(0) = 3$, por lo que es inestable y $f'(40) =3-4 = -1$ y $f''(40) = -0,05$
    pero como la derivada vale $-1$, no podemos determinar con claridad lo que ocurre en este caso. Sin embargo, si realizamos el dibujo en $wxmaxima$, veríamos que se realiza la tela de araña y así converge.



    %%%%%%%%%%%%%%%%%%%%%%%%%%%%%%%%%%%%%%%%%%%%%%%%%%%%%%%%%%%%%%%%%%%%%%%%%%%%%%%%%%%%%%%%%%%%%%%%%%%%%%%%%%%5

    \section{Relación 2}

    \subsection{Ejercicicio 2}
    Calcule una solución de la ecuación $x_ {n+3} -x_n = 0$ con condiciones iniciales $x_0 = 1$, $x_1 = 0$ y $x_2 = 0$.

    Lo primero es ver que su polinomio característico es:
    \[
        p(\lambda) = \lambda^3 -1
    \]
    y las raíces de este son:
    \[
        \lambda = \frac{-1 \pm \sqrt{-3}i}{2}
    \]
    y por tanto tenemos: $\lambda_1 = 1$, $\lambda_2 = cos(\dfrac{2\pi}{3})+ i sen(\dfrac{2\pi}{3})$ y $\lambda_3 = cos(\dfrac{2\pi}{3}) - i sen(\dfrac{2\pi}{3})$
    Así, la solución general es:
    \[
        x_n = c_1 * 1^n + c_2 cos(n(\dfrac{2\pi}{3})) + c_3 sen(n(\dfrac{2\pi}{3}))
    \]
    Solo faltaría resolver el sistema con las condiciones iniciales para obtener la solución particular del sistema.


    \subsection{Ejercicio 4}
    Dado un número real $\alpha$, determine una expresión para la sucesión $\{x_n\}$ que verifica:
    \[
        \begin{cases}
            x_{k+2} = x _{k+1} + x_k \ \ k \in \mathbb N
            x_1 = 1, \ \ \ x_2 = \alpha
        \end{cases}
    \]
    Su polinomio característico es $p(\lambda) = \lambda^2 - \lambda - 1$, cuyas soluciones son $\lambda_1 = \frac{1+\sqrt 5}{2}$ y $\lambda_2 = \frac{1+ \sqrt 5}{2}$. Con las condiciones iniciales, podemos obtener la expresión de la sucesión hallando $c_1$ y $c_2$.

    Apartado $b$): estudiar el comportamiento de la sucesión $\{y_n\}$ definida por:
    \[
        y_{k+1} = \frac{1}{1+y_k} \ \ \ k \geq 2 \quad y \quad y_2 = 1
    \]
    Hallando su polinomio característico, sus raíces son $y = \frac{-1 \pm \sqrt 5}{2}$

    \subsection{Ejercicio 9}

    La sucesión $\{\overline{x_n}\}_{n\ge0} = \{1, 2, 5, 12, 29, \dots\}$ es solución de cierta ecuación en diferencias lineal homogénea de orden 2 con coeficientes constantes.\\

    \textbf{a)} Deduzca dicha ecuación en diferencias.

    Sabemos que la ecuación en diferencias será de la forma
    \begin{equation} \label {1}
        x_{n+2} + a_1x_{n+1} + a_0x_n = 0, \quad n \ge 0
    \end{equation}

    Si sustuimos en (\ref 1) los valores $n=0$ y $n=1$ para la solución particular $\{\overline{x_n}\}_{n\ge0}$, obtenemos el siguiente sistema de ecuaciones: 

$$\begin{cases} 5 + 2a_1 + a_0 = 0\\ 12 + 5a_1 + 2a_0 = 0 \end{cases}$$

    cuyas soluciones son $a_1 = -2$, $a_0 = -1$. Por tanto, la ecuación en diferencias es:
    \begin{equation} \label{2}
        x_{n+2} - 2x_{n+1} - x_n = 0, \quad n \ge 0
    \end{equation}

    \textbf{b)} Dé una expresión de la solución general de la ecuación en diferencias.

    Resolvemos la ecuación (\ref 2) como ya sabemos. Comenzamos hallando las raíces del polinomio característico: 

$$p(\lambda) =  \lambda^2 - 2\lambda - 1 = 0 \iff \lambda = \frac{2 \pm \sqrt{8}}{2} \iff \begin{cases} \lambda_1 = 1 + \sqrt{2}\\ \lambda_2 = 1 - \sqrt{2} \end{cases}$$

    Como $\lambda_1$ y $\lambda_2$ son dos raíces distintas de $p(\lambda)$, sabemos que $\{ X_{\lambda_1}, X_{\lambda_2} \}$ es una base del espacio de soluciones $\Sigma$, donde 

    $$X_{\lambda_1} = \{ \lambda_1^n \}_{n\ge 0} \quad ; \quad X_{\lambda_2} = \{ \lambda_2^n \}_{n\ge0}$$
    Por tanto, la solución general $\{x_n\}_{n\ge0}$ será una combinación lineal de la base $\{ X_{\lambda_1}, X_{\lambda_2} \}$, es decir:
    \begin{equation} \label{3}
        x_n = c_1\lambda_1^n + c_2\lambda_{2}^n, \quad c_1,c_2 \in \mathbb C, \quad n\ge 0
    \end{equation}

    Para determinar las constantes $c_1$ y $c_2$, sustituimos los valores $n=0$, y $n=1$ en (\ref 3) para la solución particular $\{\overline{x_n}\}_{n\ge0}$, obteniendo el siguiente sistema de ecuaciones:

$$\begin{cases} \overline{x_0} = c_1 + c_2 = 1\\ \overline{x_1} = c_1(1+\sqrt{2}) + c_2(1-\sqrt{2}) =2 \end{cases}$$

    cuyas soluciones son $c_1 = \frac{2+\sqrt{2}}{4}$, $c_2 = \frac{2 - \sqrt{2}}{4}$. Por tanto, la expresión de la solución general es:

    \begin{equation} \label{4}
        x_n = \frac{2+\sqrt{2}}{4}\left(1+\sqrt 2 \right)^n + \frac{2 - \sqrt 2}{4}\left( 1-\sqrt 2\right)^n, \quad n \ge 0
    \end{equation}

    \textbf{c)} Deduzca de forma razonada el valor, si existe, de $\displaystyle \lim_{n\to\infty} \frac{x_{n+1}}{x_n}$.

    Como $\sqrt 2 < 2$, se tiene que $\left| 1 - \sqrt 2 \right| < 1$, y por tanto $(1-\sqrt2)^n \to 0$ cuando $n \to \infty$. Utilizando esa información, y atendiendo a (\ref 4), podemos expresar el límite de la siguiente forma:

    $$\lim_{n\to\infty} \ddfrac{x_{n+1}}{x_n} = \lim_{n\to\infty} \ddfrac{\frac{2+\sqrt{2}}{4}\left(1+\sqrt 2 \right)^{n+1} + \frac{2 - \sqrt 2}{4}\left( 1-\sqrt 2\right)^{n+1}}{\frac{2+\sqrt{2}}{4}\left(1+\sqrt 2 \right)^n + \frac{2 - \sqrt 2}{4}\left( 1-\sqrt 2\right)^n} =$$

    $$= \lim_{n\to\infty} \frac{\left( 1 + \sqrt 2 \right)^{n+1}}{\left(1+\sqrt 2 \right)^n} = \lim_{n\to\infty} \left(1 + \sqrt 2 \right) = 1 + \sqrt 2$$
    \vspace{0.2em}

    Por tanto, $\displaystyle \exists \lim_{n\to\infty} \frac{x_{n+1}}{x_n} = 1 + \sqrt 2.$
    \newpage
    \subsection{Ejercicio 13}

    Presentamos el modelo de Samuelson modificado:
    $$Y_n = C_n + I_n$$
    $$C_n=b\cdot I_{n-1}$$
    $$I_n = C_n-kC_{n-1}+G$$

    Donde $Y_n$ es la renta nacional, $C_n$ es el consumo e $I_n$ es la inversión. G es el gasto público que suponemos constante.

    Debemos escribir en primer lugar la ley de recurrencia para las inversiones anuales $I_n$ y posteriormente calcular las condiciones sobre los parámetros k y b para que haya convergencia.

    Para comenzar tomamos la expresión de $I_n$ e intentamos expresarla exclusivamente con términos de $I_n$. Para ello sustituimos en la tercera ecuación los $C_n$ por la expresión correspondiente que observamos en la segunda ecuación.

    $$I_n = b\cdot I_{n-1}-k\cdot (b\cdot I_{n-2}) +G\quad \quad n \ge 2$$

    Como nos interesa tener la ecuación en recurrencia para valores mayores o iguales que 0 hacemos una pequeña transformación en los índices (le sumamos 2 a todos)

    $$I_{n+2} = b\cdot I_{n+1}-k\cdot (b\cdot I_{n}) +G\quad \quad n \ge 0$$

    Y reordenamos para obtener la \textbf{ley de recurrencia:}

    $$I_{n+2} - b\cdot I_{n+1} + k\cdot b\cdot I_{n} = G$$

    \newpage
    Ahora procederemos a calcular las condiciones sobre k y b para que haya convergencia.

    Sabemos que la $I_n = I_n^h+I_n^{(1)}$, es decir, que podemos expresar $I_n$ como suma de la homogénea asociada más una solución particular. Y para que se produzca esta situación de convergencia la homogénea asociada debe converger a cero ( $\lim \limits_{n \to +\infty} X_n = 0 $ ).

    $$I_n^h:\ I_{n+2} - b\cdot I_{n+1} + k\cdot b\cdot I_{n} = 0 $$

    Cuyo polinomio característico es:
    $$p(\lambda) = \lambda^2 - b\lambda + kb$$

    Al estar ante una ecuación en diferencias lineal homogénea de orden dos, para que converja a 0 tiene que cumplirse lo siguiente:

    \[
        \begin{rcases}
            (1)\ \quad p(1) = 1-b + kb\ \textgreater\ 0\\
            (2)\ p(-1) = 1+b + kb\ \textgreater\ 0\\
            (3)\ \quad\quad \quad\quad p(0) = kb\ \textless\ 1\\
        \end{rcases}
    \]

    Y a partir de esto obtenemos las condiciones:
    Partimos de $1 \textgreater b \textgreater 0$ y de $k\textgreater 0$.

    Probamos que la condición (1) se cumple siempre.
    Restamos b  en la primera expresión:

    $$1-b > b-b$$
    $$1-b > 0$$

    Y de la segunda multiplicamos por b, un número positivo y la expresión se mantiene:

    $$k\cdot b > 0$$

    Si sumamos estas dos expresiones obtenemos:

    $$1-b +k\cdot b > 0$$

    La condición (2) es trivial. Suma de números positivos es positiva.

    Y de la condición (3) obtenemos:
    $$k < \frac{1}{b}$$ 

    Con esto deducimos finalmente que converge si:
    \[
        \begin{rcases}
            0 < k < \frac{1}{b}\\
            0 < b < 1
        \end{rcases}
    \]

    \section{Relación 3}

    \begin{ejer}[1]
        Demuestre el Teorema de Cayley-Hamilton.
    \end{ejer}

    \begin{sol}
        Sea $A \in \mathcal M_{r}(\mathbb C)$ y sea $p(\lambda) = |A - \lambda I| = (-1)^r(\lambda^r + c_{r-1}\lambda^{r-1} + \hdots + c_1\lambda + c_0)$ su polinomio característico, con $c_i \in \mathbb{C}$. Llamemos $C = A - \lambda I$. Sabemos que: $$C^{-1} = \frac{1}{|C|} \cdot Adj(C)^T \implies |C|I = C \ Adj(C)^T$$

        Entonces, tendremos que 
        \begin{equation} \label {th_cayley}
            p(\lambda)I = (A - \lambda I) \ Adj(A-\lambda I)^T
        \end{equation} 

        Consideremos la matriz $B = Adj(A - \lambda I)$. Como las entradas de $B$ son polinomios en $\lambda$, podemos separar para cada $0 \le i \le r$ los coeficientes del término $\lambda^i$, y formar una matriz de números $B_i$ tal que:
        \begin{align*}
            Adj(A - \lambda I)^T &= B_{r-1}\lambda^{r-1} + B_{r-2}\lambda^{r-2} + \hdots + B_0
        \end{align*}

        Ahora, sustituyendo en \eqref{th_cayley} nos queda que
        \begin{align*}
            (-1)^r(\lambda^r + c_{r-1}\lambda^{r-1} + \hdots + c_0)I &= (A -\lambda I)(B_{r-1}\lambda^{r-1} + \hdots + B_0) \\
                                                                     &= -(B_{r-1}\lambda^r + \hdots + B_0\lambda) + (AB_{r-1}\lambda^{r-1}+\hdots + AB_0) \\
                                                                     &= -B_{r-1}\lambda^r + (-B_{r-2}+AB_{r-1})\lambda^{r-1} + \hdots + (-B_0 + AB_1)\lambda + AB_0
        \end{align*}

        Identificando término a término en la igualdad anterior, notamos que se verifica lo siguiente:
        \begin{align*}
            \label{}
            (-1)^rI &= -B_{r-1} \\
            (-1)^rc_{r-1}I &= -B_{r-2}+AB_{r-1} \\
                           &\ \ \vdots \\
            (-1)^rc_1I &= -B_0 + AB_1 \\
            (-1)^rc_0I &= AB_0
        \end{align*}
        Por tanto, concluimos que
        \begin{align*}
            \label{}
            p(A) &= (-1)^r \left[ A^r + c_{r-1}A^{r-1} + c_{r-2}A^{r-2} + \hdots + c_1A + c_0I \right] \\
                 &= (-1)^r \left[ IA^r + (c_{r-1}I)A^{r-1} + \hdots + c_1IA + c_0I \right] \\
                 &= -A^{r}B_{r-1}+A^{r-1}(-B_{r-2}+AB_{r-1}) + \hdots + A(-B_0 + AB_1)A + AB_0 \\ 
                 &= 0
        \end{align*}
    \end{sol}

    \begin{ejer}[3] 
        Sea $A$ una matriz $4 \times 4$ con un valor propio $\lambda = 3$ de
        multiplicidad algebraica 4. Escriba todas las formas canónicas de Jordan
        asociadas a $A$ posibles.
    \end{ejer}

    \begin{sol}
        $$J_1 =
        \begin{pmatrix}
            3 & & & \\
              & 3 & & \\
              & & 3 & \\
              & & & 3 
        \end{pmatrix} \ \ \ 
        J_2 =
        \begin{pmatrix}
            3 & 1 & & \\
              & 3 & 1 & \\
              & & 3 & 1 \\
              & & & 3 
        \end{pmatrix} \ \ \
        J_3 =
        \begin{pmatrix}
            3 & 1 & & \\
              & 3 & & \\
              & & 3 & \\
              & & & 3 
        \end{pmatrix}$$

        $$J_4 = 
        \begin{pmatrix}
            3 & 1 & & \\
              & 3 & 1 & \\
              & & 3 & \\
              & & & 3 
        \end{pmatrix} \ \ \
        J_5 =
        \begin{pmatrix}
            3 & 1 & & \\
              & 3 & & \\
              & & 3 & 1 \\
              & & & 3 
        \end{pmatrix}$$
    \end{sol}

    \begin{ejer}[5]
        Resuelva la ecuación matricial $X^2 = I$, donde $X$ es una matriz real $2
        \times 2$.
    \end{ejer}

    \begin{sol} Sabemos que $X^2 = X \cdot X = I \iff X^{-1} = X$. Además, por el \textit{teorema de la forma canónica de Jordan}, sabemos que $X = PJP^{-1}$, y $J$ puede ser una de estas tres matrices, donde $\lambda_i \in \sigma(X)$:
        \begin{nlist}
        \item $J =
            \begin{pmatrix}
                \lambda_1 & 0\\
                0 & \lambda_2
            \end{pmatrix},\ $ con $\lambda_1 \neq \lambda_2$.
        \item  $J =
            \begin{pmatrix}
                \lambda_1 & 0\\
                0 & \lambda_1
            \end{pmatrix}$
        \item  $J =
            \begin{pmatrix}
                \lambda_1 & 1 \\
                0 & \lambda_1
            \end{pmatrix}$
        \end{nlist}

        Ahora, realizando manipulaciones algebraicas, tenemos que
        $$X^2 = X\cdot X = (PJP^{-1})(PJP^{-1}) = PJP^{-1}PJP^{-1} = P^{-1}J^2P$$

        Entonces, se verifica la siguiente equivalencia: $$X^2 = I \iff PJ^2P^{-1} = I \iff P^{-1}(PJ^2P^{-1})P = P I P^{-1} \iff J^2 = I$$

        Es decir, para que se verifique la ecuación del enunciado, necesariamente debe ser $J^2 = I$. Evaluemos ahora los casos que teníamos:
        \begin{nlist}
        \item $J^2 =
            \begin{pmatrix}
                \lambda_1^2 & 0 \\
                0 & \lambda_2^2
            \end{pmatrix} =
            \begin{pmatrix}
                1 & 0 \\
                0 & 1
            \end{pmatrix} \Rightarrow \lambda_1^2 = \lambda_2^2 = 1 \Rightarrow
            \lambda_1 = \pm 1, \lambda_2 = \pm 1$, con $\lambda_1 \ne \lambda_2$. Supongamos $\lambda_1 = 1,\ \lambda_2 = -1$. Entonces, la matriz $X$ que verifica $X^2 = I$ será de la forma
            $$X = P
            \begin{pmatrix}
                1 &  0\\
                0& -1
            \end{pmatrix}P^{-1},\quad P \in M_2(\R)\ \text{regular}$$

        \item $J^2 =
            \begin{pmatrix}
                \lambda_1^2 & \\
                            & \lambda_1^2
            \end{pmatrix} =  \begin{pmatrix}
                1 & 0 \\
                0 & 1
            \end{pmatrix} \Rightarrow J = I$ ó $J = -I$. Luego tenemos que $X = PJP^{-1} \Rightarrow X= I$ ó
            $X = -I$.
        \item $J^2 =
            \begin{pmatrix}
                \lambda_1 & 1 \\
                0 & \lambda_1
            \end{pmatrix}
            \begin{pmatrix}
                \lambda_1 & 1 \\
                0 & \lambda_1
            \end{pmatrix} =
            \begin{pmatrix}
                \lambda_1^2 & 2\lambda_1 \\
                0 & \lambda_1^2
            \end{pmatrix} = \begin{pmatrix}
                1 & 0 \\
                0 & 1
            \end{pmatrix} \ $.
            Vemos que no puede darse este caso.
        \end{nlist}
    \end{sol}

    \begin{ejer}[7] \hfill
        \begin{nlist}
        \item Demuestre que una matriz $A \in \mathbb C^{d\times d}$ es nilpotente si
            y sólo si $\sigma(A) = \{0\}$.
        \item Describa todas las matrices nilpotentes $2 \times 2$.
        \end{nlist}
    \end{ejer}

    \begin{sol}
        Decimos que una matriz $A$ es \textbf{nilpotente} si $\exists n \in \mathbb{N}$ tal
        que $A^n = 0$. Al menor $n$ tal que $A$ es nilpotente se le llama \textit{orden de nilpotencia}. Por ejemplo, las matrices
        $$ N = \begin{pmatrix}
            0 & 1 \\
            0 & 0
        \end{pmatrix}, \quad M = 
        \begin{pmatrix}
            0 & 1 & 0 \\
            0 & 0 & 1\\
            0 & 0 & 0
        \end{pmatrix}$$

        son nilpotentes de orden $2$ y $3$, respectivamente. En general, si $J$ es un bloque de Jordan con todos los elementos diagonales nulos, entonces $J$ es nilpotente de orden $\dim J$. Esto es una consecuencia directa del \textit{Teorema de Cayley-Hamilton}.

        Resolveremos ahora los dos apartados del ejercicio.

        \underline{Comenzamos probando $(i)$}. Sea $A \in \mathbb{C}^{d\times d}$. Aplicando el \textit{Teorema de la forma canónica de Jordan}, sabemos que $A$ es semejante a una matriz diagonal por bloques $J$, y también sabemos que $A^n = 0 \iff J^n = 0$. Además, es obvio que $J^n = 0$ si y solo si cada uno de los bloques que la forman es nulo.



        Consideremos un bloque arbitrario de $J$, digamos $$J_\lambda = \begin{pmatrix}
            \lambda & 1\\
                    & \lambda & 1 \\
                    & & \ddots & \ddots \\
                    & & & \lambda & 1 \\
                    & & & & \lambda
        \end{pmatrix}, \quad \lambda \in \sigma(A)$$


        Tenemos que $J_\lambda = \lambda I + F$, donde $$F = \begin{pmatrix}
            0 & 1\\
              & 0 & 1 \\
              & & \ddots & \ddots \\
              & & & 0 & 1 \\
              & & & & 0
        \end{pmatrix}$$

        Esta matriz recibe el nombre de \textit{shift matrix}, y tiene la propiedad de que sus sucesivas potencias desplazan una posición hacia la derecha la diagonal con $1$, dejando el resto de entradas a $0$. Como $F$ es un bloque de Jordan con todos los elementos diagonales nulos, es nilpotente de orden $\dim F = \dim J_\lambda = p$. Se puede demostrar por inducción la siguiente fórmula para la potencia $n$-ésima de la matriz $J_\lambda$: $$ J_\lambda^n = (\lambda I + F)^n = \sum^{n}_{k=0} \begin{pmatrix}
            n \\
            k
        \end{pmatrix} \lambda^{n-k} F^k$$
        %%% TODO: Demostrar la fórmula por inducción
        Ahora, como $F$ es nilpotente de orden $p$, tenemos que $$ J_\lambda^n = \sum^{p-1}_{k=0} \begin{pmatrix}
            n \\
            k
        \end{pmatrix} \lambda^{n-k} F^k, \quad  \forall n \geq p $$

        \begin{align*}
            \label{}
            J^n_\lambda = 
            \begin{pmatrix}
                \lambda^n & n\lambda^{n-1} & \begin{pmatrix}
                    n\\
                    2
                \end{pmatrix} \lambda^{n-2} & \hdots & \begin{pmatrix}
                    n \\
                    p-1
                \end{pmatrix} \lambda^{n-p +1} \\
                & \lambda^n & n\lambda^{n-1} & \hdots & \begin{pmatrix}
                n\\
                p-2
            \end{pmatrix} \lambda^{n-p+2}\\
            &  & \ddots & \ddots & \vdots \\
            & & & \lambda^n & n\lambda^{n-1}\\
            & & & & \lambda^n
        \end{pmatrix},\quad n \ge p
    \end{align*}

    Teniendo en cuenta todo el desarrollo anterior, concluimos que $J$ es nilpotente de orden $p$ si y solo si $\lambda = 0$.

    \textit{Demostración alternativa}. En primer lugar, sea $A \in \mathbb{C}^{d\times d}$ tal que su único valor propio es $0$. Entonces, el polinomio característico asociado será $p(\lambda) = (-1)^d\lambda^d$. Por el \textit{Teorema de Cayley-Hamilton}, tenemos que $p(A) = (-1)^dA^d = 0$, por lo que $A$ es nilpotente de orden $d$.

    Supongamos ahora que $A$ es nilpotente de orden $p$. Si $\lambda$ es un valor propio de $A$, ya sabemos que $\lambda^p$ es valor propio de $A^p$. Por tanto, existirá un $v \neq 0$ tal que $$ 0 = A^p v =  \lambda^p v,$$ lo cual es absurdo salvo en el caso $\lambda = 0$.

    \underline{Probamos ahora $(ii)$}. Consideremos $A \in \mathbb{C}^{2\times 2}$. Ya sabemos por $(i)$ que $A$ es nilpotente si y solo si su único valor propio es el $0$. $$A = \begin{pmatrix}
        a & b\\
        c & d
    \end{pmatrix} \implies p(\lambda) = |A - \lambda I| = \lambda^2 - (a+d)\lambda - bc + ad$$

    Como queremos que $0$ sea la única raíz de $p(\lambda)$, imponemos las siguientes condiciones: $$\begin{cases}
        a + d = 0\\
        bc = 0\\
        ad = 0
    \end{cases} \implies 
    \begin{cases}
        a = d = 0\\
        b = 0 \text{ ó } c = 0
    \end{cases}$$

    Es decir, las matrices nilpotentes de orden $2$ son de la forma: $$\begin{pmatrix}
        0 & 0\\
        c & 0	
    \end{pmatrix} \quad \text{ó} \quad 
    \begin{pmatrix}
        0 & c\\
        0 & 0	    
    \end{pmatrix}, \quad c \in \mathbb{C}.$$
\end{sol}

\begin{ejer}[9]
    Supongamos que el polinomio $p(\lambda) = \lambda^k + a_{k-1}\lambda^{k-1} + \hdots + a_1\lambda + a_0$ tiene $k$ raíces distintas $\lambda_1, \lambda_2, \hdots, \lambda_k$. Sea $A$ la matriz \textit{compañera} de $p(\lambda)$ definida en el problema anterior y sea $V$ la matriz de \textit{Vandermonde} $$V = \begin{pmatrix}
        1 & 1 & \hdots & 1 \\
        \lambda_1 & \lambda_2 & \hdots & \lambda_k \\
        \vdots & \vdots & & \vdots \\
        \lambda_1^{k-1} & \lambda_2^{k-1} & \hdots & \lambda_k^{k-1} \\
    \end{pmatrix}$$
    Demuestre que la matriz $V^{-1}AV$ es diagonal.
\end{ejer}

\begin{sol}
    Vamos a probar que: $V^{-1}AV = D = \begin{pmatrix}
        \lambda_1 & & & \\
                  & \lambda_2 & & \\
                  & & \ddots & \\
                  & & & \lambda_k
    \end{pmatrix}$. En realidad, lo que vamos a demostrar es que $AV = VD$.

    Vamos a multiplicar $A$ por $V$.

    \begin{align*}
        \label{}
        \begin{pmatrix}
            0 & 1 & 0 & \hdots & 0 \\
            0 & 0 & 1 & \hdots & 0 \\
            \vdots & \vdots & \vdots & & \vdots \\
            0 & 0 & 0 & \hdots & 1 \\
            -a_0 & -a_1 & -a_2 & \hdots & -a_{k-1}
        \end{pmatrix} \cdot \begin{pmatrix}
            1 & \hdots & \hdots & \hdots & 1 \\
            \lambda_1 & \hdots & \hdots & \hdots & \lambda_k \\
            \vdots & & & & \vdots \\
            \vdots & & & & \vdots \\
            \lambda_1^{k-1} & \hdots & \hdots & \hdots & \lambda_k^{k-1}
        \end{pmatrix} &= \begin{pmatrix}
            \lambda_1 & \lambda_2 & \hdots & \hdots & \lambda_k \\
            \lambda_1 & \lambda_2 & \hdots & \hdots & \lambda_k \\
            \vdots & & & & \vdots \\
            \lambda_1^{k-1} & \lambda_2^{k-1} & \hdots & \hdots & \lambda_k^{k-1} \\
            \lambda_1^k & \lambda_2^k & \hdots & \hdots & \lambda_k^k
        \end{pmatrix}\\
        &= VD \\
        &= \begin{pmatrix}
        1 & 1 & \hdots & 1 \\
        \lambda_1 & \lambda_2 & \hdots & \lambda_k \\
        \vdots & \vdots & & \vdots \\
        \lambda_1^{k-1} & \lambda_2^{k-1} & \hdots & \lambda_k^{k-1} \\
    \end{pmatrix} \\
        &\cdot \begin{pmatrix}
        \lambda_1 & & & \\
                  & \lambda_2 & & \\
                  & & \ddots & \\
                  & & & \lambda_k
    \end{pmatrix}
    \end{align*}
\end{sol}

\begin{ejer}[11]
    Sea $A$ una matriz $k \times k$ tal que $|\lambda| < 1$ para cualquier valor propio $\lambda$ de $A$. Demuestre que:
\end{ejer}

\begin{sol}\hfill\\
    \begin{itemize}
        \item Si $|A-I| = 0$ entonces $\lambda = 1$ es valor propio de $A$.
        \item $I-A$ es regular, luego $\exists (I-A)^{-1}$.
            $$\rho(A) < 1 \Leftrightarrow \{ A^{k}\}_{k \geq 0} \rightarrow 0$$
            \begin{align*}
                \label{}
                s_n &= \sum^{n}_{k=0} A^{k} = I + A + A^2 + \hdots + A^n \\
                (I - A)s_n &= (I-A) + (A-A^2) + (A^2-A^3) + \hdots + (A^n - A^{n+1}) \\
                           &= I - A^{n+1} \\
                s_n &= (I-A)^{-1}(I-A^{n+1}) \\
                \sum^{\infty}_{k=0} A^k &= \lim_{n \to \infty} s_n = (I -A)^{-1}
            \end{align*}
    \end{itemize}
\end{sol}

\begin{ejer}[15]
    Se considera la matriz de Leslie $$ \begin{pmatrix}
        \frac{1}{2} & 2 & 1 \\
        \alpha & 0 & 0 \\
        0 & \beta & 0
    \end{pmatrix},$$
    donde $\alpha, \beta \in ]0, 1[$. Encuentra las regiones donde se produce extinción, crecimiento ilimitado o convergencia a equilibrio en el plano de parámetros $(\alpha, \beta)$. 
\end{ejer}

\begin{sol}
    La tasa de fertilidad del primer grupo es $\frac{1}{2}$, la del segundo $2$, y la del tercero, $1$.
    \begin{align*}
        \label{}
        X_n = L^nX_0 \\
        \left\{ \frac{1}{\lambda_1^n} L^nX_0 \right\}_{n \geq 0} \rightarrow V \ ; \ \ LV = \lambda_1V
        \end{align*}

        Entonces \begin{itemize}
            \item $0 < \lambda_1 < 1 \Rightarrow$ extinción.
            \item $\lambda_1 = 1 \Rightarrow$ equilibrio.
            \item $\lambda_1 > 1 \Rightarrow$ crecimiento ilimitado.
        \end{itemize}
\end{sol}

\begin{ejer}[17]
    Se considera la matriz de Leslie $$\begin{pmatrix}
        \frac{1}{2} & \delta & \gamma \\
        \frac{1}{2} & 0 & 0 \\
        0 & \frac{1}{2} & 0
    \end{pmatrix}$$ donde $\delta, \gamma \geq 0$. \begin{enumerate}[label=\alph*)]
        \item Encuentre las regiones donde se produce extinción, crecimiento ilimitado o convergencia a equilibrio en el plano de parámetros $(\delta, \gamma)$.
        \item Describa la pirámide de edad a largo plazo correspondiente a los valores $\delta = \frac{1}{2}, \gamma = 1$.
        
    \end{enumerate}
\end{ejer}

\begin{sol}\hfill
    \begin{enumerate}[label=\alph*)]        
        \item $R = \frac{1}{2} + \frac{1}{2}\delta + \frac{1}{4}\gamma$
            \begin{align*}
                \label{}
                \frac{1}{2} + \frac{1}{2}\delta + \frac{1}{4}\gamma = 1 \\
                \frac{1}{2}\delta + \frac{1}{4}\gamma = \frac{1}{2} \\
                2\delta + \gamma = 2 \\
                \gamma = 2 - 2\delta
            \end{align*}
        \item $\delta = \frac{1}{2}, \gamma = 1$

            $$L = \begin{pmatrix}
                \frac{1}{2} & \frac{1}{2} & 1 \\
                \frac{1}{2} & 0 & 0 \\
                0 & \frac{1}{2} & 0
            \end{pmatrix}$$

            $\exists \lambda_1$ valor propio dominante (único positivo). ¿Es $\lambda_1 = 1$?

            \begin{align*}
                \label{}
                |L-\lambda I| = \begin{vmatrix}
                    \frac{1}{2} - \lambda & \frac{1}{2} & 1 \\
                    \frac{1}{2} & -\lambda & 0 \\
                    0 & \frac{1}{2} & -\lambda
                \end{vmatrix} &= \left( \frac{1}{2} - \lambda \right)\lambda^2 - \frac{1}{2}\left( \frac{-\lambda}{2} - \frac{1}{2} \right) \\
                              &= -\lambda^3 + \frac{1}{2}\lambda^2 + \frac{1}{4}\lambda + \frac{1}{4}
            \end{align*}
        
            $p(1) = 0 \Rightarrow \lambda_1 = 1$

            Buscamos $v_1 : (L - I)v_1 = 0$
            \begin{align*}
                \label{eq:}
                \begin{pmatrix}
                    -\frac{1}{2} & \frac{1}{2} & 1 \\
                    \frac{1}{2} & -1 & 0 \\
                    0 & \frac{1}{2} & -1
                \end{pmatrix} \begin{pmatrix}
                    v_1 \\
                    v_2 \\
                    v_3
                \end{pmatrix} = \begin{pmatrix}
                    0 \\
                    0 \\
                    0
                \end{pmatrix}
            \end{align*}

            $$ \begin{rcases}
                \frac{1}{2}v_1 - v_2 = 0 \\
                \frac{1}{2}v_2 - v_3 = 0
            \end{rcases} \Rightarrow \begin{rcases}
                v_2 = \frac{1}{2}v_1 \\
                v_3 = \frac{1}{2}v_2
            \end{rcases}$$

            Por ejemplo $v_1 = 4 \Rightarrow \begin{cases}
                v_2 = 2 \\
                v_3 = 1
            \end{cases}$

            $$V = \begin{pmatrix}
                4 \\
                2 \\
                1
            \end{pmatrix} \ \ \ \ \frac{V}{||V||_1} = \begin{pmatrix}
                \frac{4}{7} \\
                \frac{2}{7} \\
                \frac{1}{7} 
            \end{pmatrix}$$
     \end{enumerate}
\end{sol}

\begin{ejer}[20]
  Una determinada población está estructurada en base a tres grupos diferentes de edad: crías (hasta los 3 años),
jóvenes (de 3 a 6 años) y adultos (de 6 a 9 años). Es conocido que cada cría engendra en media una nueva cría, cada
joven engendra en media 1,5 crı́as y cada adulto engendra en media 0,5 crı́as. Además, las observaciones arrojan el
dato de que la mitad de las crı́as llegan a jóvenes, en tanto que sólo el 20 \% de los jóvenes sobrevive.

\begin{itemize}
\item Construya la matriz del modelo.
\item Si la distribución de tamaños iniciales es $P_0 = (3, 1, 0)^t$ (en las unidades adecuadas), calcule cuál será la distribución de tamaños al cabo de seis años.
\item Explique el comportamiento a largo plazo de la población (incluyendo su distribución porcentual por grupos
\end{itemize}
de edad).

\end{ejer}

\begin{sol}\hfill\\
  \begin{enumerate}[label=\alph*)]
  \item    La matriz de Leslie del modelo queda como sigue:

    \begin{center}
      \begin{pmatrix}
        
        1 & 1.5 & 0.5\\
        0.5 & 0 & 0\\
        0 & 0.2 & 0
        
      \end{pmatrix}
    \end{center}

 
 \item    Al cabo de seis años, la distribución será


    \[
    % P_6 = A^6 P_0 = \left(36.183125,\,11.930625,\,1.57375\right)
    P_6 = A^2 P_0 = \left(6.85,\ 2.25\ ,\ 0.3\right)
    \]

  \item    Para calcular la distribución porcentual por grupos de edad a largo plazo, primero vamos a comprobar la
    \emph{tasa de natalidad}:

    \[
    R = 1 + 1.5\times 0.5 + 0.5\times 0.5 \times 0.2 = 1.8 > 1
    \]

    Por tanto, la población no se extinguirá a largo plazo, sino que crecerá indefinidamente.
    Calculemos la pirámide de población. Para ello, tenemos que calcular el valor propio dominante de
    la matriz del modelo. Un breve análisis del polinomio característico basta para comprobar que no podemos calcular
    sus raíces mediante procedimientos algebraicos. El método de la potencia nos da el
    valor propio $\lambda_1 \approx 1.51$, y el vector propio asociado $v_{\lambda_1} \approx \left(1.000000,\,0.3297385,\,0.04349099\right)$.

    La pirámide de población a largo plazo será, aproximadamente:

    \[
    \frac{v_{\lambda_1}}{\|v_{\lambda_1}\|_{1}} \approx (0.72821, 0.24011, 0.03167)
    \]
  \end{enumerate}
\end{sol}

\end{document}


\end{document}
